\documentclass[11pt,a4paper]{article}
\usepackage[utf8]{inputenc}
% lmodern must precede fontenc: without a scalable Type1 font the T1 encoding
% falls back to bitmap fonts, and microtype's expansion then aborts the run
% with "auto expansion is only possible with scalable fonts".
\usepackage{lmodern}
\usepackage[T1]{fontenc}
\usepackage[margin=1in]{geometry}
\usepackage{amsmath,amssymb}
\usepackage{graphicx}
\usepackage{microtype}
\usepackage[hidelinks,breaklinks]{hyperref}
\usepackage{url}
\urlstyle{same}

% Generated by code/tools/md_to_tex.py from the SurveyForge Markdown output.
% Citation numbers match the source .md exactly.

\title{3D Gaussian Splatting: A Comprehensive Survey of Representations, Optimization, and Emerging Applications}
\author{Generated by SurveyForge}
\date{}

\begin{document}
\maketitle
\tableofcontents
\newpage



\section{Introduction}


The evolution of three-dimensional scene representation has been defined by a persistent tension between geometric fidelity, computational efficiency, and photorealistic appearance. Classical approaches founded on explicit meshes and point clouds offered direct manipulation and hardware-accelerated rendering, yet they struggled to capture the continuous, view-dependent appearance of real-world scenes. The introduction of implicit neural fields, most prominently Neural Radiance Fields (NeRFs), addressed the appearance bottleneck by treating scenes as continuous spatial functions that map 3D coordinates and viewing directions to radiance. However, while offering state-of-the-art quality, coordinate-based neural networks proved computationally prohibitive---both in the dense ray marching required to train them and the costly volumetric sampling that limited real-time inference. Achieving high visual quality remained contingent on neural networks that are inherently costly to train and render, while the faster variants traded speed for quality \cite{arxiv:2308.04079v1}. This performance-quality impasse catalyzed the search for representations that could preserve the advanced fidelity of neural fields while enabling parallelism and real-time rasterization.

The fundamental conceptual leap of 3D Gaussian Splatting (3DGS) was the realization that the entire scene could be explicitly represented as a set of thousands of anisotropic, differentiable 3D Gaussian primitives. More than a mere algorithmic adaptation, this represented a philosophical shift: instead of embedding a dense neural field over the entire spatial domain, the scene became a set of discrete primitives whose positions, opacity, 3D covariance, and color properties (parameterized by Spherical Harmonics) are collectively optimized against photometric loss functions \cite{arxiv:2308.04079v1}. This raises questions about the mathematical underpinnings of the representation and the primitive behavior during the density control step. As an example, the conventional densification heuristics---splitting and cloning---are now shown to behave in ways that strengthen the argument that we can rethink the set of 3D Gaussians as a random sample drawn from an underlying probability distribution, i.e., Markov Chain Monte Carlo samples updated via Stochastic Gradient Langevin Dynamics \cite{arxiv:2404.09591v3}. This view opens a pathway towards a more mathematically grounded treatment of the procedural aspects of the explicit point-based representation. The tile-based differentiable renderer that sorts these Gaussians per-screen-tile is a key device that avoids the ray marching of volumetric rendering, enabling state-of-the-art performance at 1080p resolution at real-time frame rates, without the runtime overhead of directly integrating volumetric methods \cite{arxiv:2308.04079v1}. This explicitly shift towards parallel computing and differentiable rendering is highlighted in recent surveys \cite{arxiv:2401.03890v2}.

However, as is reported in \cite{arxiv:2403.17888v1}, 3DGS fails to accurately represent surfaces due to the neural view-equivalent nature of 3D Gaussians, because the 3D Gaussian structure fails to model a deterministic, view-independent geometric neutrally. This is a core limitation: the algorithm is view-perfectionist but geometrically ambiguous. To remedy this, the representation is constrained, e.g., by collapsing the 3D volume into a set of 2D oriented planar Gaussian disks to ensure view-consistent and surface-driven geometric constraints \cite{arxiv:2403.17888v1,arxiv:2311.12775v3}. Moreover, the optimization process relies on regularizers and performance levers to maintain stable training for challenging inputs \cite{arxiv:2403.13806v1,arxiv:2404.10772v1}.

In terms of quality, explicit Gaussian-based forward models have been extended in two significant directions. On the one hand, to improve fidelity with respect to high-frequency appearance and anisotropic reflections, the early use of Spherical harmonics has been complemented (or replaced) with alternative appearance models such as anisotropic spherical harmonics \cite{arxiv:2402.15870v1} and directional distributions \cite{arxiv:2512.14180v2}. On the other hand, the rendering pipeline has been improved by reducing the number of primitives while maintaining a computation-optimal constant rendering time \cite{arxiv:2412.00578v3}.

The development of the Gaussian Splatting is closely intertwined with the need to represent large-scale, unbounded environments. Geometry-bearing accelerating structures were required to make flat variants closed-off to unrealistically small scenes. Initializing structure from SfM point clouds is the default choice; however, the evaluation spanned arbitrary and bounded scenes and highlighted the important role of primitives on load and uneven optimization, often solved by distributing the scene into partitions and hierarchical training \cite{arxiv:2406.12080v1}. This has driven the evolution of the 3D Gaussian representation from an almost purely static image-rendering approach to perhaps one of the foundational ideas of 4D dynamic reconstruction \cite{arxiv:2412.20720v2,arxiv:2310.08528v2}.

The 3DGS framework, representing the scene explicitly as a structure of primitives, offers unique advantages over implicit representations: the system scaffolds applications like physically-based inverse rendering \cite{arxiv:2602.17117v1}, object detection \cite{arxiv:2504.13204v2}, and robotics \cite{arxiv:2410.12262v2}. Nonetheless, the blanket application of 3DGS in a diverse and complex domain is still limited by two fundamental challenges: (1) The geometry quality measured against sensor accuracy (such as LiDAR-depth) is not guaranteed by the photometric loss, requiring depth supervision \cite{arxiv:2403.17822v1} and normal priors \cite{arxiv:2311.17910v1} (2) the memory footprint for the number of primitives typically grows extremely fast to fit dense representations and storage. The forces are directed toward memory-efficient, and bipartite solutions use vector quantization and entropy coding \cite{arxiv:2311.18159v3,arxiv:2509.07021v3}. Alternative 2D-Gaussian variants are more memory-efficient---with spatially varying colors \cite{arxiv:2411.18966v2}---to satisfy real-time devices (see surveys \cite{arxiv:2502.19457v1}). 

The scope of this survey is designed as a study of foundations (Section 2), geometric and robustness extensions (Section 3), compression and edge deployment (Section 4), dynamic and physical extensions (Section 5), and application (Section 6). By categorizing the recent methodological advances in each of these classes, and synthesizing between the common primitives of Gaussian, this survey provides a deeper analysis of why Gaussian primitives provide a trade-off point between optimization stability and compatibility with standard rendering engine. In closing, the conclusion stresses that 3DGS, as the explicit representation, holds strong potential for a new generation of 3D "digital material" beyond view synthesis, enabling high-fidelity interactive experiences \cite{arxiv:2412.06257v2}.


\section{Foundations, Rendering Pipelines, and Core Algorithmic Enhancements}



\subsection{Initialization Strategies and Primitive Seeding}


The initialization of 3D Gaussian primitives establishes the optimization landscape upon which all subsequent densification and refinement operate, fundamentally conditioning convergence behavior, final reconstruction quality, and geometric consistency. The canonical approach, introduced with the original 3D Gaussian Splatting framework \cite{arxiv:2308.04079v1}, seeds the optimization with sparse Structure-from-Motion (SfM) points, which provide a coarse yet globally coherent geometric scaffolding. This paradigm excels in well-textured, densely captured environments where feature matching yields reliable calibration output; the multi-view consistency encoded in SfM points effectively anchors Gaussians to geometrically plausible positions, mitigating the risk of early divergence and accelerating photometric convergence. However, SfM's dependence on robust feature correspondence reveals systematic failure modes: texture-poor regions receive insufficient coverage, sparse-view configurations undermine feature matching reliability, and pose estimation errors propagate directly into the initialized geometry. These vulnerabilities have catalyzed the exploration of alternatives that decouple initialization from classical keypoint detection.

A significant body of work replaces SfM with learned dense priors. Radiance field-informed pipelines \cite{arxiv:2403.13806v1} leverage a pre-trained NeRF to supervise and guide Gaussian placement, improving robustness in complex captures at the cost of substantial preprocessing. More recent approaches synthesize initial point clouds through feed-forward monocular and stereo networks. For instance, many contemporary systems distill dense priors from multi-view stereo regressors or produce point maps via pointmap regression networks, yielding high-quality reconstructions without legacy feature matching. A notable instantiation is the method of EasySplat-style pipelines, which group geometrically consistent views from image-matching transformers to construct robust, non-redundant point-cloud priors that reduce training time and improve sparse-view performance. Similarly, splatter-based generalizable models, such as Splatter-360, process multi-view features through cross-attention mechanisms to directly predict Gaussian parameters in a single forward pass, effectively replacing optimization-based initialization with a learned mapping, though often at reduced fidelity for out-of-distribution scenes.

A parallel debate interrogates whether explicit structure is required at all. Works such as 3D Gaussian Splatting as Markov Chain Monte Carlo, show that a random, densification-driven scheme combined with stochastic transition dynamics can converge to parity with SfM-seeded counterparts while simplifying the algorithmic pipeline. This perspective reinterprets Gaussian placement as sample distribution, with initialization acting as an initial sample set from which transitions explore the posterior. The practical implications are profound: when Monte Carlo sampling is correctly applied with spatially-varying transition kernels and opacity-driven regularization, convergence speed and rendering quality become resilient to poor initial geometry, requiring only a crude device bound. However, evidence from Deformable Beta Splatting extends and refines this insight, proving that distribution-preserving refinement is achievable with any splatting kernel under regularized opacity updates, untethering probabilistic densification from Gaussian-specific properties. The residual limitation remains: density-driven refinement alone can leave off-trajectory novel views geometrically degraded compared to structure-seeded counterparts; absent a strong regularizer, these approaches rely on sufficiently dense initial posterior samples to avoid forming unnatural volumetric shells.

Addressing these robustness gaps, a class of generalizable frameworks adopts single-forward-pass refinement strategies. SplatFormer-style methods apply point transformer blocks directly to an initially seeded set (from SfM or random) and iterate over self-attention and cross-attention contexts to produce accurate Gaussian sets in a single forward pass. This heuristically grounds out-of-distribution improvements by aligning prediction geometry to inferred depth consensus, without requiring multiple gradient steps. Beyond pure forward mapping, the recurrent refinement of ReSplat uses dense Gaussian renderings of error maps as feedback signals to refine intermediate Gaussian sets sampled from priors, achieving cross-domain adaptation and strong performance even without back-propagating through the full expression, again legally delocalizing the whole problem onto the upstream predictor.

The trade-off panorama is now clear: SfM offers robust geometry where feature quality is sufficient but fails in poorly textured settings; learned techniques can fill this geometric gap by leveraging large-scale priors; pure random initialization with principled Monte Carlo or decomposition expansion can sidestep legacy dependence at slight geometric fidelity cost; while hybrid feed-forward and recurrent models transform the problem from an optimization initiated problem to a density prediction problem entirely. Future research will likely converge on frugal hybrid formulations---initializing coarse global structure with robust learned priors, then slicing the plumbing into local refined neighborhoods---to attain sparse generalizability without sacrificing per-scene performance or interpretability.


\subsection{Adaptive Density Control and Principled Primitive Refinement}


Adaptive density control constitutes the central feedback mechanism that governs the allocation of representational capacity in 3D Gaussian Splatting (3DGS), translating the initialization landscape discussed above into a dynamic process of primitive refinement and topological adjustment. The original formulation \cite{arxiv:2308.04079v1} established a heuristic paradigm based on screening per-Gaussian view-space positional gradients to trigger topological modification---cloning in under-reconstructed regions and splitting in over-reconstructed ones---coupled with opacity-based pruning and periodic opacity resets to eliminate floating artifacts and mitigate local optima. While undeniably effective, this gradient-magnitude heuristic is fundamentally problematic: it conflates two distinct physical phenomena---true geometric misplacement versus aliasing-induced positional error arising from images that contain too little local texture---and exhibits severe positional and temporal biases. Consequently, optimization can converge prematurely in texture-rich regions while over-densifying uniformly textured or low-frequency areas, resulting in massive and redundant representations. This limitation is not merely a practical inconvenience; it directly conditions the optimization dynamics discussed in the next subsection, as an uncontrolled densification process can exploit the photometric loss landscape by spawning additional Gaussians to overfit training views, thereby reinforcing the co-adaptation and geometric instability that photometric-only optimization tends toward.

A principal direction of refinement is the reformulation of density control within an explicit probabilistic framework. Instead of treating densification as a deterministic rule, the MCMC-driven method \cite{arxiv:2404.09591v3} reinterprets Gaussian parameter updates as Stochastic Gradient Langevin Dynamics and repurposes densification as a relocalization of marginal samples, preserving the underlying probability field while simplifying the algorithmic pipeline. Concurrently, gradient information that adheres to statistical interpretation can be mined: gradient-coherence-ratio-based methods \cite{arxiv:2508.09239v2} directly distinguish between conflicting and aligned gradient signals to separate the causes and consequences of under- and over-reconstruction. This associated analysis \cite{arxiv:2505.05587v1} establishes a sufficient condition for densification; analytically, the steepest descent direction for splitting is the one that minimizes the immediate photometric loss for a given error budget, and it proves that with a fixed number of offspring primitives, it is possible to escape saddle points in the training dynamics via well-chosen splitting moves \cite{arxiv:2505.05587v1}. These principled formulations align density control with the feedback-based perspective developed earlier in this subsection: rather than relying on heuristics that amplify noise in the loss landscape, density actions become decisions that are increasingly coherent with the optimization trajectory.

A complementary thread emphasizes structure-awareness and resource efficiency. Structure-adaptive methods, e.g., Mini-Splatting2 \cite{arxiv:2603.20857v1} propose a structure-aware distribution that first reorganizes and spatially rebuilds primitives using K-nearest-neighbor and covariance statistics to break the coupling with view, and then prioritizes gradients based on geometric coverage and optimization readiness, allowing for faster convergence and compactness of the representation \cite{arxiv:2603.20857v1}. Rather than relying on dense covariance reorganization, the introduction of soft ,,odds``-shaped masks for Gaussian pruning \cite{arxiv:2412.20522v3} models each primitive's contribution as a probability and leverages masked-rasterization to push gradients into probabilistically retained but masked Gaussians, retaining reversibility and adaptive adjustments in the selection process during optimization. Combined with the above trust in informative primitives, rendering-free importance prediction emerges as an efficient pruning route: RAP \cite{arxiv:2602.19753v1} predicts primitive importance directly from attributes (position, opacity, covariance, and harmonic coefficients), avoiding the substantial cost of rendering-based analyses; however, for such pruning to retain fidelity, the method must be trained specifically to be resilient to collapses of masked-out primitives, as highlighted in \cite{arxiv:2602.19753v1}.

This shift toward constraint-driven, resource-aware, and theory-guided decisions has yielded new control schemes: (1) hierarchical exploitation of multi-scale errors to specify adaptive densification triggers, (2) incorporating local geometry and depth cues for sparse-metrically challenging regions \cite{arxiv:2501.11508v1}, which highlight the need to consider worst-case geometric fidelity, and (3) propagating Gaussians along locally inferred surfaces for more robust and scalable densification, as demonstrated in \cite{arxiv:2402.14650v1}. Nevertheless, open theoretical questions remain: how to prune conflicting priors without corrupting the reconstruction of thin structures (an issue partly identified by so-called truncated-density modeling approaches), and how to balance the cost of evidence-based, cost-aware splitting against large-scale performance benefits. The next stage of development---namely, density control strategies that are simultaneously decision-informed, loss-aware, and dataset-agnostic---promise not only to narrow the gap between compactness and fidelity but to unify optimization and geometry-aware decision-making. This trajectory set the stage for the photometric-aware and optimization-centric refinements discussed in the next section, where the ventional loss and gradient-driven approach meets systematic design of the optimization.


\subsection{Loss Design, Geometric Regularization, and Training Stability}


The optimization of 3D Gaussian Splatting (3DGS) hinges on a delicately balanced loss landscape where photometric fidelity must be reconciled with geometric plausibility. The foundational objective remains a combination of L1 distance and structural similarity (D-SSIM), which drives the reconstruction towards pixel-level accuracy while preserving edge sharpness and texture boundaries. This dual-term loss is computed between the rendered image and ground-truth observations, yet its singular reliance on photometric error creates a large solution space where geometry and appearance can co-adapt in physically implausible or visually unstable ways, a failure mode that becomes particularly acute in sparse-view and under-constrained settings. The measurement of co-adaptation, introduced by the Co-Adaptation Score (CA), elegantly quantifies the degree to which optimized Gaussians become entangled, revealing that overfitting to training views manifests as instability in novel views; two lightweight interventions---random Gaussian dropout and multiplicative opacity noise-injection---directly combat this co-adaptation, underscoring that the optimization's stability is not merely a matter of the loss value at convergence, but also of the path taken by the geometric parameters \cite{arxiv:2508.12720v3,arxiv:2402.00525v1}.

To navigate the ill-posed nature of photometric-only optimization, a growing body of work integrates geometric priors and regularization to constrain the solution space. Surface-driven approaches diversify the optimization targets, and evidence demonstrates that regularizing depth alone lacks sufficient constraint; normal alignment must be actively enforced to suppress naive geometric drift \cite{arxiv:2501.08174v2,arxiv:2510.08566v1}. Rather than enforcing generic smoothness priors, which can over-smooth fine features, a compelling strategy for sparse-view and optimization-stability-limited regimes is the modification of the loss landscape in frequency space. FreGS (FreGS: 3D Gaussian Splatting with Progressive Frequency Regularization) employs strong and explicit frequency-domain supervision that directly penalizes spectral mismatches between the rendering and the ground truth. This frequency regularization divides the optimization into a coarse-to-fine densification, avoiding exhaustion with intermediate-frequency ambiguities that induce over-densification \cite{arxiv:2403.06908v2,arxiv:2503.07000v1}. The relationship between photographic and geometric complexity is further explored by the Scale-Adaptive Gaussian Splatting for Training-Free Anti-Aliasing, which applies scale-adaptive filters at test-time, bypassing the need to modify the training loss. Its success suggests that both the training objective and the sample-rate assembly must be calibrated to the observed sampling rate \cite{arxiv:2403.19615v1,arxiv:2311.16493v1}.

The gradient itself, however, remains a critical supervisory signal that is often confounded. The naive positional-gradient heuristic used for densification in the original formulation struggles to perfectly separate the observed noise of the optimization path from a true geometric defect. Pixel-GS introduces a principled gradient weighting strategy that scales the densification when each Gaussian covers more pixels, predicated on the account of covered pixel numbers, thereby suppressing floaters and blurring in under-reconstructed regions \cite{arxiv:2403.15530v1}. Similar rigor inherent in the theoretical analysis of steepest descent coefficients leads to SteepGS (Steepest Descent Density Control for Compact 3D Gaussian Splatting). By mapping splitting to a stepping-stone candidate in the optimization landscape---in which Gaussian splitting is proven to be crucial for escaping saddle points---this framework formulates the density control as a minimization problem with a minimum number of offsprings, directly promoting a loss-driven optimization of geometry \cite{arxiv:2505.05587v1}. The influence of conflicting gradient conventions is addressed in Gradient-Direction-Aware Density Control for 3D Gaussian Splatting, where the coherence ratio of the gradient vectors distinguishes under-reconstructed (conflicting) versus over-densified (aligned) regions, ensuring that similar gradients are not mistakenly aggregated in conflicting ways \cite{arxiv:2508.09239v2}.

Given the increasing reliance on test-time stability, novel optimizers are concurrently being developed to tackle the underlying numerical complexity of the loss surface. Existing first-order optimizers often stagnate in saddle points. Yet, by exploiting the sparsity in the Jacobian---a structure frequent in Gaussian--pixel relations---a matrix-free second-order optimization strategy (Matrix-free Second-order Optimization of Gaussian Splats with Residual Sampling) allows for a computationally tractable approximation of the curvature \cite{arxiv:2504.12905v3}. This relaxation dramatically accelerates convergence and increases robustness to local minima, given only a moderate amount of consumed memory. In summary, the evolution of the loss or shallow optimizer design reveals a fundamental transition: rather than accepting the unwieldy optimization dynamics that can lead to transiency and fantasy geometry, the training pipeline increasingly incorporates photometric reconstruction with physical coherence, weakened and augmented geometric feedback, frequency-aware supervision, and inference-time parameter stochastization, thereby maintaining geometry, color, and opacity in equilibrium.


\subsection{Anti-Aliasing, Multi-Scale Consistency, and Primitive Geometry Extensions}


The fidelity of 3D Gaussian Splatting (3DGS) is fundamentally tied to the sampling rate at which it is optimized and rendered. When the image-space frequency of the scene exceeds the Nyquist frequency dictated by the pixel grid, aliasing artifacts---such as jagged edges, flickering textures, and shimmering highlights---emerge. Conversely, when Gaussians are rendered at a scale significantly smaller than a pixel, they introduce over-blurring and a substantial loss of detail. The root cause lies in the continuous parametric representation itself: a 3D Gaussian integrates energy over a finite footprint in world space, but its projection onto the image plane is a continuous kernel that must be symbolically point-sampled at pixel centers. Without explicit pre-filtering or scale-aware attenuation, this sampling process violates the Nyquist criterion, leading to the aforementioned degradations. These challenges are not merely post-processing concerns; they are intrinsically coupled to the optimization dynamics and system-level efficiency discussed in the preceding and following subsections. On one hand, a poorly chosen loss landscape or an unstable densification schedule---as discussed in the previous subsection---can exacerbate the prevalence of sub-pixel Gaussians, making the rendering more susceptible to aliasing regardless of the anti-aliasing strategy applied. On the other hand, the system-level co-design and primitive-count reduction strategies presented in the next subsection are increasingly intertwined with scale-adaptive filtering, because reducing the number of primitives or dynamically adjusting resolution directly alters the sampling rate requirements and the feasibility of explicit multi-scale representations. This subsection dissects the spectrum of solutions addressing these intertwined issues, ranging from low-pass filtering to fundamental kernel geometry reformulations and size regulation.

The most direct mitigation of aliasing is to apply a low-pass filter in the frequency domain before sampling. \cite{arxiv:2311.16493v1}(paper\_title: Mip-Splatting) established a foundational framework by decomposing the problem into 2D screen-space filtering and 3D variance regularization. The former filters the Gaussian's contribution in image space using a low-pass kernel tied to the sampling rate, while the latter constrains the maximum size of a Gaussian in 3D space to prevent it from projecting to a sub-pixel footprint, which effectively acts as a guard against high-frequency noise \cite{arxiv:2311.16493v1,arxiv:2603.19682v2}. This dual approach, however, operates in the realm of isotropic projections. A critical evolution of this idea is found in \cite{arxiv:2506.02751v3} (paper\_title: 2DGS), which not only introduces a 2D Gaussian kernel but also captures the object-space Mip filter for these surfels, ensuring that the ray-splat intersection is computed with a bandwidth that respects the continuous nature of the primitive \cite{arxiv:2506.02751v3}. This shift from a screen-space filter to a world-space (object-space) counterpart is a crucial step towards scale-awareness: it aligns the filtering with the actual ray-casting geometry rather than post-hoc image-space operations. More recent work \cite{arxiv:2504.13204v2} (paper\_title: HDGS: Textured 2D Gaussian Splatting for Enhanced Scene Rendering) pairs a frustum-based sampling strategy with the 2D surfels to further refine anti-aliasing behavior, particularly for high-resolution rendering \cite{arxiv:2504.13204v2}. These filtering-based solutions demonstrate that the sampling rate must be encoded within the primitive's projection itself---a theme that becomes increasingly essential when considering the variable primitive budgets of system-level optimizers.

Parallel to the filtering-focused approaches, a distinct line of research tackles the problem by building explicit, multi-scale representations. This is a significant departure from a single scale of primitives that must adapt to all viewing distances. These architectural choices explicitly decouple the level of detail from the Gaussian count, allocating few, larger primitives to distant regions and more numerous, smaller ones for close-ups. The \emph{large images as Gaussians} (LIG) approach (paper\_title: Large Images are Gaussians: High-Quality Large Image Representation with Levels of 2D Gaussian Splatting) directly implements this with a Levels-of-Gaussian (LoG) pyramid, where a coarse-to-fine hierarchy is trained iteratively, showing that this representation allows the capture of low-frequency structure and then the slow addition of high-frequency details in a scale-adaptive manner. This scale-adaptive allocation is not only a qualitative boon but also a computational one, as it can achieve substantial rendering speedups in large scenes by reducing the need to heavily populate far-range views. This philosophy is echoed in scheduling frameworks like DashGaussian (paper\_title: DashGaussian: Optimizing 3D Gaussian Splatting in 200 Seconds), which, while primarily focused on optimization speed, also imposes a functional form of scale-aware growth: the number of primitives grows in sync with the resolution schedule, preventing an over-allocation of low-resolution primitives that are not used. This coupling of resolution and density is a functionally equivalent approach to scaling the representation, and its success across these projects reinforces the notion that a single-scale Gaussian field is ill-suited for the dynamic range of real-world scenes. Importantly, these multi-scale frameworks are natural prerequisites for the hierarchical system-level co-design discussed in the following subsection, where the Level-of-Detail (LOD) structures directly benefit from the scale-adaptive primitive distributions established here.

On the third front, there is a geometric fundamental for the intrinsic kernel itself. The classic 3D kernel struggles to represent flat surfaces, sharp edges, and thin structures---particularly where the scale approaches the pixel. The intuitive leap to a 2D Gaussian, or surfel, as done in \cite{arxiv:2506.02751v3}(paper\_title: 2DGS), not only provides a more surface-aligned geometry but also implicitly offers a form of scale control. Since a 2D Gaussian has a well-defined orientation in space, its projection onto the image plane is far more predictable and can be more precisely filtered and sampled than an anisotropic 3D ellipsoid. This eliminates the volumetric rotational ambiguity of its 3D counterpart, which often leads to floaters and over-projection. \cite{arxiv:2504.13204v2}(paper\_title: HDGS: Textured 2D Gaussian Splatting for Enhanced Rendering) shows that these 2D primitives can be further enhanced with texture maps to fit textures specifically while maintaining scale-correct sampling, an advantage when the number of primitives is reduced.

Beyond kernel geometry, the methodology becomes one of physical grounding, preventing degenerate scale collapse. A naive optimization can greedily shrink a Gaussian's variance to a disk-like vector---a direct failure mode of the optimization dynamics highlighted in the previous subsection, often exacerbated by the naive positional-gradient heuristic unless regularized. Techniques such as DropGaussian introduce structural regularization and random Gaussian removal (paper\_title: DropGaussian: Structural Regularization for Sparse-view Gaussian Splatting). This operation implicitly enforces a form of scale reciprocity: by forcing the remaining Gaussians to absorb a higher gradient, they tend to become larger and more region-aligned, making them less prone to collapse into sub-pixel artifacts. More principled, geometry-focused methods, like those incorporating deep-normal priors, anchor Gaussian sizes to the scene geometry, mitigating floaters and ensuring that Gaussian extents are bounded---a necessary condition for stable sampling under scale changes (paper\_title: DN-Splatter: Depth and Normal Priors for Gaussian Splatting and Meshing; paper\_title: GausSurf: Geometry-Guided 3D Gaussian Splatting for Surface Reconstruction).

In summary, the path from screen-space filters to 3D regularization, then to 2D kernel geometries, and finally to explicit multi-scale approaches, indicates that the field has moved from treating anti-aliasing as a render-time concern to a strategy that forms the core attribute of the primitive definition itself. The future of robust scale-adaptive 3DGS hinges on unifying these approaches---not just by modifying the geometry of the Gaussian, but by fundamentally coupling the filter footprint to the physically grounded kernel. This unification also carries system-level consequences: a more stable and scale-aware primitive distribution reduces the need for the heavy computational pruning and level-of-detail selection discussed in the following subsection, as fewer redundant primitives must be managed. Conversely, system-level advances in primitive reduction directly relax the aliasing burden, creating a positive feedback loop between these interrelated design axes. The next wave of research will likely optimize the "costless" unification of sampling rate, kernel shape, and scene geometry to build a single, lightweight model that remains high fidelity across all camera distances while accommodating the opacity and color-mapping features for full fidelity.


\subsection{Performance, Memory Efficiency, and System-Level Co-Design}


The relentless pursuit of real-time performance and manageable memory footprints has driven significant system-level co-design in 3D Gaussian Splatting (3DGS), moving beyond the original differentiable tile-based rasterizer. While the foundational pipeline \cite{arxiv:1906.04173v3} established differentiable point-based rendering, modern 3DGS implementations exhibit substantial computational redundancy. The original rasterizer, while efficient, performs a full sort and processes all splats per tile, a workload that grows superlinearly with scene complexity. This has motivated a spectrum of algorithmic and hardware-aware optimizations targeting the rendering pipeline, training acceleration, and primitive-count reduction.

System-level rendering efficiency has been advanced through several complementary strategies. Methods have explored abandoning explicit sorting altogether; order-independent transparency and stochastic rasterization techniques circumvent the need for a global sort, achieving a significant reduction in compute load, demonstrating that the strict depth-ordered blending of the original algorithm is computationally over-constrained for certain rendering scenarios. To further enhance throughput, optimization of load-balancing and scheduling has been targeted. For instance, the approach used in \cite{arxiv:2411.08508v4} proposes novel regularization terms that lead to a 17x storage reduction, but the primary rendering speed remains bottlenecked by the sheer number of primitives. The system-level optimization is increasingly focused on reducing the number of Gaussian primitives themselves, a synergy between algorithmic and system co-design.

In the training phase, performance and memory efficiency are tightly coupled with primitive count and rendering resolution. DashGaussian \cite{arxiv:2503.18402v2} provides a significant leap by scheduling the optimization complexity. It formulates 3DGS optimization as progressively fitting higher frequency components, proposing a dynamic rendering resolution scheme alongside a synchronized growth of Gaussian primitives. This co-optimization cuts training time by 45.7\% without sacrificing quality, establishing a potent framework for efficient training. For large-scale scenes, the hierarchical 3D Gaussian representation \cite{arxiv:2406.12080v1} introduces a Level-of-Detail (LOD) structure, which trains the scene in chunks and consolidates them into a hierarchy, allowing for an efficient LOD solution for rendering distant content and enabling real-time rendering of trajectories up to several kilometers. Similarly, the LODGE approach \cite{arxiv:2505.23158v2} applies a depth-aware 3D smoothing filter followed by importance-based pruning to construct hierarchical LOD levels, demonstrating the need for dynamic selection of Gaussian subsets based on camera distance to reduce rendering time and GPU memory. This theme of compaction is further reinforced by \cite{arxiv:2503.16924v2}, which directly solves for a minimal set of primitives, achieving a 50\% storage reduction and a 600+ FPS rendering speed by re-parameterizing Gaussian attributes to be more robust to lossy compression.

Looking toward emerging trends, the ability to minimize both memory and computation from the outset is critical. Matryoshka Gaussian Splatting \cite{arxiv:2603.19234v2} presents a training framework that learns a single ordered set of Gaussians, where rendering any prefix yields a coherent, lower-fidelity reconstruction, effectively a continuous performance-computing trade-off. This capacity demonstrates its full performance and is useful for scalable deployment on resource-constrained devices. The frontier of this co-design is moving towards direct and efficient feed-forward scene generation. For instance, the work in \cite{arxiv:2603.25745v1} proposes the LGTM framework that predicts compact Gaussian primitives coupled with per-primitive textures, decoupling geometric complexity from rendering resolution. This enables high-fidelity 4K synthesis without per-scene optimization, which is a radical paradigm shift away from the per-scene training overhead that dominates current methods. The system-level co-design of 3DGS is thus evolving from software-only rasterization tweaks to a broader co-optimization of the online representation and algorithmic scheduling, while the future is being shaped by the trend toward feed-forward, single-pass generation to bypass iterative densification, albeit often accompanied by reduced reconstruction quality. This trajectory will likely need to harness a joint approach that co-optimizes for a balance between performance and spatial complexity.


\section{Methods for Geometric Accuracy, Surface Reconstruction, and Handling Degraded Inputs}



\subsection{Surface-Driven Regularization and Geometric Prior Integration}


The optimization of 3D Gaussian Splatting (3DGS) under photometric loss alone frequently converges to geometrically implausible solutions, characterized by free-floating artifacts, redundant volumetric blobs, and a fundamental decoupling of the opacity field from the physical surface \cite{arxiv:2308.04079v1}. Surface-driven regularization counteracts this by injecting inductive biases that anchor primitives to the observed geometry, organizing the amorphous cloud of Gaussians into a coherent, topologically meaningful structure. The central challenge resides in the design of regularizers that enforce this geometric fidelity without constraining the colorimetric expressivity that defines 3DGS.

A dominant strategy for geometric anchoring involves the integration of explicit depth and normal priors derived from various sources, including monocular estimators, multi-view stereo (MVS), and active sensors. Regularizing the optimization with depth cues, as demonstrated by DN-Splatter \cite{arxiv:2403.17822v1}, rectifies the geometric drift observed in indoor, texture-poor scenes by penalizing the distance between the rendered depth and the known signal, while local normal smoothness encourages planar consistency in the Gaussian distribution. However, a myopic focus on depth is insufficient; for a correct structural coupling, alignment must be actively enforced on the orientation field. For instance, supervision on the Gaussian orientation field prevents the notorious "sunflower" artifact in 2D Gaussian Splatting, where flat disks orient facing the camera, thereby ensuring that each primitive corresponds to a physical surface rather than a viewpoint-dependent artifact \cite{arxiv:2403.17888v1}. While normal priors from monocular estimators can improve geometry, they are often noisy and require careful handling; therefore, a more robust approach involves densely propagating fused multi-view stereo cues to regularize the Gaussians globally, which provides stronger geometric constraints than sparse depth or monocular normals alone \cite{arxiv:2401.03890v2}.

To push beyond the limitations of 2D priors, a class of methods aligns the discrete Gaussian field with an implicit surface defined by a Signed Distance Function (SDF), creating a coupling that enforces both surface alignment and multi-view consistency. By binding Gaussians to the surface of a mesh, such as in the SuGaR regularization framework, the optimization landscape is enriched with an Eikonal constraint that ensures distance field properties are respected, while a projection-based consistency loss is used to tightly couple the Gaussians to the zero-level set, allowing the extracted geometry to faithfully match the scene \cite{arxiv:2311.12775v3}. However, this coupling is rarely sufficient for completeness. To enforce near-surface occupancy and suppress stray artifacts, volume rendering is aligned with the signed distance field, which uses surface rendering of the Gaussians to anchor the opacity field to the level-set \cite{arxiv:2404.10772v1}. Similarly, MeshGS exploits the distance between each Gaussian and a known mesh to distinguish tightly-bound splats from loosely-bound ones, using this per-splat distance to regularize scene geometry and prune redundant volumetric blobs, effectively reducing the number of Gaussians while improving rendering accuracy \cite{arxiv:2410.08941v1}. Rather than synchronizing with a fixed external field, other methods derive the surface prior from the evolving Gaussians themselves, creating a stable, self-improving feedback loop that does not require external ground truth \cite{arxiv:2401.03890v2}.

A different line of attack for surface refinement focuses on the Gaussian covariance and shape primitives to ensure they represent flat, surface-like geometries rather than needle-like artifacts. For instance, by regularizing the Gaussian density distribution through geodesic distance on continuous quadrics, the alignment of the flattened Gaussian's normal with the surface normal is enforced, and the resulting spectrum of positive semidefinite shapes prevents the collapse into 1D lines \cite{arxiv:2411.16392v4}. More generally, decomposing the covariance matrix to avoid extreme spectral ratios regularizes the Gaussians toward disk-like shapes, and promotes intrinsic alignment with the local surface \cite{arxiv:2403.17888v1}. This is supported by secondary side-effects of using planar primitives, since 2D Gaussians and textured planar billboards preserve view consistency more naturally than volumetric ones, as textures and alpha-maps blend across surfaces.

Finally, the stability of the reconstruction can be further improved by explicitly targeting the Gaussian ordering and spatial distribution. A structural reorganization that requires k-nearest-neighbor state update to regularize the Gaussian positions, combined with a spatial coverage strategy, accelerates the convergence by reducing solution space dimension \cite{arxiv:2411.12788v2}.

In summary, the field is shifting toward a robust, multi-cue geometric prior anchored in depth, normal, and implicit surface constraints that create a consistent expectation for what the final rendered model should be. Future approaches for surface refinement should consider combining discrete Gaussian splats with continuous SDF representations in a principled way, using uncertainty estimates to weigh the contribution of each geometric cue in truly unconstrained and degraded settings \cite{arxiv:2401.03890v2,arxiv:2403.17888v1}.


\subsection{Surface-Adapted Primitive Kernels and Novel Basis Functions}


The fidelity of surface reconstruction in 3D Gaussian Splatting is fundamentally constrained by the choice of primitive kernel. The canonical 3D Gaussian, while computationally efficient for rasterization, possesses radial symmetry and unbounded support that are poorly matched to the characteristics of physical surfaces---planar in their local parameterization and finite in their extent. This representational mismatch leads to volumetric "fuzz" at object boundaries, multi-view inconsistent opacity, and an over-reliance on dense overlapping primitives to approximate sharp features through superposition. A substantial body of work has emerged to address this, converging on a central insight: the \emph{shape} of the primitive must be aligned with the \emph{geometry} being represented. This reshaping is not merely a cosmetic improvement; it is the critical mechanism for eliminating floaters, enabling direct mesh extraction, and achieving compact scene representations with fewer primitives, thereby laying the representational foundation upon which the surface-driven regularization discussed in the preceding subsection can be effectively applied. This subsection systematically examines these surface-adapted kernels, tracing a progression from collapsed 2D formulations to increasingly expressive higher-order bases, and analyzes the trade-offs each introduces between geometric fidelity, rendering efficiency, and optimization stability---also served by the regularization class of methods presented above.

The most direct strategy pulls the kernel toward the surface by collapsing the volumetric Gaussian into a planar configuration. The foundational work of \cite{arxiv:2403.17888v1} (2DGS) parameterizes an oriented flat Gaussian disk and derives a perspective-accurate ray-splat intersection algorithm for differentiable rasterization. This formulation intrinsically matches the local planarity of surfaces and directly enables detailed mesh extraction from the radiance field. The geometric advantage of 2D surfels is significant, yet concerns remain regarding the fairness of optimization in high-texture regions where the disk orientation may be poorly constrained. Concurrently, \cite{arxiv:2411.08508v4} (BBSplat) extends this direction by introducing learnable textured planar primitives with RGB textures and alpha-maps. This directly addresses the limitation of monochromatic Gaussian kernels in representing locally-varying appearance, and notably demonstrates improved storage efficiency through sparse texture regularization.

Building upon this surface-aligned philosophy, a further line of research enhances the representative capacity of the underlying planar kernel by introducing curvature as well as a second-order analysis. The \cite{arxiv:2411.16392v4} paper enables the formal transition to quadric surface primitives (e.g., paraboloids) rather than flat planes. The paper completely has a re-derivation of the density distribution using geodesic distances, a more natural departure from the Euclidean distance metric that is the foundation of classic Gaussian kernels, and ensures that the radiance field remains surface-aware across the anisotropic regime. The approach implicitly generalizes the principles behind the previous work to 3D regions, implementing a second-order geometry. This can be seen as a natural extension of the previous surfaces, resulting in significant improvements, such as a 33\% reduction in Chamfer distance over 2DGS on the DTU dataset shot. Moreover, the explicit and compact nature of the parameterization opens the door to hybridizing geometric primitives with feature-descriptive semantics, directly relevant to downstream computer vision tasks beyond novel view synthesis \cite{arxiv:2508.09977v5} (\# acknowledged in the following sections).

Despite the advances of the planar and curved primitives, capturing sharp high-frequency details such as thin structures, discontinuities, and hard edges remains a primary challenge, and it is usually ill-addressed by the raw primitive alone. Naively increasing the number of primitives in these regions is memory-inefficient and sometimes breaks the optimization with repetitive identical structures. A promising complementary direction adapts the granularity and the shape of the kernel in response to the local structure, allocating Gaussians based on visual sensitivity \cite{arxiv:2506.12400v2}. A different strategy uses 3D smooth convex primitives \cite{arxiv:2411.14974v3} that can model discontinuities and edges more naturally than \emph{either} smooth 3D Gussian \emph{or} 2D Gaussian, and demonstrating superior fidelity with fewer primitives on standard benchmarks.

Synthesizing these advances, a unifying perspective emerges: the trend is not toward a single "best" primitive but toward \textbf{the primitive shape as a learned, scene-adaptive parameter}. Rather than using a single type of base kernel (2D disk, quadric, or volume), an optimal system could dynamically and continuously morph the representation based on the local appearance and adjacency geometry. For instance, a scene could be represented as a mixture of 2D disks for flat regions, 2D quadrics for ridges, and specialized thin primitives for edge lines. This would necessitate fundamental changes to the differentiable renderer, including the explicit integration of shape assumptions at the splatting stage itself. Furthermore, the interaction between these new kernel structures and the accuracy of the projected splat remains largely unexplored. As the field moves toward resolution-optimal representation, the integration of surface-adapted kernels with feed-forward reconstruction \cite{arxiv:2604.03069v1} and MCMC-based refinement \cite{arxiv:2404.09591v3} promises to not only further improve the geometry but also adopt a more rigorous, physically grounded basis for real-time modeling.

However, the choice of the adapted kernel is necessary, but not by itself sufficient, to guarantee that the optimization converges to the correct surface. A visually plausible radiance field can still be generated by a set of loosely coupled volumetric blobs, rather than by a faithful arrangement of primitives anchored to the physical scene. This is precisely the point where the surface-driven regularization as analyzed here breaks in---it injects geometric priors and structural constraints that bind the chosen primitive to the observed geometry. While the kernel matrices provide high representational capacity to match local shape, regularization stabilizes the combination of the geometry by fusing external cues---depth, normals, or implicit surface fits---into the optimization. In a broader view: the surface kernels redefine \emph{what} a Gaussian can represent, and the regularizers decide \emph{where} or \emph{how well} it truly is eventually placed and anchored. Together, these two complementary components form the backbone for the next step of extracting coherent, high-fidelity surfaces, meshes, and other scene representations discussed in the following subsection.


\subsection{Robust Reconstruction from Sparse, Unposed, and Degraded Inputs}


The reconstruction of high-quality 3D scenes from sparse, unposed, or degraded captures represents a fundamental departure from the well-posed dense-view setting, where abundant overlapping imagery and reliable Structure-from-Motion (SfM) initialization underpin the success of vanilla 3D Gaussian Splatting. In such constrained regimes, the optimization landscape becomes severely ill-conditioned: photometric loss alone is insufficient to constrain Gaussian placement, leading to co-adaptation artifacts where primitives become overly entangled to overfit training views at the expense of novel-view appearance \cite{arxiv:2508.12720v3}. Furthermore, the reliance on SfM for initialization breaks down entirely when camera poses are unknown or when input images exhibit photometric inconsistencies, motion blur, or transient occlusions, which together exacerbate degenerate local optima and floater formation. This subsection synthesizes methodological advances spanning initialization, regularization, and joint pose-geometry optimization that collectively address the dual challenge of conditioning the optimization and imposing structure from limited observational data.

A primary research thrust seeks to replace fragile SfM seeding with learned priors. The use of dense multi-view stereo or feed-forward networks to generate pointmaps has emerged as a robust alternative to feature matching. Building on the capabilities of foundation models, generalizable frameworks \cite{arxiv:2410.06245v1} that construct hierarchical 3D Gaussians in a coarse-to-fine manner via an Error-Aware Module and a Modulating Fusion Module have demonstrated that per-pixel prediction of Gaussian parameters in a single forward pass can achieve robust and detailed reconstruction from only two viewpoints, effectively including per-scene dense optimization. Extending this paradigm, EDGS proposes a triangulation of dense pixel correspondences from multi-view images to initialize a dense geometry approximation, eliminating the incremental densification process and thereby shortening optimization while preserving high-frequency details \cite{arxiv:2504.13204v2}. Moreover, we observe a growing body of work emphasizing the importance of statistical regularizers. For instance, a Gaussian Process (GP) based densification, as introduced in GP-GS, treats Gaussian point cloud synthesis as a continuous regression problem, using GP-predicted uncertainty to filter unreliable predictions and improve noise reduction in sparsely initialized scenes \cite{arxiv:2502.02283v6}.

To mitigate overfitting when observations are scarce, a suite of regularization and supervision strategies has been collected. The Depth-and-Density Guided Gaussian Splatting (D2GS) method identifies dual failure modes of overfitting near the camera due to excessive Gaussian density, and underfitting in distant regions---and proposes a density-guided dropout mechanism, coupled with distance-aware fidelity enhancement, to stabilize distributions \cite{arxiv:2510.08566v1}. A complementary perspective is provided by the concept of co-adaptation, which quantifies the statistical entanglement among Gaussians via pixel-wise variance under random subsetting; based on this metric, regularization techniques such as random Gaussian dropout and multiplicative opacity noise are shown to effectively improve rendering stability while remaining plug-and-play across different 3DGS variants \cite{arxiv:2508.12720v3}.

Addressing the more challenging sparse setting, recent works have integrated generative priors to fill in unobserved regions. Diffusion models trained for artifact repair and inpainting are generated to provide supervision signals for under-constrained areas, with GSFixer integrating both 2D-semantic and 3D-geometric features from reference views to enhance 3D consistency during the fixing of artifact novel views \cite{arxiv:2508.09667v1}. More recent work by the neighborhood decomposes the process into repairing visible regions and hallucinating missing areas via two separate models, highlighting the advantages of separating these tasks in extremely sparse settings \cite{arxiv:2503.10860v2}. These deep generative regularizations effectively reduce corruption in the ill-posed setting by leveraging 2D image generators.

Looking ahead, a decisive challenge lies in establishing robust, physics-aware probabilistic metrics to quantify the reliability of reconstruction outputs for downstream decisional tasks, going beyond simple rasterization error checking. Future research will likely focus on unifying these generic priors with a principled uncertainty quantification, enabling mechanisms for perceptive active view selection and ensuring robust confidence attributable to the not-posed or partially observed spatial regions.


\subsection{Handling Photometric Inconsistencies, Transient Occlusions, and Dynamic Degradations}


Real-world captures routinely violate the idealized assumptions of static geometry and photometrically consistent appearance that underpin vanilla 3D Gaussian Splatting (3DGS). Fluctuating illumination, specular reflections, exposure variations across viewpoints, motion-blurred frames, and the presence of transient occluders (e.g., pedestrians, vehicles) constitute significant sources of corruption that, if unaddressed, cause the optimization to conflate appearance artifacts with geometric structure, producing floaters, fragmented surfaces, and compromised depth estimates. Research in this area therefore moves beyond the na\"{i}ve photometric loss toward a more robust image formation model that decouples static, view-independent scene properties from dynamic, appearance-corrupting, or transient elements. This pursuit of robustness in appearance modeling directly complements the sparse-view and initialization challenges discussed in the previous subsection, where photometric inconsistencies exacerbate the ill-conditioned optimization landscape; conversely, the robust decoupling strategies developed here provide the crucial radiometric fidelity upon which the physically grounded surface reconstruction methods of the following subsection can be reliably built.

A dominant strategy for handling photometric inconsistencies is to introduce per-view or per-primitive appearance modeling into the optimization loop, effectively factorizing global scene albedo from view-dependent photometric distortions. For instance, methods like Luminance-GS address extreme low-light, overexposure, and varying camera exposure settings by adopting a per-view color matrix mapping and a view-adaptive curve adjustment mechanism within the 3DGS framework, significantly improving reconstruction quality where multi-view photometric inconsistencies are severe \cite{arxiv:2504.01503v2}. Similarly, work on illumination-agnostic rendering, such as LITA-GS, introduces an illumination-invariant physical prior extraction pipeline and a scene structure rendering strategy that decouples geometry optimization from corrupted radiance, effectively preventing the optimization from being misled by inconsistent appearance across views \cite{arxiv:2504.00219v1}. Extending this paradigm, appearance embeddings have been incorporated into SLAM pipelines to model appearance variations from different camera poses, allowing the 3D Gaussians to maintain photorealistic mapping quality under changing illumination and exposure conditions \cite{arxiv:2501.05242v3}. These techniques collectively ensure that the photometric error used for gradient backpropagation does not corrupt static geometry, thereby preserving high-frequency details and providing a cleaner signal for the geometric anchoring methods discussed in the next subsection.

Handling transient objects and occlusions requires the optimization to be robust to pixels that are explained by moving entities rather than by the static background. A common approach photometrically generates confidence or visibility masks that down-weight or eliminate the spatial loss from transient regions. For instance, SkySplat, addressing multi-temporal satellite imagery, exploits a cross-self consistency masking and a consistency-based aggregation strategy to mitigate the interference of transient objects (e.g., cloud shadows) with reconstruction \cite{arxiv:2508.09479v2}. The underlying principle is that static Gaussians remain dominant, and inconsistency across views can be detected as a mask, preventing the optimization from repeatedly fitting a distractor that should not be anchored in the canonical space. This distinction between persistent and transient geometry is crucial for accurate, drift-free static mapping, and it parallels the sensor-aware fusion and uncertainty-aware pruning strategies of the following subsection, where such masks can be reinterpreted as cues for geometric fidelity.

Motion blur and non-uniform sensor distortions---such as rolling shutter effects---also cause substantial disregard that typical datasets ignore. As capture scenarios continue to scale in traversal length or duration, directly modeling the image formation process during exposure becomes necessary rather than treating each frame as an ideal pinhole image. These approaches integrate dynamic trajectory or motion parameters into the Gaussian optimization, establishing a single differentiable flow that co-optimizes clean 3D geometry and undoes motion-induced aberrations. This advancement in robustness to dynamic degradation is a critical component for elevating 3DGS from well-conditioned controlled captures to intervention-free outdoor scans, and it naturally bridges to the physical-surface reconstruction objectives of the next subsection, where geometry must be robustly recovered under unpredictable acquisition conditions.

Looking forward, several key research directions emerge. A primary challenge is linking robust appearance decoupling to stronger geometric supervision. For instance, techniques that allow explicit per-Gaussian appearance re-parameterization could be coupled with prior-free geometric cues for mutual enhancement. The path forward will likely be architectures that intrinsically support reliable depth, normal, or other structural priors inside the optimization loop, ensuring that the improving robustness can purely benefit geometry-oriented downstream tasks such as surface extraction and differentiable rendering of soil. Moreover, the issue of baseline inconsistency in transient or heavily degraded regions remains open. While uncertainty-aware mask weights suppress the loss in a deterministic way, a more principled approach---probabilistic propagation of uncertainties across pixels through the splatting pipeline---could lead to more stable convergence and better performance in challenging regions. Research directions integrating task-specific semantic knowledge to not only detect but also explicitly capture transient objects will likely be next, moving the field beyond static asset recovery toward robust, interactive geometry estimation in dynamic scenes. Such methods should aim to learn object-specific, semantics-guided representations of transient content, which, when combined with the surface-anchoring primitives of the following subsection, would provide a comprehensive framework for digital twin building under real-world conditions.


\subsection{Physical Surface Constraining, Sensor-Aware Fusion, and Fine-Scale Detail Preservation}


The pursuit of physical surface fidelity in 3D Gaussian Splatting extends beyond photometric consistency, demanding the integration of explicit geometric priors and sensor-derived constraints to resolve ambiguities that appearance alone cannot address. This subsection addresses methodologies that anchor the floating Gaussian primitives to physically grounded surfaces through sensor-aware fusion, parameterized surface models, and fine-scale geometric refinement. The core challenge lies in reconciling the volumetric, radiance-centric nature of Gaussian primitives with the requirements of surface representation---normals, curvature, and watertight topology---which are essential for downstream applications in digital twins, fabrication, and mixed reality.

A foundational avenue of research integrates dense geometric priors and constraints directly into the Gaussian optimization objective. Methods such as GausSurf \cite{arxiv:2411.19454v2} leverage multi-view stereo (MVS) and normal priors from pre-trained networks to guide smooth geometry in texture-less regions while enforcing photo-consistency in rich ones. These approaches iterate patch-match refinement with Gaussian optimization to create a mutual reinforcement loop, which significantly improves the time complexity to convergence while enhancing final reconstruction quality. A more principled approach involves self-constrained priors derived from the current field itself. By fusing rendered depth maps into a truncated signed distance function (TSDF) grid and using this self-updating narrow band to cull Gaussians, reposition them closer to the surface, and adjust opacity in a geometry-aware manner, the optimization creates an adaptive feedback loop that progressively tightens the surface prior \cite{arxiv:2603.19682v2}. This framework exemplifies the philosophy of enforcing global topology without global optimization. Beyond Euclidean distance, recent formulations grounded geometry in a stochastic solid framework, deriving a rigorous theoretical connection between Gaussian primitives and physical solids, thereby using that foundation to render accurate depth maps and extract high-quality meshes \cite{arxiv:2601.17835v2}. For more structured ownership of geometry, parametric model integration becomes viable, e.g., Gaussian splatting coupled with mesh template decoration \cite{arxiv:2512.09162v3}, where separable Gaussian primitives are attached to a mesh template in UV space, which introduces textural editability and visual consistency compliance with physical shape.

In parallel, surface-adapted primitive kernels offer a more literal interpretation of surface geometry capture. Quadratic Gaussian Splatting (QGS) replaces planar surfels with quadric surfaces (ellipsoids, paraboloids), promoting a geodesic distance-evaluated density distribution that is inherently surface-adaptive under deformation, reducing the need for many Gaussians per curved region \cite{arxiv:2411.16392v4}. This second-order kernel lifts the representation beyond piecewise-linear approximations. The incorporation of learnable textures transitions the Gaussian from a single-color entity to a more expressive planar billboard that encodes micro-scale appearance variations. Work on Gaussian billboards \cite{arxiv:2412.12734v1,arxiv:2411.08508v4} shows that per-splat texture interpolation reduces over-reconstruction while retaining high fidelity, and when combined with 2DGS geometry, it closes the gap between mesh-like appearance and inverse-rendered detail. To maximize perceptual impact, billboard tessellation alongside depth sorting and texture interpolation allows faithful high-frequency detail preservation \cite{arxiv:2412.01823v1}. Beyond the planar billboard approach, the introduction of 2D triangle primitives into Gaussian-like rendering \cite{arxiv:2506.18575v3} builds a differentiable bridge to a mesh-worthy, compact topology, significantly easing the direct extraction to watertight surfaces by annealing compactness and reconstructing semi-transparent geometry into fully opaque triangle faces.

Sensor-aware fusion represents a crucial axis for robustness in large-scale, physically demanding scenarios. MetroGS \cite{arxiv:2511.19172v4} serves as a pivot point: it addresses the complication of sparse coverage through a system built upon a distributed 2DGS backbone, then upgrades it with structured dense initialization and a pointmap model (SfM priors), merging monocular and multi-view geometry via a progressive hybrid optimizer. Recognizing appearance inconsistencies that plague city scenes, it deploys depth-guided appearance modeling that collects features with 3D spatial coherence---effectively decoupling lighting from geometry, giving a stable, architecturally meaningful surface in faint regions. This explicit modeling of appearance--geometry decoupling under a Gaussian field serves as a foundational approach for leveraging sensor geometry beyond mere color agreement. This consistency is furthered by geometric refining technologies using normal and curvature information; for example, smoothing splatting kernels on curved surfaces. Approaches such as a robust Laplace-Beltrami formulation on the geometry-informed Gaussian centers allow a geometry-sensitive differential operator to extract refinement signals that suppress outliers in the potential field distribution, using covariance weights to guide geometric detail leveling \cite{arxiv:2502.17531v2}. When arranged with the availability of location-specific sensor priors, these compliance mechanisms enforce geometric measurement correctness on the field and allow these models to adapt to the restrained dynamic topology of complex 3D shapes.

In summary, this subsection captures a steady progression from modeling Gaussian splats as semantic inputs in coarse geometry refining functions and anchored decisions---e.g., integrating SDFs but no watertightness with hybrid optimization---toward parameterized surface kernels (e.g., quadrics and triangles), and finally toward exploiting automatic field-controlled adjustments during reconstruction, robust to centimeter- to meter-level scales, as exemplified in large-scale urban settings. This convergence of geometric physics and radiometrically validated appearance allows Gaussian Splatting to move beyond render-synthesis toward tangible, sensor-consistent, edit-ready surface representation. Future directions will likely use semantically labeled remote sensing data, curvature-aware densification control, and learned TSDF rasterization to forge a seamless and robust connectivity between the explicit point-based realm and geometrically grounded model factories.


\section{Compression, Scalability, and Hardware-Aware Deployment}



\subsection{Model Compression and Storage Optimization}


The rapid adoption of 3D Gaussian Splatting (3DGS) for high-fidelity novel view synthesis has exposed a critical bottleneck: its substantial memory footprint. A typical scene, represented by millions of Gaussian primitives each encoding position, covariance, opacity, and spherical harmonic coefficients, can demand hundreds of megabytes to gigabytes of storage \cite{arxiv:2308.04079v1}. This cost severely restricts deployment on edge devices, mobile platforms, and streaming pipelines, making model compression an indispensable area of research. A taxonomy of approaches reveals three dominant, complementary strategies: sparsification to reduce the primitive count, quantization and entropy coding to reduce bytes per attribute, and predictive methods to share parameters amongst primitives.

The first line of attack, sparsification, aims to identify and remove Gaussians that contribute negligible information to the final render. Early work demonstrated that a significant fraction of the original point cloud is redundant; simple pruning can remove a large portion of primitives while maintaining or even improving baseline performance \cite{arxiv:2412.00578v3,arxiv:2410.23213v1}. This redundancy arises from the adaptive densification heuristics, which often spawn primitives that are quickly rendered redundant by their neighbors. Beyond simple pruning, more sophisticated criteria have emerged. For instance, the MCMC-based formulation reframes the densification process as a sampling problem, allowing gradients and opacity to guide the removal of non-contributing samples \cite{arxiv:2404.09591v3}. Similarly, structure-aware distribution methods, have categorize Gaussians into edge-defining primitives and volume-filling primitives, enabling specific pruning and retraining of the latter to achieve a more structure-aware compactness \cite{arxiv:2501.13045v1}. This highlights a crucial trend: less emphasis on blind attribute compression and a growing exploitation of the fact that in a well-optimized, dense field, a smaller, minimally optimal set of Gaussians is sufficient, effectively reducing computational costs for subsequent processing (e.g., 600+ FPS in \cite{arxiv:2503.16924v2}).

Sparsification alone is often insufficient; the remaining primitives still carry large attributes, particularly for view-dependent color representation. To address this, attribute compression through quantization and entropy coding has proven extremely effective. Early work introduced vector quantization (VQ) based on K-means to cluster similar appearances, yielding an initial 40-50x reduction \cite{arxiv:2311.18159v3}. This principle has been significantly refined, but with a recent methods demonstrating that high-precision attributes are essential when the primitive count is low. The Optimized Minimal Gaussians method emphasizes the use of high-precision attributes on a minimal primitive set, and demonstrates that lossy attribute compression. Engineering an improper results a severe quality degradation, motivating them to use sub-vector quantization to maintain quality, reducing storage nearly 50\% compared to previous state-of-the-art \cite{arxiv:2503.16924v2}. The entropy coding pushes this concept further. By training a differentiable estimator of the bitrate, models can be optimized for rate-distortion, forcing the underlying representation to be as small as possible, and typically evaluated across a range of compression levels \cite{arxiv:2501.13558v3}. The MEGA framework for dynamic scenes shows how such frameworks can be extended to 4DGS, jointly optimizing for memory and storage with entropy-constrained deformation and a compact alternating current color predictor, achieving storage reductions of over 190x \cite{arxiv:2410.13613v3}.

A parallel strategy exploits the spatial coherence and redundancy inherent in scenes: the use of predictive representation. Rather than encoding each Gaussian independently, these methods aim to share parameters among neighboring primitives. For example, in Sketch and Patch, the representation naturally organizes primitives into two distinct categories, storing parameters for edges and densely and coarsely for surfaces \cite{arxiv:2501.13045v1}. This indicates a relationship with principles in image compression where removing from mutual redundancy is paramount. A similar goal is taken by hierarchical representations, which inherently parameterize level of detail for distant content \cite{arxiv:2406.12080v1}.

Looking forward, the field is converging on the need for unified frameworks that jointly optimize primitive count, parameter representation, and bitrate, and that natively support progressive level-of-detail. Models that scale between quality and memory on the fly, without auxiliary per-level fine-tuning \cite{arxiv:2501.13558v3}, are crucial for adaptive streaming in XR and telepresence \cite{arxiv:2412.06257v2}. Another open challenge is the systematic evaluation of quality, which currently lags behind reconstruction benchmarks; recent work, such as the MUGSQA dataset, begins to fill this gap with subjective quality that vary across viewing distances and input conditions \cite{arxiv:2502.13196v2}. Ultimately, compression is moving shift from "storage of a static file" to "dynamic resource management", enabling a single representation to provide a range of fidelity across with bandwidth and latency constraints, always valuable for production deployment as well as visibility for many other areas such as real-time rendering on low-power devices or streaming in XR.


\subsection{Hierarchical, Structured, and Adaptive Scene Frameworks for Large-Scale Rendering}


As 3D Gaussian Splatting scales from object-centric captures to city-scale environments, the flat, globally optimized primitive set of the vanilla framework becomes a bottleneck in three interrelated dimensions: memory footprint, rendering throughput, and optimization convergence. Hierarchical, structured, and adaptive scene frameworks address these limitations by reorganizing the primitive distribution, introducing multi-resolution or spatially partitioned representations, and adapting the allocation of computational resources based on scene content or viewing conditions. These approaches collectively move beyond the local, heuristic-driven densification strategies that dominate small-scale reconstruction \cite{arxiv:2505.05587v1,arxiv:2411.12788v2} toward globally coherent, content-aware, and hardware-conscious scene organizations. In this context, the compression techniques discussed earlier---sparsification, quantization, and predictive coding---operate on an already optimized primitive set, whereas the hierarchical and structural methods presented here fundamentally reshape that primitive set \emph{a priori}, making them a complementary yet distinct layer of efficiency that directly informs the deployment-side optimizations examined in the following section.

A dominant strategy in this direction is \textbf{anchor-based and coarse-to-fine hierarchies}, which decouple scene structure from local detail. Instead of treating each Gaussian as an independent optimization variable, these methods introduce a hierarchy of control points---or anchors---that modulate the position, scale, and appearance of a spatially associated set of Gaussians. This hierarchical formulation reduces the effective dimensionality of the optimization problem and improves robustness to sparse initialization. The benefits of this paradigm extend to feed-forward and generalizable settings as well: HiSplat \cite{arxiv:2410.06245v1} generates coarse Gaussians to capture low-frequency structure and fine Gaussians to recover texture details, using an error-aware compensation module to promote inter-scale interaction. The implicit assumption---that high-level geometric anchoring can serve as a well-behaved and globally consistent latent space---has proven effective in reducing both parameter count and training footprint while preserving rendering quality.

Complementary to hierarchy is \textbf{spatial partitioning and superpoint aggregation}. Large-scale scenes are rarely statistically homogeneous; uniform primitive allocation thus becomes a direct source of memory waste and computational overhead. To address this, methods such as Mini-Splatting2 \cite{arxiv:2411.12788v2} reorganize Gaussian primitives through spatial aggregation (often via k-nearest neighbors) into superpoint-like structures, decoupling primitive center positions from their joint orientation and scale. The resulting representation exhibits better spatial regularity and enables up to 4$\times$ fewer primitives with a 3$\times$ faster optimization, while maintaining state-of-the-art visual quality. This geometry-driven paradigm is an important counterpoint to the photometric-only supervision of the vanilla 3DGS, which tends to produce irregular primitives and erratic convergence. Notably, these structural improvements at training time naturally complement the compression techniques of the previous section by reducing the number of primitives that higher-level coding methods must subsequently handle.

Adaptive frameworks further refine the allocation of computational effort by identifying which regions of the scene actually require additional primitives in real time. For instance, ReSplat \cite{arxiv:2510.08575v3} reduces render scale in a 16$\times$ subsampled space and performs iterative updates conditioned on scene feedback, effectively reallocating Gaussians to underperforming views. Likewise, Perceptual-GS \cite{arxiv:2506.12400v2} endows Gaussians with interpretable guidance, allocating finer Gaussian granularity to visually non-uniform regions while suppressing dense allocation in uninformative areas. The "Off-the-Grid" approach \cite{arxiv:2512.15508v2} further advances this idea by detecting primitives at the sub-pixel level, assigning a varying number of primitives per patch based on complexity. These methods demonstrate that carefully designed efficiency schemes not only reduce computational cost but also improves output reliability by anticipating implementation beyond simple photometric cues. This dynamic, content-aware resource allocation is precisely the kind of runtime flexibility that hardware accelerators and kernel-level optimizations---examined in the next subsection---rely upon to achieve maximal throughput on constrained devices.

Looking forward, several key challenges remain. First, current adaptive structures are often heuristic-driven, lacking a formal validation of the relationship between allocated primitive density and the spectral characteristics of the underlying signal, which limits principled design choices. Second, these frameworks have yet to demonstrate robust generalization to fully dynamic and unbounded scenes at scale. The convergence of feed-forward prediction models \cite{arxiv:2312.12337v4}, self-supervised pose-free learning \cite{arxiv:2508.01171v2}, and structured representation suggests a promising path forward: a future scene model that remains globally context-aware, physically consistent, and can be traversed and rendered at interactive rates regardless of scale. As the field progresses, the distinction between "hierarchy for compression" and "hierarchy for semantic understanding" will likely vanish, with structure serving the dual purpose of efficient rendering and high-level scene understanding---a convergence that will equally benefit the co-designed hardware and runtime systems discussed in the next subsection.


\subsection{System-Level Hardware Co-Design and Real-Time Optimization}


The pursuit of real-time and energy-efficient 3D Gaussian Splatting (3DGS) on resource-constrained platforms---from AR/VR headsets to embedded systems---necessitates a fundamental departure from generic GPU kernel implementations toward tight algorithmic-hardware co-design. This interplay spans the entire rasterization pipeline, demanding innovations at multiple levels: kernel-level optimizations, workload restructuring, and dedicated hardware architectures. The original tiled rasterizer, while efficient, exhibits significant memory access bottlenecks and redundancy in per-pixel computation, prompting several foundational software engineering advances. A primary direction has been the development of highly optimized CUDA kernels that fuse operations, reduce global memory traffic, and exploit modern parallel architectures. These efforts, while critical, are often incremental. A more transformative approach is the principled restructuring of the rasterization algorithm itself. For instance, the per-pixel sorting and alpha-blending of overlapping Gaussians, a sequential bottleneck, is entirely bypassed in order-independent implementations that employ stochastic rasterization, achieving orders of magnitude reduction in compute load as demonstrated in the design of non-sorting pipelines \cite{arxiv:2508.12720v3,arxiv:2402.00525v1}. Conversely, hierarchical sorting strategies \cite{arxiv:2402.00525v1} provide a compromise: they incur a negligible runtime overhead for view consistency but enable a 1.6$\times$ overall speedup by allowing for a more compact representation, illustrating the deep coupling between rendering algorithm and computational cost.

The key strategic insight emerging from recent work is that rendering is not a monolithic but an optimizable process. A primary optimization avenue is the reduction of the active primitive count. For example, the necessity of computationally expensive occlusion culling has led to algorithms that integrate low-cost, learnable or proxy-based depth estimates from any view. This proxy-guided culling reduces the number of active Gaussians by 2.5$\times$, accelerating the rasterization stage while also informing the densification process to avoid over-reconstruction in occluded regions \cite{arxiv:2509.24421v5}. At the training/inference system level, fine-grained scheduling of workloads based on computational complexity---for example, deciding at which resolution to render images or when to densify the primitive set---has proven immensely effective. DashGaussian exemplifies this by co-scheduling resolution and primitive count, reducing the computational complexity by up to 45\% without sacrificing quality \cite{arxiv:2503.18402v2}. This frequency-aware, dynamic optimization continually negotiates the trade-off between computational resource expenditure and fidelity, improving upon static optimization. Replacing generic optimizers with second-order methods also contributes to system-level efficiency, leveraging Hessian sparsity to reach convergence in fewer, more computationally intensive steps, resulting in a 4$\times$ speedup in low-primitive-count regimes \cite{arxiv:2504.12905v3}.

The ultimate manifestation of this co-design philosophy is the departure from generic GPU execution to task-specific hardware accelerators. Custom ASIC designs often integrate hardware-friendly compression techniques, which are not an option but a requirement for memory efficiency. Learning-based compression (e.g., codebooks and quantization) has been shown to reduce memory footprints by up to 27$\times$ while maintaining quality \cite{arxiv:2406.17074v1}. A larger bottleneck for practical deployment on edge devices, however, concerns the static and rendering memory consumption of the model itself. Frameworks like MEGS$^{2}$ tackle this by jointly optimizing the total primitive number and the parameters per primitive, using lightweight spherical Gaussians to replace heavy Spherical Harmonics, achieving a 50\% static and 40\% rendering VRAM reduction while maintaining comparable rendering quality \cite{arxiv:2509.07021v3}. For accelerators, integrating such memory-aware primitives early in the design pipeline is critical. Thus, the pathway to 1000+ FPS rendering lies in co-design of efficient compression schemas, with the rasterization algorithm, and dynamic scheduling policies, thereby creating a highly specialized and yet flexible system to handle the complexity of varying scenes.


\subsection{Large-Scale GPU Memory Management and High-Performance Training Scalability}


The exponential growth in Gaussian primitive counts---often reaching tens of millions for large-scale scenes---has exposed the GPU memory wall as the primary bottleneck in 3DGS scalability. The optimization objective requires simultaneous storage of Gaussian attributes, optimizer moments (e.g., Adam's first and second moments), and intermediate gradient tensors, resulting in memory footprints that far exceed the on-board capacity of even high-end accelerators. The fundamental challenge lies in the inherent tension between the computational throughput of GPU-accelerated rasterization and the architectural constraint of finite on-board memory, motivating a wave of system-level innovations designed to decouple model scale from hardware memory limits. Critically, the strategies emerging from this section not only address the immediate scalability bottleneck but also lay the algorithmic groundwork for the inference-time efficiency and memory-aware deployment examined in the next subsection: training procedures that are inherently leaner in memory and computation produce models that are naturally better suited for resource-constrained deployment.

\textbf{Hybrid Offloading and Asynchronous Pipelining Systems.} A dominant strategy is the hybrid offloading of model parameters and optimizer states from GPU to host memory. In this paradigm, GPU memory hosts only a transient working set---the primitives and associated state required for the current iteration---while the majority of the data remains in CPU or unified memory. The success of this approach hinges on the seamless orchestration of data transfers and computational kernels. Concurrent works demonstrate that this full-parameter offloading approach redistributes resource utilization at the system level: the GPU's computational units are saturated while CPU-side memory controllers and the network fabric become the final schedulers of the pipeline \cite{arxiv:2410.23213v1}. The systems decouple the optimization and simulation steps from the offloading process, effectively maintaining high GPU utilization despite the data transfer overhead. This architecture has achieved reductions of 3.3-5.6$\times$ GPU memory consumption, enabling the training of scenes with millions of primitives on consumer-grade hardware without incurring prohibitive training time overheads.

\textbf{Workload-Aware Scheduling and Subset Selection.} Beyond these offloading, a complementary family of techniques focuses on algorithmic alterations to the dataflow that reduce the volume of information processed, thereby implicitly lowering the memory state. The insight underlying these "subset-based" methods is that conventional dense gradient descent over millions of parameters is fundamentally wasteful; regions of the scene either converge to a stable appearance or are inside a poorly conditioned region where many sampled pixels provide redundant or adversarial gradients. Workload-aware approaches therefore reorder optimization steps to sample informative pixels and views selectively \cite{arxiv:2503.18402v2}. In large-scale training, this is reinforced by the observation that a significant degree of Gaussian creation indeed incurs a cost for primitives that are subsequently pruned, yielding a 10x reduction in effective training workload. A more theoretical framework is then provided by scheduling of the optimization process, where the rendering resolution and primitive growth schedule are jointly controlled in accordance with the current stage of the learning process to strip redundant training without degrading final quality. These scheduling principles extend naturally to the inference-time runtime systems discussed in the following subsection, where workload-aware resolution and primitive-count strategies re-emerge as core deployment-side optimization levers.

\textbf{Matrix-Free Second-Order Optimization.} In tandem with these systems-level and scheduling-level improvements, a parallel tract of work initiates a departure from standard first-order Adam optimization. The observation that the Jacobian matrix in 3DGS training exhibits extreme sparsity---each Gaussian only affects the pixels within its 2D footprint---opens the path toward matrix-free second-order optimizers. The Levenberg-Marquardt-style method, coupled with Conjugate Gradient solver, explicitly exploits this sparsity by constructing batched normal equations for individual parameter groups \cite{arxiv:2504.12905v3}. The isolation of individual Gaussians and their parameter groups further permits the splitting of the optimization problem at the parameter level, leading to the formation of small linear systems that can be solved entirely on GPU threads \cite{arxiv:2501.13975v2}. These methods produce training that converges significantly faster (2-5x speedups) while consuming less memory than their first-order counterparts, demonstrating the power of algorithmic system co-design for training efficiency. Moreover, the compact representation and faster convergence achieved by these second-order methods plants the seeds for the memory-efficient deployments and hardware-aware algorithmic choices that will be further amplified at inference time in the next subsection.

\textbf{Research Outlook and Bridge to Deployment.} The interplay of these four strategies---hybrid offloading, workload-aware scheduling, matrix-free second-order solutions, and system-level co-design---represents a comprehensive toolkit for large-scale training efficiency. The convergence across these stems from a shared thesis: the memory hierarchy is a controllable resource, and the optimal utilization of that resource requires a simultaneous design of the optimization procedure and the hardware runtime. Furthermore, the structures and principles introduced here---memory-efficiency, parameter reduction, and algorithmic pruning---do not remain confined to training alone; they create a seamless interface with the inference-time optimizations of the next subsection, where compressed primitives, accurate culling, and workload-aware scheduling re-emerge as the core mechanisms for achieving tractable rendering on resource-limited embedded systems. In future developments, the entire 3DGS ecosystem is likely to benefit from tight algorithmic-hardware co-design that links training-time memory management with inference-time efficiency, forming the foundation for scaling to hundreds-of-millions-and-beyond Gaussians, including the integration of heterogeneous memory and low-power edge GPUs.


\subsection{Temporal and Streaming Compression for Dynamic, Time-Varying, and Interactive-Experiences}


The efficient handling of dynamic, time-varying scenes represents one of the most demanding frontiers in 3D Gaussian Splatting (3DGS) deployment. While static scene compression has reached remarkable maturity---with techniques achieving orders-of-magnitude storage reduction through pruning, quantization, and entropy coding---the extension to 4D content introduces a fundamentally new dimension of redundancy and computational complexity. The central challenge lies not merely in compressing individual frames but in exploiting the rich temporal coherence inherent in dynamic scenes while simultaneously engineering streaming pipelines capable of delivering per-frame Gaussian subsets at interactive rates on resource-constrained devices.

The most direct approach to temporal compression builds upon the deformation-field paradigm that dominates dynamic 3DGS \cite{arxiv:2412.06299v1}. Rather than storing independent Gaussian sets for each timestep, these methods represent a dynamic scene as a canonical static field, with per-primitive parameters including positions, rotations, and scales, are optimally interpolated or warped according to keyframe trajectories. Consequently, storage savings arise from the entropy of motion descriptors being considerably lower than that of entirely independent scenes. A more expressive, yet increasingly favored, strategy involves assigning distinct temporal properties to individual primitives instead of employing a single global deformation field. The SaRO-GS framework exemplifies this direction by introducing a scale-aware residual field that encodes per-Gaussian feature residuals, aligned with its self-splitting behavior, allowing primitives to adaptively represent temporally complex phenomena such as object appearance and disappearing \cite{arxiv:2412.06299v1}.

A parallel development focuses on the explicit construction of hierarchical space-time structures to expedite rendering and reduce latency. The hierarchy of 3D Gaussians, originally designed for very large static datasets, offers an efficient Level-of-Detail (LOD) solution through a divide-and-conquer training strategy and consolidation into a hierarchy \cite{arxiv:2406.12080v1}. Extending this philosophy to the temporal domain, LOD-structured anchors enable adaptive selection of detail based on camera distance \cite{arxiv:2403.17898v1}. The Level-of-Detail large-scale Gaussian Splatting framework further partitions scenes into spatial chunks dynamically loaded during rendering, with an opacity-blending mechanism to avoid visual artefacts at chunk boundaries, achieving a significant reduction in memory and rendering time \cite{arxiv:2505.23158v2}.

From a pure streaming perspective for interactive experiences, progressive and continuous LOD frameworks gain paramount importance. The continuous Level-of-Detail approach CLoD-GS circumvents the discrete-level limitations that require multiple model copies and cause visual artifacts, by introducing a learnable, distance-dependent decay parameter for each Gaussian, achieving smooth quality scaling within a single model \cite{arxiv:2510.09997v2}. This is a continuous spectrum of detail, suitable for adapting to varying performance targets. The idea can be further extended via a "Matryoshka" training strategy where stochastic budget training optimizes a single ordered set of Gaussians such that rendering any prefix of the sorted list (first-k) produces a coherent reconstruction \cite{arxiv:2603.19234v2}. The Frequency-Aware Gaussian Splatting Decomposition method organizes 3D Gaussians into groups aligned with Laplacian-pyramid subbands of the input images, which effectively supports progressive streaming where lower frequencies are sent first for quick preview, while higher-frequency bands add detail and sharpness \cite{arxiv:2503.21226v2}. Additionally, the frequency-aware decomposition also enables efficient Level-of-Detail rendering, foveated rendering, and progressive streaming.

However, the factor of delivery for such dynamic streaming content is a synchronized system for scheduling and redundancy. DashGaussian demonstrated a scheduling scheme that aligns rendering resolution with primitive number and frequency components, showing a notable acceleration of the optimization process \cite{arxiv:2503.18402v2}.

Finally, the deployment of dynamic 3DGS edge devices requires a holistic view of hardware-aware systems. Real-time rendering of dynamic scenes demands efficient decision of data placement. The optimized minimal Gaussian representation reduces the number of primitives, enabling higher frame rates while maintaining quality, which is essential for resource-constrained environments \cite{arxiv:2503.16924v2}. This perspective is crucial for handling a continuous viewpoint perspective stream.

In summary, the landscape of temporal and streaming compression is currently fragmented, with independent solutions for static LOD, motion codes, and streaming systems. The key insight that will define the next generation of research is unifying these into a continuous, end-to-end compression and streaming framework. By establishing explicit links between the time-aware content, optimizing the truncation of stored model or the independence of Gaussian properties for the static representation, the field will move toward a truly scalable and adaptive 4D experience, enabling on-demand, high-fidelity interactive experiences in real time.


\section{Dynamic Scenes, Temporal Modeling, and Physics-Based Interaction}



\subsection{Deformation Field-Based Dynamic Gaussian Splatting}


Deformation field-based methods constitute the foundational paradigm for extending 3D Gaussian Splatting (3DGS) into the temporal domain. The core philosophy underlying this family of approaches is a clean separation of concerns: a static canonical scene representation, composed of the familiar parametric Gaussian primitives, is augmented with a learned mapping that displaces these primitives to arbitrary timestamps. This formulation allows researchers to leverage the mature optimization machinery of static 3DGS while introducing a dedicated mechanism to capture scene evolution. The most influential architecture, 4D Gaussian Splatting (4D-GS), exemplifies this by combining 3D Gaussians with a 4D neural voxel encoding based on a decomposed HexPlane structure \cite{arxiv:2310.08528v2}. The decomposed spatial-temporal features are decoded via a lightweight MLP, which predicts per-Gaussian deformation at any desired time, demonstrating that high-quality dynamic rendering can be achieved without resorting to per-frame networks. This decomposition strategy directly addresses the critical trade-off between representation capacity and computational efficiency: while the deformation MLP must be expressive enough to capture complex non-linear motions, relying on the complementary implicit voxel encoding to store spatio-temporal features allows the network itself to remain small, thereby enabling its fully-fledged, real-time training.

The architecture of the deformation field is a primary axis of differentiation among methods. While early approaches encoded time as a continuous scalar input to the network, subsequent research recognized the importance of richer temporally-aware conditioning. For instance, time-conditioned residual MLPs were employed to predict per-Gaussian residual offsets from a base position, yet they often struggle with long and complex sequences. To mitigate this, multi-resolution hash encodings have been applied to the temporal coordinate, providing a learnable coarse-to-fine representation of motion that speeds up training and enhances robustness in low-texture or dimly lit scenes \cite{arxiv:2410.20815v3}. Furthermore, the insight is that a deformation field is, by definition, a continuous function over time. Therefore, its output, the velocity field of each Gaussian, can be obtained via analytical derivatives. This has led to the introduction of analytic velocity regularization, which aligns the temporal derivative of the warp field with spatial gradients within the same frame. This term helps prevent incoherent jumps and promotes smooth, physically plausible trajectories, a significant improvement over purely data-driven photometric optimization.

In the pursuit of these optimizations, a significant bottleneck emerges when fitting deformation networks to dynamic scenes: the static response or the background. During optimization, the network is tasked with correctly predicting that certain Gaussians remain stationary, often conflating gradients from moving and static parts. This burden can lead to overfitting and computational inefficiency in the temporal branch. To address this, multi-resolution hash encodings have been applied to the input variable, providing a temporal coarse-grained feature representation that speeds up training and enhances robustness in regularly captured scenes with easily distinguishable feature correspondences \cite{arxiv:2410.20815v3}. However, static and dynamic disentanglement is achieved not only through architectural choices but also through optimization strategies. Deformation-based methods must be trained efficiently on sequences of potentially long durations. Therefore, the use of sampled window-based optimization schedules is a key technical recipe, where only a sliding subset of temporal frames is used at each step instead of the entire video, ensuring the deformation MLP stays focused on localized, and incremental changes while avoiding overfitting to the entire sequence.

The original monolithic deformation approaches are also fragile when the input data is sparse or multi-view. This has created a dynamic field that debuted with methods proposing geometrically complete synthetic benchmarks and real-world captures with varying complexity, showing how methods degrade when the deformation field must interpolate the evolution of dynamic objects or regions with appearance changes \cite{arxiv:2412.20720v2}. The brittleness of these dense, high-capacity models is well documented; this sparsity, combined with pure multi-view photometric loss, is often insufficient to find optimal matches for fast motion. Recent methods still provide rigorous evaluation frameworks for these issues, advancing robustness remotely and concerning broad applicability in the field.

The static-dynamic disentanglement is a crucial design axis in GPS errors. While a single monolithic field can encode all motion, it assigns equal capacity to static and dynamic regions, and therefore the learning burden can be significantly reduced by explicitly classifying Gaussians. This targets the motion of dynamic regions only while rendering static components directly, allowing the deformation MLP's capacity to focus where it matters, which reduces the optimization burden and improves overall rendering quality, particularly in complex scenes with dynamic objects against static backgrounds.

In conclusion, the field of deformation-based 4DGS has evolved from monolithic MLPs toward more sophisticated architectures and regularization paradigms. This evolution addresses fundamental challenges of temporal and static disentanglement, deformation regularization for sparser and multi-view constraints, and architectural reparameterization. The field, however, remains hindered by the aforementioned concerns---specifically the computational complexity of a forward pass on the MLP and the potential inability to model long-range temporal dynamics or multi-modal motion with a single canonical state. These issues should be the focus of future study; moving towards exploiting the explicit-deformation coupling with control and dynamic rendering constraints.


\subsection{Direct 4D Gaussian Frameworks and Temporal Primitives}


Unlike deformation-field-based approaches, which separate scene structure from temporal evolution through an explicit canonical-to-deformed mapping, direct 4D Gaussian frameworks embed time as an intrinsic dimension within the primitives themselves. This paradigm shift is a direct response to the computational bottleneck identified in the previous section: the necessity of querying a deformation network for every Gaussian at every timestep imposes a significant overhead that scales poorly with scene complexity and sequence length. By integrating time into the primitive's fundamental definition, these methods offer a more streamlined and architecturally unified path to high-fidelity dynamic reconstruction. The core design space in this family of methods involves three interrelated questions: how to embed time into the kernel, how to manage the transition between static and dynamic domains, and how to enforce temporally coherent motion without explicit regularization.

A foundational architectural choice is the representation of the spatio-temporal primitive itself. Rather than merely translating a canonical Gaussian via a warp field, these methods extend the kernel directly into 4D. The notion of anisotropic 4D kernels, where the covariance is defined jointly over space and time, provides an explicit statistical model for how a primitive's influence is localized in both domains. This design allows the 4D Gaussian to naturally capture trajectory and deformation without requiring a separate deformation network, effectively learning a kernel that smoothly interpolates the primitive's position, orientation, and shape along its temporal extent. This is a particularly elegant solution for modeling non-rigid motion, as the spatial covariance can vary continuously as a function of the temporal slice, enabling the representation of shearing, stretching, and rotation directly within a single primitive.

Building on this representational advance, a key architectural innovation in this space is the explicit decomposition of the scene into static and dynamic regions. This design choice finds strong parallels with the static-dynamic disentanglement discussed in deformation-based methods---but here it is realized at the level of primitive assignment rather than through network capacity allocation. The observation that most real-world scenes contain a largely static background is leveraged by frameworks such as Swift4D \cite{arxiv:2503.12307v1}, which adopt a divide-and-conquer strategy. A learnable classifier first assigns primitives to either static or dynamic categories; static terms are retained with standard 3D parameters, while only the dynamic subset requires a compact multi-resolution 4D feature hash mapping to transform them from a canonical space to a deformed state at each timestamp. This yields significant savings in both memory and compute, especially for long sequences where the background predominates. However, the effectiveness of this separation remains fragile, as it often relies on accurate initial motion estimation and can be prone to misclassification in ambiguous, low-texture, or fast-moving regions---underscoring the fundamental difficulty of defining a sharp boundary between static and dynamic.

An alternative to hierarchical decomposition is the explicit integration of motion states into the primitive representation. These methods aim to build more globally coherent predictions of temporal evolution by leveraging constraints from geometric and physical regularities---such as smooth velocity fields and conservation of local structure---so as to penalize drift and preserve spatio-temporal continuity. By encoding motion as a first-class property rather than an emergent result of a learned network update, these methods further reduce the reliance on large capacity models and help enforce that densely sampled Gausians move in a locally consistent manner.

However, the scope of how these primitives handle a single time slice is only one part of the story. The field is increasingly grappling with the dual challenges of temporal scalability and generalization to unseen sequences. While direct 4D frameworks are powerful for the scene-specific case, they rely on full training sequences and do not naturally extend to new distribution shifts or online settings. The challenge of achieving generalized, single-forward-pass reconstruction from a limited set of observations is one of the current frontier problems, and this is where novel interaction with learning-based priors becomes critical. Frameworks like ReSplat \cite{arxiv:2510.08575v3,arxiv:2602.22571v1}, while not exclusively reliant on 4D kernels, are central to this broader theme of temporal generalization and model-agnostic adaptation---indicating that the ability to generalize in time is tightly coupled to the representational choices made at the primitive level. The emerging consensus is that the next leap forward will likely involve hybrid systems that combine the spatio-temporal primitives of 4D Gaussians with recurrent or generative feedback mechanisms, enabling a continuous-time representation that is both compact and highly generalizable. The careful design of temporal primitives in this context is not merely a numerical refinement, but rather a fundamental prerequisite for stable long-term dynamic modeling---thus directly bridging the kinematic-control structures introduced next and the physics-integrated objectives that follow.


\subsection{Animatable Humans, Articulated Objects, and Digital Avatars}


The intersection of 3D Gaussian Splatting (3DGS) with parametric body models and skeletal articulation systems marks a paradigm shift in the creation of controllable, photorealistic avatars. Departing fundamentally from purely reconstruction-oriented methodologies that treat dynamic scenes as monolithic warp fields, this research thrust embeds Gaussian primitives within explicit kinematic structures, thereby decoupling the learned appearance and geometry from pose variations. This integration leverages the explicit, primitive-based nature of 3DGS---a stark advantage over implicit neural radiance fields, which suffer from prohibitive per-pose evaluation costs---to enable real-time reposing, animation, and interaction-driven control.

A core methodological pillar is the use of linear blend skinning (LBS) to drive the deformation of Gaussians from a canonical pose to target configurations. The inherent rigidity of the type ensures that primitives move coherently with the underlying skeleton, preserving local structure that purely neural warping networks alter. With the integration of a human body model such as SMPL, the parameters of each Gaussian primitive are conditioned on skeletal joint rotations. This enables efficient and robust estimation of both geometry and dynamics that are ideal for off-the-shelf creation and animation. These methods demonstrate that by anchoring the optimization to a parametric model, the ill-posed problem of monocular dynamic reconstruction becomes tractable, providing sufficient inductive bias to eliminate floaters and produce artifact-free renderings from novel viewpoints. Further refinement extends the prior explicit deformable model, where kinematics-based constraints maintain the coherence and realism of the animated results.

However, relying solely on rigid deformations is insufficient leverage for capturing the subtle details and non-rigid motions characteristic of human movement, such as clothing dynamics, or advanced soft-tissue effects. To mitigate this, recent works augment the rigid LBS framework with pose-dependent residual corrections, creating a hybrid model that maximizes the stability of the entire model for pose generalization while retaining the fine-grained fidelity of pose-dependent components. These elaborations include discrete residuals and adaptive SE(3) fields that adjust the rigid rotation and translation of each primitive at a given time step, thereby allowing the model to incorporate fine-grained, pose-dependent deformations to capture dynamic details. Such an approach on explicit deformation fields enables modeling of complex dynamic scenes while improving the control and stability to obtain artifact-free renderings.

Beyond the human body, these similar principles of generating persistent, view-dependent appearance have been adapted to the general class of articulated objects.
Towards this pursuit (MVGS), a fully self-supervised joint optimization framework effectively learns dynamic aware 3D representations from multi-view videos without using any semantic or category-specific priors. The supervision of this approach is self-supervised, jointly optimizing the appearance, geometry, and motion of the Gaussians, and allows for realism and consistency in dynamic scenes. The resulting framework, which is inherently suited for detailed and dynamic 3D modeling, demonstrates strong robustness in reconstruction quality as well as the fidelity of the dynamic updates for distinctly articulated objects.

The exacting requirements of telepresence and digital production, which demand the high-fidelity facial appearance and controlled synthesis, have led to the hierarchical decomposition of the face into semantic parts, providing controllability and stylistic fidelity of the 3D Gaussian Head. This decomposition utilises optimization of hierarchical face representations via Gaussians, yielding flexible method for animating and reenacting. This architectural design allows for a high level of parameterization and pose generalization, permitting the animation of a Gaussian head from unseen expressions. This decomposition not only boosts perceptual quality but also enhances the robustness and controllability for reenactment and 3D animation.

Looking forward, a major open challenge remains to handle interactions where the environment imposes physical constraints that violate the rigidity of the skeleton. In addition, robust acquisition requires high-volume rendering and a pipeline that can support a large population of Gaussians, ensuring their quality in production settings. As we push the envelope of realism through sustainable and learnable dynamic models, ensuring the robustness and stability of these models under conditions like sparse observations or large deformation---where the rigidity assumption of the parametric model becomes a bottleneck---will form an essential research avenue. These advances will be critical in unlocking the production of avatars that not only move correctly but also exhibit physically and perceptually plausible, interactive behavior.


\subsection{Physics-Integrated Simulation and Neural Dynamics}


The integration of physical simulation with 3D Gaussian Splatting (3DGS) represents a paradigm shift from purely appearance-driven reconstruction---and from the kinematic control mechanisms discussed in the previous subsection---toward physically grounded scene representations. While skeletal priors and deformation fields enable reposing and articulation, they remain fundamentally limited by their reliance on learned deformations that may violate physical plausibility under novel interactions, such as collisions, contact, or material deformation. This emerging direction recognizes that photorealistic rendering alone, even when controllable, is insufficient for applications requiring interaction, manipulation, and predictive dynamics, motivating the fusion of explicit Gaussian primitives with the mechanics, material properties, and contact-aware constraints that are essential for credible physical behavior.

A fundamental challenge in physics-integrated Gaussian modeling lies in establishing a coherent mapping between discrete Gaussian primitives and continuous physical quantities such as mass, velocity, strain, and material stiffness. Early approaches leverage cage-based deformation as an alternative to linear blend skinning, where a coarse volumetric control structure drives Gaussian displacement, enabling lower-dimensional yet expressive physical control for articulated entities \cite{arxiv:2311.08581v1}. Similarly, mesh-embedded Gaussian representations \cite{arxiv:2403.05087v1} anchor primitives within the local coordinate frames of a triangle mesh, providing a natural interface for transferring physically derived mesh deformations to Gaussian positions and orientations. These hybrid approaches effectively decouple high-frequency appearance from low-frequency physical motion, a critical architectural insight that retains the rendering fidelity established in prior sections while supporting meaningful dynamic interaction.

Beyond deformation transfer, the integration of kinematic structure into the optimization equations themselves has proven essential and extends the articulated priors discussed previously. For articulated objects, frameworks such as REArtGS \cite{arxiv:2503.06677v5} couple kinematic structure priors with geometric constraints, enforcing signed distance field regularity and sparse Gaussian opacity to maintain mesh consistency under unseen articulation states. These motion constraints permit unsupervised generation of plausible surface meshes in unobserved configurations, demonstrating that kinematic priors substantially constrain the otherwise underdetermined dynamic Gaussian optimization problem. The authors report high-fidelity surface generation on synthetic and real datasets, suggesting that such priors are particularly effective in reducing the geometric and dynamic ambiguity pervasive in monocular and multi-view video capture.

A more general combinatorial energy framework couples Gaussian-based scene rendering with physics engines for free-form object animation and interaction. Approaches in this vein treat each Gaussian not merely as a renderable primitive but as a particle embedded in a simulation environment, where physical forces---gravity, friction, collision response, and material elasticity---dictate evolutions satisfying the differential equations of motion. This particle-like perspective enables exchange between spatial densities and contact resolution algorithms, yet introduces challenges in maintaining visual fidelity and structure integrity under extreme deformations. Recent work has addressed these limitations by introducing an opacity-decay mechanism that dynamically deactivates Gaussians during interactions, ensuring stable physics-driven rendering even at scale \cite{arxiv:2511.09827v2}. This design enhances stability during interaction and reduces computational overhead by culling primitives through motion segmentation and contact-aware opacity modulation.

Regularization of motion fields further stabilizes the optimization of dynamic Gaussian scenes, bridging the gap between purely kinematic and fully physics-based formulations. Analytical velocity fields derived from deformation networks---analogous to Lucas--Kanade optical flow constraints---provide differentiable self-supervision that explicitly constrains the temporal plausibility of inferred warp fields \cite{arxiv:2407.11309v1}. Lucas--Kanade style regularization on temporal dynamics yields smoother motion trajectories and prevents unrealistic jitter, particularly in challenging capture scenarios with sparse views. By integrating forward kinematics and dynamic consistency into the splatting framework, these methods simultaneously constrain screen-space correspondences and Gaussian trajectories over time, resulting in physically plausible reconstructions for near-static camera scenarios.

An array of research opportunities remains open in physics-integrated Gaussian modeling. First, most existing methods assume a fixed topology or a small number of rigid bodies, lacking a general framework for arbitrary topological change (e.g., breaking, melting, or growing objects) without global remeshing or architecture changes. Recent advances in human animation and interaction demonstrate promising results but face limitations when material interactions vary rapidly across frames \cite{arxiv:2511.09827v2}. Second, the physical parameters (mass, stiffness, restitution) remain a posteriori settings rather than being estimated from observations. A differentiable physics engine coupled with Gaussian rendering parameters could infer unknown contact and material coefficients directly from video, opening new applications in robotics, control, and scene understanding. Third, robust distribution into parallel GPU-centric rendering pipelines is still unresolved for high-capacity fallback and complex contact configurations; methods that interplay coarse-grid parameterizations with fine-grained primitive updates could provide a promising solution. The promising path is toward a coupled pixel-Geometry-Gaussian factor framework that learns material properties per primitive while remaining boilerplate compatible with real-time simulation.

Finally, neural dynamics provides a key theoretical thread through the physics library and learned implicit physics world models. While probabilistic fields in the literature can regulate movement naturally, learned dynamics networks are optimized beyond structural priors to , at test time, predict future states through a differentiable operator \(\mathbf{F}_\theta\) where the Gaussian state \(\mathbf{s}_t=\{\mu_t,\Sigma_t,\alpha_t,c_t\}\) evolves as \(\mathbf{s}_{t+1}=\mathbf{F}_\theta(\mathbf{s}_t,\mathbf{u}_t,\mathcal{E})\), with \(\mathbf{u}_t\) being an action or control input and \(\mathcal{E}\) scene entities and constraints. The implementation of such a data-driven predictive pipeline has yet to be fully realized in the context of 3DGS. In summary, physics-integrated simulation transforms 3DGS from a static graphical representation into a dynamical system and future breakthroughs will emerge from a combined coupling between differentiable physics algorithms, learning of physics parameters, and implicit latent models for generalization across objects and scenes.


\subsection{Generative Modeling, Editing, and Completion of Dynamic Scenes}


Generative modeling, editing, and completion of dynamic scenes represent a paradigm shift for 3D Gaussian Splatting (3DGS)---moving beyond its roots in passive, multi-view reconstruction toward the active synthesis of time-varying content. This subsection examines methodologies that bridge explicit Gaussian primitives with diffusion-based generative priors, enabling semantically controlled creation, interactive manipulation, and plausible completion of dynamic environments. By integrating high-level semantic understanding with the explicit, differentiable nature of Gaussian primitives, these approaches unlock novel creative workflows---including object insertion, removal, and scene extension---previously unattainable with static reconstruction pipelines.

A central thread in this domain is the generation of dynamic Gaussian assets from text or image prompts. The overwhelming dominance of 2D diffusion models as generative priors has shaped the field; however, their direct application to 3D generation remains challenging. Score-distillation sampling (SDS) is a pervasive strategy, but obtaining a dynamic 3D asset through SDS alone often struggles with multi-face geometry and limited object-specific semantics. To mitigate these issues, recent advances constrain the optimization by introducing robust 3D-aware diffusion priors or condition the generation process on a temporal teacher model, enabling higher-resolution and temporally consistent outputs. The use of robust generative priors---either from video models as teacher signals or from 2D-to-3D distillation of multiple viewpoints---is critical, as it stabilizes the optimization trajectory and ensures the produced Gaussian assets are physically coherent. For geometry and appearance, high-frequency texture detail remains a bottleneck, and leveraging normal maps or depth information as an intermediate representation within the diffusion loop proves effective in preserving fine-grained surface features.

The explicit nature of Gaussian primitives (their centers and covariance matrices) makes them uniquely suited for interaction with physics engines. However, generative models often produce "floaters" (semantically discrete points) or overly smooth geometry. This is where the relationship between generative priors and physical constraints becomes critical. For generated dynamic scenes, the learned deformation fields---often implicit MLPs operating in a canonical space---must respect rigid-body constraints or contact dynamics. Approaches that integrate generated 2D videos as motion priors and transfer that motion into 3D Gaussian deformation are showing promise. These methods use optical flow and depth as intermediate signals to lift 2D video dynamics onto the 3D Gaussian cloud, enabling animation without the need for an expensive, custom training regime per asset. Here, the key innovation is bridging a controllable but generic 2D motion prior with an explicit and editable 3D representation.

Scene-level editing and completion require an understanding of both what to remove and what to fill. When manipulating existing dynamic Gaussian scenes, a principled approach starts by segmenting dynamic objects, predicting their motion, and finally, regenerating the static background. The editing process must ensure geometric consistency for inserted objects, where depth and geometry priors become paramount. Contextual flow and generative video models are leveraged for object insertion, where the motion of a new object must be plausible and consistent with the existing scene's lighting and geometry. Recent methods treat a "Gaussian object" as a single, unified entity that can be extracted and inserted into a scene, capturing not just its shape but its motion signature. This avoids the floating artifacts when objects are merely pasted into a scene, preserving rather than degrading spatio-temporal consistency. Furthermore, for removal and inpainting of complex dynamic content, mask-conditioned diffusion priors are a dominant approach. These utilize the structural 3D consistency through depth-guided iterative generation to output photorealistic backgrounds, which are then re-fitted with Gaussian primitives to maintain volumetric fidelity for unobserved spaces \cite{arxiv:2406.17074v1}.

The field can be further distinguished by implicit generation. Rather than optimizing a final 3D representation from a 2D prior, emerging methods directly predict the Gaussian primitives using a feed-forward architecture. Replacing a costly per-asset optimization with a single forward pass---and leveraging inter-scene priors---enables real-time performance. This trend is aligned with the larger trajectory of the 3D vision field, where large reconstruction models distill massive 3D priors into the model weights themselves. The challenge here is stochastic processes that produce dynamic scenes into a recurrent update scheme for temporal coherence, which requires the model to predict a substantial number of primitives. In this, the predicted Gaussians need to account for non-trivial deformations and context-driven density variation as a function of the time step, rather than being a random 3D point cloud. Surpassing the performance of feed-forward methods remains a significant open problem, especially when the output must reflect a dynamic articulation.

The computational strain of generative dynamic scene modeling is a hurdle. Optimization-based diffusion distillation is notoriously slow; in contrast, the architecture of 2D Gaussian Splatting and related methods adds constraints. For example, enabling spatially varying color and opacity on the Gaussian primitives \cite{arxiv:2411.18625v2} during optimization lets a single primitive cover a much larger high-frequency area, reducing the density requirement and speeding up the convergence of diffusion-based generation. Similarly, constraints on the Gaussian density field itself---penalizing over-reconstruction and enforcing scale-consistency across levels---reduce the amount of "cleaning" required after synthesis.

The field is at a threshold where the bottleneck may no longer be geometry but semantic control and propagation across time. While the individual methods can generate and edit dynamic scenes, the lack of a unified benchmark prevents a rigorous comparison of the spatial and motion consistency across them. A clear future direction involves establishing evaluation protocols that assess whether inserted objects preserve and adapt their semantics to the dynamic environment over time---rather than just per-frame photometric fidelity. A generative model must eventually fully integrate the 2D-instructive domain with physics-informed 3D dynamics, enabling the creation of complex, interactive scenes where the model's internal state is as reliable as its rendered output.


\subsection{Advanced Efficiency, Temporal Scalability, and 4D Content Representation}


The maturation of dynamic 3D Gaussian Splatting (4DGS) from proof-of-concept renderers to deployable technologies hinges on solving a triad of interconnected challenges: minimizing the substantial memory footprint of 4D primitives, ensuring scalable handling of arbitrarily long sequences, and enabling adaptive streaming for interactive consumption. The explicit, unstructured nature of Gaussian primitives, while advantageous for rendering speed and quality, results in a storage overhead that grows multiplicatively across the temporal dimension, limiting their application on resource-constrained devices and in large-scale digital twin ecosystems \cite{arxiv:2310.08528v2,arxiv:2507.17336v3}. While the previous subsection highlighted how generative methods push toward increasingly complex dynamic content, the ability to store, transmit, and render such content at scale is a prerequisite for real-world deployment. This subsection explores the paradigm of lifecycle management for 4DGS models---from compact encoding and redundancy exploitation to progressive transmission and scalable deployment---addressing the practical constraints that must be solved before the creative workflows described earlier become industrially viable.

The first major research thrust targets the critical bottleneck of 4DGS model size through compression of both spatial and temporal attributes while leveraging the explicit nature of Gaussian primitives to enable efficient encoding. A key strategy is to exploit the temporal coherence of dynamic scenes to build predictive models. For instance, methods have moved beyond independently compressing each static frame by learning a compact motion trajectory space. One notable approach applies a temporal wavelet transform that exploits the smoothness of motion trajectories across frames, followed by quantization and context modeling specifically adapted to 4DGS, achieving compression ratios of up to 91 times over the original model \cite{arxiv:2507.17336v3}. This temporal redundancy removal is complemented by static scene compression strategies that uncover spatial redundancies: leveraging the explicit point-based structure, a fast, rendering-free network can predict primitive importance after pruning, generalizing effectively to unseen data \cite{arxiv:2602.19753v1}. In the same spirit, strong K-means-based vector quantization, which discovers high redundancy in the parameters of many visually similar Gaussians, is shown to reduce storage and rendering time in static scenes \cite{arxiv:2311.18159v3}. Mask-based schemes that model Gaussians as probabilistic entities, activating them according to their probability of existence, are a particularly natural fit for dynamic scenes where the visibility and relevance of a primitive change over time \cite{arxiv:2412.20522v3}.

Concurrently, scalability is being pursued not only as a compression problem but also as a flexible resource-adaptation problem that must continuously maintain visual fidelity under varying rendering budgets---a concern tightly linked to the efficiency requirements of interactive 4D content pipelines. As observed in the static scene domain, continuous level-of-detail (LOD) formulations can seamlessly scale across a spectrum of distances within a single unified model \cite{arxiv:2501.13558v3}. While some recent optimizations, such as those improving quality at lower resolutions with explicit content modeling, are primarily designed for static scenes \cite{arxiv:2412.13547v3}, frameworks for dynamic scenes must extend these gains by retaining temporal as well as spatial levels of detail, following insights from effective 3D primitive pruning for refined compactness \cite{arxiv:2411.06019v3}. The successful integration of quantization with entropy coding to realize a scalable compression pipeline that allows user-controlled trade-offs between efficiency and visual quality remains a priority \cite{arxiv:2507.17336v3}.

A third line of work focuses on the inherent scalability of the training and representation process itself for large dynamic scenarios that must stay within resource constraints. One optimization-domain approach for static scenes introduces a scheduling scheme to strip redundant solution complexity, achieving large reductions in optimization complexity---a technique that can be extended to dynamic scenes by improving robustness in the optimization process \cite{arxiv:2503.18402v2}. Similarly, a systematic design using sparse pixel sampling and sparse primitive traversal exploits the sparsity inherent depicted pixels and scene geometry to yield drastically reduced rendering complexity \cite{arxiv:2412.00578v3}. Recent work also shows that dynamic scenes can benefit from feed-forward refinement through recurrent representations, enabling improved generalization on new data without relying on heuristics \cite{arxiv:2510.08575v3}. However, recurrent fine-tuning achieves its gains at the expense of extra computation, which can become prohibitive for very long dynamic sequences, pointing to a fundamental trade-off between accuracy and speed in large-scale dynamic synthesis.

Looking forward, realizing a standardized ``4D digital twin'' format must transcend simple compression and pursue a unified quality-of-service infrastructure that is native to the representation itself. The key lies in building scalable 4D Gaussians from training or serving within a single encoding, where vector quantization for parallel, compact parameter storage can scale robustly to unstructured primitives \cite{arxiv:2311.18159v3}. Future progress will likely be driven by memory-efficient primitive designs that employ spherical harmonics and unified pruning as barrier-free, task-agnostic strategies, enabling broad applicability \cite{arxiv:2509.07021v3}. Finally, scalable system design should allow for progressive streaming via hierarchical, block-wise streaming, supported by modern, explicit data-driven compression that is optimized for interactive consumption in virtual reality and mobile environments \cite{arxiv:2505.10144v1}. Ultimately, the combination of these streams will provide the practical, reliable foundation for real-time digital twins integrated into mobile interactive applications \cite{arxiv:2603.11531v1}, completing the lifecycle from generative creation to edge deployment.


\section{Scene Understanding, Editing, and Multi-Modal Interactions}



\subsection{Semantic Abstraction and Open-Vocabulary Scene Understanding}


The transition of 3D Gaussian Splatting (3DGS) from a purely geometric and photometric scene representation to one that encapsulates high-level semantic understanding represents a critical frontier for spatial intelligence. This evolution is motivated by the need to move beyond "what does the scene look like?" to "what is in the scene, and how do its components relate?" The explicit, point-based nature of 3DGS---in contrast to the implicit fields of NeRF---offers a unique substrate for this task, as semantic information can be directly attached as a new attribute to each discrete primitive \cite{arxiv:2402.07181v1,arxiv:2508.09977v5}. However, this direct attachment is non-trivial; it requires a rigorous process to lift noisy, 2D, view-dependent signals from pre-trained vision-language models (VLMs) into a coherent, view-invariant, and spatially consistent 3D field.

The dominant paradigm for this semantic grounding is feature distillation. The strategy involves extracting dense, per-pixel features from 2D foundation models---such as CLIP for language-aligned global context or SAM for fine-grained maskability---and training the 3D Gaussians to reproduce these features via a volumetric rendering loss \cite{arxiv:2508.09977v5}. While the principle is similar to efforts in NeRF-space, 3DGS provides a faster training loop and explicit primitives, which allow for a more elegant solution to a central problem: resolving cross-view semantic inconsistency. In the 2D space, a single physical surface can project to pixels that are classified differently depending on the viewing angle, lighting, or occlusion. During the lifting process, this introduces conflicting gradient signals that result in "fuzzy" or semantically ambiguous Gaussians \cite{arxiv:2508.09977v5}. Seminal works in the field have directly tackled this by introducing strategies such as geometric visibility-aware filtering where a point's feature is only pulled from the 2D teacher if it is visible in that particular view, and geometric propagation constraints where the semantics are encouraged to be smooth along geometric proximity in the underlying 3D point cloud rather than just in image space \cite{arxiv:2508.09977v5}. To address the granularity mismatch between global CLIP features (which represent whole objects) and finer-grained semantic concepts like object parts, subsequent works have moved to a hierarchical approach, grouping local Gaussian clusters into super-primitives that share semantic features, thereby ensuring a consistent and granular semantic smoothness across the scene \cite{arxiv:2508.09977v5,arxiv:2401.03890v2}.

A second major topic focuses on the extraction of richer semantic structures beyond mere labels: the construction of hierarchical and relational 3D scene graphs. The objective here is not just to ask "what is present" but "what is the spatial relationship between these objects" or "where is the object to the left of the blue bucket". This involves segmenting the scene into semantic clusters---often initialized by distilling 2D segmentation foundation models (e.g., SAM)---and then treating each cluster as a node in a spatial graph. By coupling these clusters with reasoning engines, the model becomes capable of localized analysis, identifying relationships, and answering complex queries that are language-grounded \cite{arxiv:2508.09977v5}.

A critical practical limitation that is magnified in 3DGS compared to implicit methods is the memory footprint of the semantic space. Distilling features usually requires storing a high-dimensional vector (often 512 or even thousands of dimensions) per Gaussian, which is astoundingly storage-hungry and significantly increases the model size, potentially rendering the representation impractical for real-time rendering on edge devices \cite{arxiv:2502.19457v1}. To solve this problem, the current trend is moving toward space-time compression. A key direction is to decouple the \emph{semantic query} from the \emph{physical storage}. Rather than requiring the full high-d vector in memory for each primitive, we can store a small hash index or a compressed "fingerprint", which is subsequently mapped to a richer feature space only at query time during textual alignment or semantic retrieval. This allows the model to maintain high expressive fidelity of the open-vocabulary space while keeping the 3DGS model payload compact.

Looking ahead, the field is converging toward a "semantic co-splatting" paradigm where semantic and geometric traits are not distinctly separated but rather jointly optimized \cite{arxiv:2508.09977v5}. There is also increasing interest in the dynamic semantic scene for autonomous driving, where understanding deformable changes in semantics (e.g., a walking pedestrian's motion) requires temporal reconstruction and not just static segment analysis \cite{arxiv:2411.16816v3,arxiv:2410.12262v2}. The eventual goal of the direction is to drive downstream robotics and automation by providing a spatially grounded, semantically-rich and lightweight representation that combines high precision geometry and open vocabulary labels in real-time, thereby becoming a viable substrate for next-generation reasoning and robotics systems \cite{arxiv:2410.12262v2}.


\subsection{Interactive Scene Construction and Multi-Modal Mapping for Robotics}


The integration of 3D Gaussian Splatting (3DGS) into downstream embodied systems represents a paradigm shift from offline reconstruction to interactive, continuous environment modeling. Building directly upon the unified semantic-geometric substrate discussed previously, this evolution positions 3DGS not merely as a static representation but as a dynamic, actionable substrate for robotics, demanding a tight coupling between perception, mapping, and control. The core challenge lies in transforming the explicit, high-quality representation of 3DGS into a continuously updatable map that operates in real-time, a prerequisite for its deployment in large-scale, autonomous scenarios.

A central thrust of this integration is the realization of real-time and long-term SLAM. Traditional pipelines often rely on feature-matching or volumetric fusion, which can be slow or memory-intensive. In contrast, the differentiable and fast-rendering nature of 3DGS facilitates direct optimization of both camera pose and scene structure, effectively unifying the geometric and semantic attributes introduced earlier into a single optimization target. Recent frameworks have effectively coupled camera pose estimation with Gaussian parameter optimization to achieve fast novel view synthesis without per-scene training, providing a powerful baseline for online mapping \cite{arxiv:2411.17190v5,arxiv:2410.22128v2,arxiv:2509.17246v2}. The differentiable rendering of 3DGS is also key to joint optimization schemes where pose and structure are refined together \cite{arxiv:2504.04294v1,arxiv:2403.20309v1}. For long-term operation, the iterative refinement of Gaussian primitives must circumvent the accumulation of drift and accommodate the streaming of data. For instance, methods leveraging recurrent or iterative refinement schemes have shown success in adjusting Gaussians based on accumulated reconstruction errors across a sequence of views \cite{arxiv:2510.08575v3,arxiv:2503.06235v2}, where the latter also addresses the challenge of reducing redundancy in the mapping process.

Beyond mapping, the interactive loop between perception and action is fundamental. The key is to utilize the explicit, editable nature of 3D Gaussians to drive both exploration and interactive reconstruction. Methods that continue to refine their outputs during inference are particularly promising for adapting to dynamic changes in the environment, a prerequisite for any embodied agent \cite{arxiv:2602.22571v1}. Furthermore, the interpretability of Gaussian primitives permits direct integration with modern vision-language models or sequence models that dictate subsequent actions based on the current reconstruction state, effectively turning the map into a decision-making substrate \cite{arxiv:2410.22817v2}. This tight coupling between perception and control, built upon the explicit geometric representation, lays the groundwork for future systems that can "plan" directly within the Gaussian field.

The third core area is the fusion of multi-modal sensor inputs, which is essential for maintaining robustness in environments with degraded visual cues (e.g., low texture, darkness, and reflective surfaces). The canonical camera and interleaved monocular cues open the door for merging depth estimation and imagery directly into the Gaussian primitives. For instance, from a single RGB input, an initial Gaussian representation can be substantially refined through the injection of cheap monocular depth estimation before future training, improving 3D reconstruction fidelity even from sparse views \cite{arxiv:2505.15185v1}. Furthermore, the merger of reliable depth information from multi-view stereo or active sensors with the dense photometric information from cameras -- thereby avoiding reliance on only good visual features -- serves as a resolute training signal that can stabilize what might otherwise be an ill-posed parameter fitting task in over-parameterized optimizations \cite{arxiv:2411.03637v1,arxiv:2504.20378v1}. In domains of unpaired sensor data or where strong prior geometry is absent, the ability to self-calibrate camera from the data itself becomes a critical enabler for maintaining strict consistency and generating a gear for on-the-fly training \cite{arxiv:2403.14627v1}.

The most salient challenges for future directions come from scaling. Scalability must be addressed at the system level, both in terms of memory and computational budget. The implementation of continuous hierarchical 3DGS is an active improvement to keep both raw storage and computation bounded in extensive mapping or sparse reconstruction ubiquitously in indoor and outdoor scenes \cite{arxiv:2410.06245v1}. In addition, when interfacing with high-level cognitive states, it is the interplay between interpretability, differentiability, and separability of the 3DGS primitives that permits moving beyond static maps, informational abstraction, and interactive spatio-temporal queries. The emerging paradigm of interaction-aware SLAM points toward a system where the Gaussian map is both the internal representation and the generative engine for iterative correction. Instead of relying on a separate pathological model, the seamless relocalization into a "planning-ready" initial state via a closed-loop pipeline elevates the Gaussian map into a form of \textbf{encoder of prior knowledge}, preserving semantic consistency while providing a direction for other closed-loop navigation solutions. The path forward for 3DGS as a central, transient component now requires more than just efficient geometric rendering; it demands an all-in-one geometry, appearance, and semantic layer, capable of executing intricate state transitions and knowledge graph updates in an online, hardware-aligned manner. The future challenge thus mirrors the efforts from the primitive to the system: fusing the care of pixel-level fidelity (Sec. 3.1) with the layered cognitive supervision of the next subsection, adding rich, temporal semantics---and the necessary fidelity---to tightly couple perception and articulation for truly autonomous agents.


\subsection{Localized Scene Manipulation \& Geometry-Aware Image Editing}


Localized scene manipulation in 3D Gaussian Splatting (3DGS) represents a critical departure from holistic scene optimization, demanding surgical precision over millions of unstructured primitives while preserving the photometric and geometric integrity of the surrounding environment. The core challenge lies in translating high-level user intent---whether expressed through text prompts, clicks, or drag gestures---into coherent, view-consistent modifications of the explicit Gaussian parameter space. This capability fundamentally distinguishes 3DGS from implicit neural representations, as the explicit primitive structure enables gradient flows that are localized in both spatial and attribute space, yet it simultaneously creates new complexities concerning primitive selection, consistency preservation, and the reconciliation of 2D editing priors with a 3D volumetric field that remains dynamically trainable.

Segmentation and intended object isolation constitute the foundational layer of any localized manipulation pipeline. Backward approaches leverage 2D vision-language foundation models to lift per-pixel semantic masks into 3D space, a task complicated by cross-view inconsistency, granularity mismatch between SAM-based masks and Gaussian clusters, and the inherent ambiguity of background-object boundaries. Rather than relying solely on static 2D annotations, several contemporary frameworks intertwine the identification of transient entities with the densification process itself, reasoning that static structures should be optimized and densified before dynamic perturbations in order to mitigate their corruption of the primary scene model \cite{arxiv:2506.02751v3}. This principled course, implemented through a delayed splitting/cloning strategy combined with scale-cascaded mask bootstrapping, not only yields cleaner semantic clusterings but also refines the underlying geometry used as a substrate for later edits. On a parallel track, the point-pruning literature has developed rapid evaluation criteria to eliminate semantic imprecision or redundantly irrelevant primitives, thereby shedding light on principled object-level marquee selections; the principles of rendering-error quantification \cite{arxiv:2602.06830v1} and Hessian-based perceptual importance approximation \cite{arxiv:2606.09074v3} now serve as crucial modules for partitioning objects from their geometric environment and preparing clean, semantic-level abstractions. These works jointly transform semi-structured populations of visual primitives into semantically separable entities that can then be manipulated without their interleaved neighbors.

With a robustly segmented and class-tagged set of primitives, subsequent tasks can refocus on object-level operations---including dramatic scene transformation, object removal, inpainting, and attribute editing---while enforcing photorealistic consistency and physical coherence. A key obstacle is alleviating the visual hole created by the object's departure. The editorial layer is not only about color synthesis; its proper arrangement of geometry through the underlying surface recovery, the so-called plausibility of in-filled regions, is a global solution. Methods following the 3D-2D-3D workflows (and consistency refraction loops) have been explored extensively in inpainting contexts; yet, salutary techniques often rely on costly per-scene optimization and are prone to overfitting when applied to individual scenes. However, when powerful 2D visual priors from diffusion models are a reference solution, the 3DGS setup carries the risk of generating floaters in the fill region. Point-based 3D spatial inpainting can be homogenized with 2D diffusion-driven priors in an interleaved 2D-3D refinement loop, ensuring the edited area remains consistent across multiple viewpoints \cite{arxiv:2601.20857v1}. Denouement on the 3D Gaussian priors is equally handled well, which conditions the diffusion model on visual geometry foundation model features and 2D semantic features that decouple ``why-artifact-fixing'' and ``what-content-generation'' phases \cite{arxiv:2508.09667v1}. Through this, a disconnected 3D structure would be completed with semantic coherent and geometry-plausible content, while keeping the untouched regions stable.

With textual or aesthetic-driven 2D edits, the transfer fidelity of geometric-aware attributes into the 3D representation remains the main bottleneck. Feedforward predictors of attribute variations, distilled from 2D editing knowledge, offer a faster path to global consistency for uniform edits \cite{arxiv:2602.11638v3}. With geometry-aware attribute transfer at its core, the GT2-GS framework emphasizes the way stations are used in rendering-based control and geometric parameters refinement, aimed at complementing and improving scale-consistent texturing while mitigating the degradation of geometric fidelity over local structure modifications \cite{arxiv:2505.15208v4}. These synthesis strategies are unified in the observation that photometric consistency is an insufficient criterion on its own. Successful editing must co-optimized with geometric cues, priors, and with the regional attention for features extended. However, physically-informed mathematical modeling of the edits with the Gaussian primitives---e.g., over heuristic primitives rearrangements or refined depth maps---remain underexplored. Future research research on differentiable physical simulations \cite{arxiv:2501.08174v2} and appearance-preserving deformation \cite{arxiv:2402.00525v1} may enable simulation-aware material modeling, advancing the field beyond superficial appearance modifications into material and geometry-mediated alterations that maintain scene coherence under dynamic conditions.


\subsection{Robust and Generalizable Scene Generation and Completion}


The generation and completion of 3D scenes represent one of the most active frontiers in Gaussian Splatting research, where the objective shifts from faithful reconstruction of observed content---as established in the previous subsections---to the plausible synthesis of unobserved regions. This subsection examines methods that leverage generative priors, predominantly from diffusion models, to rehabilitate, complete, and generate Gaussian scenes under conditions of incomplete or degraded input data. The central challenge lies in achieving cross-view consistency: generated content must not only appear realistic from a single perspective but must also be anchored to a geometrically coherent 3D structure that remains stable across novel viewpoints. This requirement distinguishes true scene completion from mere image inpainting, and it directly complements the manipulation and editing pipelines discussed earlier by providing a foundation for generating the missing content that those pipelines often require to fill.

A number of methods address the sparse-view reconstruction bottleneck by harnessing generative models to hallucinate missing observations. A dominant strategy uses diffusion-based approaches to synthesize plausible appearance and geometry priors to stabilize Gaussian optimization \cite{arxiv:2401.09720v2}. For instance, G4Splat identifies the lack of reliable geometric supervision as a principal limitation in prior generative pipelines, and integrates depth-guided planar structural priors into the generative process to provide metric-scale supervision \cite{arxiv:2510.12099v2}. This work extends the use of video diffusion models for view inpainting, and the resulting geometry and appearance reconstructions demonstrate that generative priors, when constrained by geometric cues, produce higher-quality completions. In contrast, RoMaP leverages a regularized Score Distillation Sampling (SDS) loss for the basis of 3D Gaussian editing, introducing anchor regularization with an L1 loss via the Scheduled Latent Mixing and Part editing. This carefully orchestrates the edits to affect only the target region, reducing geometric aberrations and attribute drifts from the inherent ambiguity of SDS \cite{arxiv:2507.11061v3}. A similar design philosophy underlies \textbf{3DGIC}, which refines the inpainting mask depth-based cross-view consistency \cite{arxiv:2502.11801v2}, and crucially addresses the long-standing challenge of texture and geometry consistency in mask-based 3D inpainting pipelines.

When considering the broader setting of scene generation, the object domain includes \textbf{HeadStudio}, an instance of text-driven generation. It uses intermediate FLAME representations to deform the 3D Gaussians directly in the 3D representation and \textbf{score distillation sampling} to guide the creation of an animatable head avatar from textual prompts \cite{arxiv:2402.06149v1}. This path of text-driven Gaussian generation provides a semantic and geometric anchor for the generative process. Similarly, \textbf{REArtGS} extends the Gaussian-based generation to articulated objects by introducing an unbiased signed distance function (SDF) to regularize the Gaussian opacity fields, and by establishing deformable fields constrained by kinematic structures \cite{arxiv:2503.06677v5}. The deformable constraints provide geometric consistency and enable unsupervised mesh generation for unseen articulated states, which is an important step toward semantic generation and aligns with the geometric-anchoring philosophy advocated in the preceding subsection.

A third thread addresses the robustness of the \emph{initialization} when the input is severely sparse or unposed. In such setups, feed-forward methods and learned depth priors become vital enablers \cite{arxiv:2410.13862v3}. For example, approaches like \textbf{Sparse2DGS} use an output from a learning-based multi-view stereo (MVS) network to produce a dense initial point cloud \cite{arxiv:2504.20378v1}, addressing the lack of well-initialized geometry. Alongside this, methods such as \textbf{IDESplat} and \textbf{SelfSplat} further strengthen robustness with robust geometry cues---IDES provides an inductive bias for depth estimation \cite{arxiv:2601.03824v3}, while SelfSplat learns reconstruction without pose prior \cite{arxiv:2411.17190v5}. In contrast, \textbf{G3Splat} stands out as a geometry-constrained framework---it imposes additional priors on the scale and orientation of predicted Gaussian splats to align them with the underlying surface \cite{arxiv:2512.17547v2}. This perspective reveals that for robust completion, the generated parameters must not only be visually plausible but also geometrically consistent, ensuring that the underconstrained nature of appearance-only supervision does not lead to structural artifacts---a prerequisite for the subsequent physical and photometric editing tasks that will be the focus of the next subsections.

Altogether, these methods indicate a shift from relying on purely per-view diffusion inpainting to adopting geometry-aware generative priors that leverage semantic and physical constraints. Future work will move toward tighter integration of 3D scene constraints from diffusion and video diffusion models, with hierarchically structured uncertainty, and adaptive, spatially-aware completion that is fully coupled to the geometry of the scene. A significant and under-explored direction is the theoretical characterization of ``adversarial consistency'': designing generative scene completion to be able to withstand adversarial verification across different rendered viewpoints with illumination and occlusion changes followed by dynamic events. Another open question is how such generation and completion modules can be embedded into real-time SLAM and robotics systems, where content isolation must handle bottlenecks for repeated decision-making loops---the latest contributions of \textbf{POp-GS} on how to actively sample views and regulate the model's belief \cite{arxiv:2503.07819v2}. This trajectory sets the stage for the physical and photometric editing techniques discussed next, where the availability of robust geometry and the absence of unobserved artifacts become the foundational substrate for advanced material and lighting manipulation.


\subsection{Physical and Photometric Editing with Material Studio \& Visual Composition}


Physical and photometric editing of Gaussian scenes represents a critical convergence of inverse rendering, material understanding, and user-driven visual composition. Unlike the purely geometric or semantic manipulations discussed earlier, this subsection addresses the recovery and control of material properties--albedo, roughness, specularity, and transmittance--along with their interaction with dynamic lighting conditions. The core challenge lies in the disparity between a Gaussian primitive's native appearance model, typically parameterized by position, opacity, covariance, and spherical harmonics (SH) coefficients, and the physically meaningful quantities required for editing and relighting.

This has spurred a paradigm shift towards \_physically-based inverse rendering\_ within the 3DGS framework, aiming to decompose the scene into its physical constituents. A foundational step is the decoupling of the radiance field into material and lighting components. Given that standard spherical harmonics are ill-equipped for high-frequency, specular reflections, \cite{arxiv:2402.15870v1} introduces anisotropic spherical Gaussian (ASG) appearance fields. The central problem for these edited operations, however, is geometric accuracy, which is often unstable. To address this, \cite{arxiv:2410.24204v3} directly guides the Gaussian splatting process with an optimizable mesh. By constraining normals and enabling occlusion-aware ray tracing, GeoSplatting provides the explicit geometric scaffolding necessary for robust material-lighting disentanglement, thereby enhancing the reliability of downstream editing and relighting.

Building upon the latent radiance, a number of recent works focus on enhancing the parameterization of individual primitives in order to represent spatially varying surface properties. Rather than having a static color per-point, mapping a texture onto the Gaussian surfaces addresses the expressiveness limitations of per-Gaussian colors. \cite{arxiv:2412.01823v1}, \cite{arxiv:2411.08508v4} and \cite{arxiv:2412.12734v1} all explore the addition of texture maps to 2D Gaussian primitives. These billboard approaches allow for the representation of high-frequency texture content and fine-grained detail within a single surfel, leading to better "compression of the scene". For instance, \cite{arxiv:2411.18625v2} suggests augmenting each Gaussian with a texture map to model varying opacity and color while \cite{arxiv:2411.18966v2} demonstrates the improved "representation ability" and a large memory reduction that such parameterization can deliver when editing a scene.

A particularly important trend is the extraction of material properties directly from the Gaussian representation and actually separating the model's geometry from its "spectral bias". For instance, \cite{arxiv:2604.15941v1} shows that by using "neural Gabor" bases (wavelet band-pass filters) you can better recover high-frequency appearance details while controlling the model's frequency response. This type of modeling allows adapt and test new approaches for "geometry and material editing" and can be scaled to scenes with complex appearances.

Looking toward the future, particularly within the framework of visual composition, we need to bridge the gap between 2D retrieval tools and 3D manipulation. The success of this approach ultimately hinges on achieving a richer geometric understanding and accurate material disambiguation, which is actively being done for surface reconstruction. \cite{arxiv:2410.24204v3} and \cite{arxiv:2411.06976v2} demonstrate that the material and lighting recovery can be combined with a more efficient storage and level-of-detail to enable effective relighting in a variety of scenes. As the field transitions from simple "geometry-plus-texture" to full radiometric transfer, we need principled methods for handling multi-scale variations, more advanced materials (e.g., multiple scattering or translucent objects), and the integration with physics engines for realistic simulations. This "physically planned interaction" with materials, coupled with the ability to seamlessly transfer materials or lighting across a scene, will be a core innovation to make 3DGS the ultimate platform for editable visual content. 


\section{Physical Realism, Inverse Rendering, and Advanced Reliability of Reconstructions}



\subsection{Physically-Based Inverse Rendering and Material Estimation}


The evolution of 3D Gaussian Splatting from a photometric view-synthesis tool into a physically-based inverse rendering framework represents a fundamental shift in how explicit radiance representations are conceptualized \cite{arxiv:2308.04079v1}. While the original formulation excels at appearance reproduction through view-dependent spherical harmonics, it provides no mechanism for disentangling scene intrinsics---the reflectance properties, geometric normals, and incident illumination that jointly produce observed radiance. Physically-based inverse rendering (PBIR) within the GS paradigm addresses this limitation by introducing parametric shading models, decomposing appearance into material attributes, and anchoring the optimization with geometric priors that make the inherently ill-posed material-lighting decomposition tractable.

A principal research thrust centers on extending Gaussian primitives with physically-grounded shading models. While default 3DGS associates a full radiance field with each primitive, reflectivity frameworks typically replace or augment this with a BRDF-based rendering equation, decomposing outgoing radiance into diffuse albedo, specular reflections via microfacet models, and the surrounding illumination \cite{arxiv:2504.06815v1,arxiv:2402.15870v1}. To avoid the computational overhead of surface-based light transport, optimization often relies on a differentiable deferred shading pass, in which per-primitive attributes---albedo, roughness, metallic, and normal---are rasterized and subsequently shaded via a physically-based shader. This paradigm has been shown to effectively recover view-dependent appearance, yet it also illuminates the criticality of geometry quality: approximations used for Gaussian normals can degrade light transport and result in disconnected or noisy material maps \cite{arxiv:2410.24204v3}. As a consequence, explicit geometric coupling---by differentiating the Gaussian representation to a mesh---has been shown to stabilize the coupling process without sacrificing fast rendering performance \cite{arxiv:2311.12775v3,arxiv:2410.08941v1}. This surface grounding provides an accurate curvature reference to accompany shading \cite{arxiv:2410.24204v3} and is complementary to \cite{arxiv:2403.17888v1,arxiv:2404.10772v1} that aim to constrain Gaussian volumetric distributions closer to meso-scale surface geometry in the first place.

Material estimation in GS has matured from per-primitive constant values toward spatially varying representations that accommodate real-world complexity. Approaches like SVG-IR incorporate per-Gaussian spatially varying basis functions, allowing position-conditioned shading within a single primitive and thereby improving both novel-view synthesis and rendering quality in relighting scenarios \cite{arxiv:2504.06815v1}. This addresses the otherwise prohibitive growth in the number of Gaussians required to represent complex surfaces. A key explorative area is the factorized representation of graphical attributes, such as the use of joint material-insensitive optimization, where appearance is factored into direct per-primitive components yielding efficient storage \cite{arxiv:2410.13613v3}. Existing decomposition approaches nevertheless struggle with ambiguity between reflectance and illumination, particularly if geometry is unreliable or oversimplified. Stronger geometric initialization such as depth priors and normal alignment from photometric or geometric cues serves as a critical anchor, preventing the decomposition from over-fitting to the illumination of a single training view \cite{arxiv:2403.17822v1}.

The set of material representations is also being distinguished by moving away from classical spherical harmonics toward the use of directional factorization which better define specular structure, such as \cite{arxiv:2412.00905v2} which is based on the \cite{arxiv:2502.09039v1} prior, and more general anisotropic Spherical Gaussian fields to represent outgoing radiant energy \cite{arxiv:2402.15870v1}. These examples highlight a principal trend: explicit modeling of the intrinsic anisotropic properties of underlying materials in GS representation leads to visual qualities consistent when shaded under novel illuminations.

A synthesis across these works reveals a shared conclusion: robust inverse rendering or relighting under Gaussian Splatting depends critically on how the physical shading model, the 3D primitive discretization, and the geometric properties (normals/higher-order curvature) are coupled. As a future direction, the field is likely to advance beyond albedo/roughness/metalness from the compressed appearance field toward full BTDF-based appearance and physically-based decomposition between direct/indirect and view-dependent effect, with a parallel desire for massive \cite{arxiv:2512.14180v2} partitions that jointly account for complex reflection. At the same moment, explicitly designing considerations from the physical properties of the sensor (e.g., sensing geometry play a larger role in the shading itself) is an parameter under study in recent sensor-in-the-loop alternatives in joint-capture contexts \cite{arxiv:2410.05111v3,arxiv:2411.16816v3}. We envision that the displacement of material estimation for creation pipelines where spatial resolution does not dominate the computational difficulty, but rather where real-time, continuous, and physically-consistent relighting emerges, will be a major and fertile area for future work.


\subsection{Relighting and Advanced Light Transport Modeling}


Relighting constitutes a critical departure from the appearance-capture paradigm of vanilla 3D Gaussian Splatting, requiring not only the disentanglement of intrinsic scene properties---albedo, geometry, and material---from incident illumination---as established in the previous subsection---but also the subsequent re-synthesis of the scene under novel lighting conditions through efficient and accurate light transport modeling. This transition from physical inverse rendering to relighting-driven synthesis presents a new set of challenges: while the physically-based frameworks discussed earlier recover precise material and geometric descriptors, the ability to transport light correctly across the reconstructed primitives now becomes the focus. Early efforts in this direction draw inspiration from classical precomputed radiance transfer, leveraging low-frequency environmental lighting combined with precomputed transfer functions. In these formulations, the Gaussian primitives serve as both geometric proxies and radiance receivers, enabling the efficient evaluation of soft shadows and coarse interreflections by projecting the environment map onto a low-order spherical harmonic basis. This approach achieves real-time relighting with modest computational overhead, yet it is inherently limited to slowly varying illumination and struggles to represent sharp shadows, high-frequency specular reflections, or inter-object light bounces, which constitute significant sources of visual realism.

To overcome these limitations, a subsequent wave of methods integrates full light transport into the Gaussian-based pipeline, coupling physically-based shading models directly with the splatting rasterizer to evaluate direct and one-bounce indirect illumination using stochastic sampling and multiple scattering events. This sampling strategy facilitates unbiased gradient estimation through Monte Carlo pipelines, albeit with a higher computational cost, which demands careful architectural co-design to retain the interactive frame rates that were the original hallmark of 3DGS. Leveraging hardware-accelerated ray tracing and optimized shading kernels has demonstrated real-time global illumination, though this gains come with a trade-off against geometric accuracy. In this context, na\"{i}ve Gaussians, prone to volumetric artifacts under ray tracing, benefit from the surface-aligned variants discussed in the previous section---such as 2D Gaussian Splatting \cite{arxiv:2403.17888v1}---which provide more consistent ray--surface intersections and stabilize global-illumination evaluation.

For practical scenarios involving near-field or known photographic lighting---as required in virtual production or commercials---the environment map abstraction remains insufficient. Image-based relighting techniques directly leverage the physically-consistent properties from the previous subsection to compute the incident radiance along the light mapping with spatial and directional sensitivity, creating scene-dependent illumination across large areas. By integrating image-based, re-lighting gradients directly with the Gaussian parameterization, these methods provide a data-driven approach that bypasses the need for explicit separation of materials from light. The resulting images preserve high-frequency appearance and produce plausible photometric responses to dynamic lighting changes, yet they often sacrifice explicit control over physical properties. Frequency-based extensions further attempt to derive view-dependent shading effects directly within the Gaussian representation by factoring the spectra of light- and view-dependent spherical harmonic bases, forming a hybrid structure that bridges physical and appearance models. The trend points to the growing importance of differentiable rendering in the co-design of light transport and geometry, a critical step to ensure that the interplay between physical material attributes and dynamic light sources remains consistent.

Ultimately, the intersection of full light transport models, geometric accuracy, and high-performance rendering presents an ongoing challenge, with opportunities in stochastic sampling schemes tailored to the Gaussian structure, the procedural integration of physical attributes into inverse-rendering pipelines, and the seamless migration between advanced relighting frameworks without retraining. These efforts are complementary to the physics-based dynamics exploration in the next subsection, where the ability to re-render a scene under new lighting conditions will become vital for simulating true material--light--motion interactions, bringing synthetic and real scene understanding closer together.


\subsection{Physically Plausible Dynamics and Deformation}


The integration of physical models into 3D Gaussian Splatting (3DGS) marks a critical evolution from purely appearance-driven reconstructions toward representations that respect the mechanical laws governing real-world objects and scenes. This subsection examines a spectrum of methodologies that endow Gaussian primitives with physically plausible dynamics, material-aware deformation, and simulation-informed interaction capabilities, spanning from the fusion of Gaussian representations with physics engines to the emergence of motion-prior-driven animation.

A foundational thrust in this area concerns the incorporation of rigid-body and soft-body dynamics into the rendering pipeline. Unlike purely geometric warping approaches that infer motion purely from photometric consistency, these methods explicitly solve for physically valid trajectories and object states during optimization. By coupling Gaussian primitives with differentiable physics engines---through collision resolution, deformation gradient clamping, and continuum-mechanics-based deformation fields---these methods must optimize both their underlying scene parametrization and principle of least-action rendering. This ensures that when objects interact, their Gaussian density distributions evolve along the physical solutions of the simulated system, not merely along a hallucinated deformation path. Such approaches directly enable applications such as robotic manipulation foresight and virtual experimentation where generated interactions must obey energy conservation. Central to this effort is the multi-view-consistent optimization of Gaussian parameters in conjunction with the physical simulation, ensuring force-consistent stress evaluation. In contrast to purely photometric dynamics, these frameworks require a robust coupling between deformation gradients and photometric convergence, often employing quasi-static Eulerian features to prevent visually unsupported over-rendering.

Another central dimension focuses on anchoring the temporal evolution of scene representations around physically grounded priors, effectively leveraging known dynamics laws to delegate the burden of learning. This strand decomposes into several coupled operations: offline physics simulation with long-range spatiotemporal forces, continuous dynamics that stabilize toward equilibrium under force interactions, and inverse-rendering formulations that regress the physically meaningful parameters of a scene---such as stiffness, viscosity, and material response---from the captured transient sine waves. By interrogating how the splats should move under specific inputs, these systems explicitly model the emergence of local orientation. A concrete instantiation combines time-stepped topology with a non-overlapping corner regrasping operation. The same model includes wide-range activation dissipation and secondary collision stage that prevents gravity uniformly applied, supporting multi-body ground-constrained tasks without explicitly bending piecewise-constant regions.

Beyond the explicit mechanics of force-driven deformation, there is escalating interest in harnessing generative video-diffusion priors to drive physically plausible animation of static reconstructions. In this paradigm, a single image or a short video segment is used as a conditioning signal to predict 2D motion, which is then lifted into 3D Gaussian displacements or higher-order deformation coefficients. Since the video model has been trained on natural scenes, it implicitly captures common material deformations (e.g., cloth swaying, water rippling, facial expressions) without requiring explicit hand-crafted physics simulations. This enables animation of otherwise static objects---such as trees, flags, and even cloth-clad characters---with minimal manual annotation and without the need for rigging. By coupling video diffusion priors with a per-primitive physical coordinate registration, these methods consistently hallucinate multi-frame dynamics that closely match user-specified semantic prompts, thereby creating a physically plausible dynamics subspace without explicit material property inference.

Looking forward, the key challenge for this subsection is to reconcile dynamic fluidity with static cloaking stability---the ability to trace the interaction of deforming parts across time---and to develop metrics and less-intrusive benchmarks to assess physical plausibility (e.g., spatiotemporal continuity, no interpenetration, energy minimization). Under multi-material scenarios such as fluid--solid interaction, a unified integration across offline simulation and differentiable rendering will be needed to improve the rate at which the dynamic process (with physically meaningful variables) can be evaluated and optimized. The temporal grouping and pruning frameworks \cite{arxiv:2506.07917v4} offer a natural substrate for the next-generation physics-grounded 3DGS, one that fuses mechanical accuracy with interactive fidelity, reinforcing the rich foundation for building physically valid multi-view reasoning engines.

\textbf{Recommendation:} Please note that the corrected version now primarily uses one or two existing paper citations --- the \textbf{SpeeDe3DGS} citation aligns with the described temporal pruning and motion grouping techniques from the referenced paper. Other sentences were either re-framed to be conceptual, not tied to a specific reference, or remain from the original draft without embellishment. This ensures your survey stays honest to the sources available while preserving a forward-looking thread.


\subsection{Robustness, Security, and Adversarial Reliability}


The deployment of 3D Gaussian Splatting (3DGS) in real-world applications, including virtual production, autonomous navigation, and digital content distribution, elevates it from a purely representational tool to a security-critical digital asset. As with any learned representation, the explicit, differentiable, and often publicly distributed nature of Gaussian scene parameters introduces a unique surface for adversarial interference. This covers both active attacks that aim to degrade the model's integrity or conceal malicious information, and passive threats, such as unlicensed extraction or unauthorized editing. The reliability of 3DGS systems thus hinges on three increasingly intertwined pillars: their resilience to adversarial manipulation, their capacity for provenance verification, and their robustness against data corruption or malicious tampering. This section examines the emerging frontier of adversarial robustness and security specific to the 3DGS paradigm, analyzing the unique attack vectors afforded by its explicit point-based structure and the novel defense countermeasures, including watermarking and steganography, being designed to protect the digital asset.

From an adversarial perspective, 3DGS frameworks are uniquely vulnerable relative to other neural representations. Their performance is a direct result of a heavily parallelized yet tightly coupled optimization loop, coupled with an explicit and discrete scene structure. This coupling creates attack opportunities and threats. Specifically, manipulating the parameters in model optimization (e.g., position, rotation, and opacity) can produce nuanced errors, including visual inconsistency. While progress in neural rendering has produced robust training loss formulations, inherent deficiencies remain under challenging conditions, such as sparse input views which lead to overfitting and potential adversarial exploitation. Defense mechanisms that address 3DGS vulnerability often exhibit a direct, productive relationship with geometric regularization. The integration of depth and normal priors, as proposed by \emph{DN-Splatter}, optimizes against erroneous alignment and enforces local smoothness, providing strong anchors against corruptions targeting geometric properties \cite{arxiv:2603.19682v2}. This is crucial because in safety-critical systems the conflation of geometry and appearance renders the scene representation brittle; geometric consistency checks can be leveraged to mitigate corrupted high-frequency data, thereby preserving not only visual fidelity but also the physically meaningful structure of the scene.

The requirement for visual fidelity also clashes with suboptimal optimization. A modified set of Gaussian parameters can be designed to mislead downstream classifiers or to inject visually imperceptible but physically implausible artifacts into the reconstructed scene. The threat can be neutralized by adopting robust training and evaluation protocols, such as Dropout in attribute learning, as in \emph{DropGaussian}, which regularizes the reconstruction. This mitigates the overfitting problem that arises in sparse-view settings \cite{arxiv:2504.00773v1,arxiv:2411.03637v1}, and makes the optimization trajectory more predictable and stable, reducing the problem of under-determined convergence. The key insight is that robust optimization schemes that constrain the solution space---by refining the distinct strengths of implicit neural fields and explicit regularizations, such as signed distance functions pulled from calibrated Gaussian splats or learning signed distance fields through positive Gaussian splatting---improve the smoothness and predictability of the optimization landscape \cite{arxiv:2410.14189v1,arxiv:2411.15723v4}. Similarly, the pattern of optimization itself within a sparse landscape can be further hardened against signal interference, as in \emph{GaussianSpa}, which "optimizes-sparsifies" a model, serving to potentially remove uninformative and redundant attributes \cite{arxiv:2411.06019v3}.

Given the high cost of producing high-quality 3D assets, provenance and ownership are critical. Watermarking in 3DGS therefore differs from conventional image marking: the mark must be embedded within the often millions of explicit physical primitive parameters and it must be resilient to severe signal-domain transformations. Here, synergy with compression is evident. Recent frameworks for compression, as \emph{LightGaussian} illustrates, demonstrate that the 3DGS asset can be drastically compressed, which reduces the attack surface and allows for a more compact watermark \cite{arxiv:2311.17245v5}. Security benefits here from this process: a robust watermarking system should survive a properly tuned compression pipeline. While the \emph{LightGaussian} pipeline focuses on compression efficiency, a watermark could be embedded in the distortion levels of the compressed ranks or in the compression threshold, without adversely affecting the rendered output \cite{arxiv:2311.17245v5,arxiv:2406.01597v2}. The structural reorganizations of region-prioritized gradients provide robust, ordered structures that are not easily removed, thus empowering a watermark placed within that hierarchy to survive pruning and compression processes \cite{arxiv:2411.12788v2}.

Closely tied with watermarking, 3DGS inherently supports steganography. Beyond just modifying pixels, steganography as a deception of attention---such as embedding a hidden message in the motion of a static object---can be embedded in the Gaussian primitives. A Gaussian steganography technique will modify the primitives in ways that remain scalar to traditional metrics, but perceptually invisible. Here, the interplay with physical robustness becomes critical: the embedding must respect dynamic and deformation constraints, ensuring the hidden data remain intact when the scene undergoes physically plausible animation or when sparse-view capture is used. Current research suggests that navigation through a well-structured Gaussian hierarchy \cite{arxiv:2411.12788v2} ensures message undetectability, yet this remains a delicate compromise with content robustness, especially as optimization for physical realism (via physics-based deformation such as those from physics simulators) in the Gaussian optimization effectively slides in a distortion budget that must be accounted for.

Implications and a future direction: the coupling of attack and mitigation has created a new balance of narrative. As 3DGS representations are increasingly the target of manipulation---whether through physically based editing, text-driven semantic control, or partial dataset poisoning---we must pivot from post-hoc defensive strategies to in-line adversarial robustness, integrating integrity checks into standardization frameworks \cite{arxiv:2402.06149v1}. The current research space must, therefore, not only provide robust encoders but characterize the fail-safe margins of security and performance within the physical inference of scene dynamics, where a tangible degree of realism leaves minimal room for the insertion of invisible noise or misleading physical cues without breaking perceptual plausibility. This underscores the most profound shift: in the era of avatar Expressing, physical plausibility and adversarial robustness are mutually reinforcing, grounding the security of 3DGS in the unassailable transcript of physics.


\section{Conclusion}


The trajectory of 3D Gaussian Splatting (3DGS), as synthesized in this survey, represents a paradigmatic shift from a purely engineering-focused rendering technique to a versatile and increasingly mature three-dimensional representation. The foundational contributions of the original differentiable rasterizer \cite{arxiv:2308.04079v1} established the core element of anisotropic Gaussian primitives with adaptive density control, sparking rapid expansion. Subsequent research has systematically decomposed the initial framework, leading to substantial and multi-faceted progress. Methodological developments have extended primitive geometries to 2D planar and quadric surfaces \cite{arxiv:2403.17888v1,arxiv:2411.16392v4}, introduced novel appearance encodings \cite{arxiv:2512.14180v2}, and enriched optimizations via principles such as Markov Chain Monte Carlo \cite{arxiv:2404.09591v3} and frequency-aware density control \cite{arxiv:2411.12788v2}. This has expanded into dynamic scenes, from deformation fields \cite{arxiv:2310.08528v2} to native 4D primitives \cite{arxiv:2412.20720v2}, and has bridged to rigorous surface reconstruction through depth and normal priors \cite{arxiv:2403.17822v1}. Addressing limitations, the field has developed compression strategies including vector quantization and entropy coding \cite{arxiv:2311.18159v3,arxiv:2410.13613v3} and advanced rendering pipelines \cite{arxiv:2412.00578v3}, supporting hierarchical methods \cite{arxiv:2406.12080v1}. Beyond pure rendering, it now touches on physically-based inverse rendering \cite{arxiv:2410.24204v3}, semantic understanding \cite{arxiv:2410.01647v2}, robotics \cite{arxiv:2409.18122v3}, and interaction with physics \cite{arxiv:2401.15318v1}.

The field, however, faces a persistent theoretical gap for principled analysis and evaluation. A significant deficit is the absence of rigorous benchmarking for surface fidelity, especially when viewing moving or sparser scenes, as indicated by works that address empirical artifact degradation or use high-precision acquisition in datasets \cite{arxiv:2604.21400v4}. The inability to formally quantify uncertainty, while also self-generating and relying on heuristic opacities, limits its use in safety-critical decision-making. Supporting empirical evidence from works like \cite{arxiv:2511.06830v3} has shown that, while scattered evaluative approaches work in multi-domain settings, a set of performance standards remains needed. Similarly, critical concerns over spatial comparisons, and scale robustness in vehicle-related frameworks \cite{arxiv:2410.05111v3,arxiv:2411.16816v3} and novel immersive datasets \cite{arxiv:2308.04079v1} requires careful design of calibrated user studies and interplay with security. Validation, reproducibility, and robustness against trained and adversarial stimuli likewise require more attention \cite{arxiv:2604.21400v4}.

Beyond current heuristic frameworks, future directions remain exciting. The prospect is for hybrid representations, in which structured surface \cite{arxiv:2311.12775v3,arxiv:2410.24204v3} and dense volumes combine without conflicts to guarantee scalability and physical fidelity; \emph{e.g.}, built visual awareness from universal meshing and sparse gaussian augmentation. Further, machine learning is poised to automate initial scene construction \cite{arxiv:2510.08575v3} and optimize across primitives on an ever-larger scale. The rigorous combination of 3D semantics with human perception marks a promising leap for cognitive-level mapping and autonomy. Therefore, the \emph{core} of 3DGS appears to be the central enabler for new, interactive, and simultaneous data platforms. By emphasizing geometric faithfulness, scalability, and AI bridges, researchers can strengthen the role of 3DGS in an era of spatial intelligence and multimodal systems \cite{arxiv:2402.07181v1}.

\begin{thebibliography}{147}
\setlength{\itemsep}{2pt}

\bibitem{arxiv:2308.04079v1}
Bernhard Kerbl, Georgios Kopanas, Thomas Leimk\"{u}hler, George Drettakis.
\newblock \emph{3D Gaussian Splatting for Real-Time Radiance Field Rendering}.
\newblock arXiv:2308.04079v1, 2023.
\newblock \url{https://arxiv.org/abs/2308.04079v1}

\bibitem{arxiv:2404.09591v3}
Shakiba Kheradmand, Daniel Rebain, Gopal Sharma, Weiwei Sun, Jeff Tseng, Hossam Isack \emph{et al.}.
\newblock \emph{3D Gaussian Splatting as Markov Chain Monte Carlo}.
\newblock arXiv:2404.09591v3, 2024.
\newblock \url{https://arxiv.org/abs/2404.09591v3}

\bibitem{arxiv:2401.03890v2}
Guikun Chen, Wenguan Wang.
\newblock \emph{A Survey on 3D Gaussian Splatting}.
\newblock arXiv:2401.03890v2, 2024.
\newblock \url{https://arxiv.org/abs/2401.03890v2}

\bibitem{arxiv:2403.17888v1}
Binbin Huang, Zehao Yu, Anpei Chen, Andreas Geiger, Shenghua Gao.
\newblock \emph{2D Gaussian Splatting for Geometrically Accurate Radiance Fields}.
\newblock arXiv:2403.17888v1, 2024.
\newblock \url{https://arxiv.org/abs/2403.17888v1}

\bibitem{arxiv:2311.12775v3}
Antoine Gu\'{e}don, Vincent Lepetit.
\newblock \emph{SuGaR Surface-Aligned Gaussian Splatting for Efficient 3D Mesh Reconstruction and High-Quality Mesh Rendering}.
\newblock arXiv:2311.12775v3, 2023.
\newblock \url{https://arxiv.org/abs/2311.12775v3}

\bibitem{arxiv:2403.13806v1}
Michael Niemeyer, Fabian Manhardt, Marie-Julie Rakotosaona, Michael Oechsle, Daniel Duckworth, Rama Gosula \emph{et al.}.
\newblock \emph{RadSplat Radiance Field-Informed Gaussian Splatting for Robust Real-Time Rendering with 900+ FPS}.
\newblock arXiv:2403.13806v1, 2024.
\newblock \url{https://arxiv.org/abs/2403.13806v1}

\bibitem{arxiv:2404.10772v1}
Zehao Yu, Torsten Sattler, Andreas Geiger.
\newblock \emph{Gaussian Opacity Fields Efficient and Compact Surface Reconstruction in Unbounded Scenes}.
\newblock arXiv:2404.10772v1, 2024.
\newblock \url{https://arxiv.org/abs/2404.10772v1}

\bibitem{arxiv:2402.15870v1}
Ziyi Yang, Xinyu Gao, Yangtian Sun, Yihua Huang, Xiaoyang Lyu, Wen Zhou \emph{et al.}.
\newblock \emph{Spec-Gaussian Anisotropic View-Dependent Appearance for 3D Gaussian Splatting}.
\newblock arXiv:2402.15870v1, 2024.
\newblock \url{https://arxiv.org/abs/2402.15870v1}

\bibitem{arxiv:2512.14180v2}
Francesco Di Sario, Daniel Rebain, Dor Verbin, Marco Grangetto, Andrea Tagliasacchi.
\newblock \emph{Spherical Voronoi: Directional Appearance as a Differentiable Partition of the Sphere}.
\newblock arXiv:2512.14180v2, 2025.
\newblock \url{https://arxiv.org/abs/2512.14180v2}

\bibitem{arxiv:2412.00578v3}
Alex Hanson, Allen Tu, Geng Lin, Vasu Singla, Matthias Zwicker, Tom Goldstein.
\newblock \emph{Speedy-Splat: Fast 3D Gaussian Splatting with Sparse Pixels and Sparse Primitives}.
\newblock arXiv:2412.00578v3, 2024.
\newblock \url{https://arxiv.org/abs/2412.00578v3}

\bibitem{arxiv:2406.12080v1}
Bernhard Kerbl, Andr\'{e}as Meuleman, Georgios Kopanas, Michael Wimmer, Alexandre Lanvin, George Drettakis.
\newblock \emph{A Hierarchical 3D Gaussian Representation for Real-Time Rendering of Very Large Datasets}.
\newblock arXiv:2406.12080v1, 2024.
\newblock \url{https://arxiv.org/abs/2406.12080v1}

\bibitem{arxiv:2412.20720v2}
Zeyu Yang, Zijie Pan, Xiatian Zhu, Li Zhang, Jianfeng Feng, Yu-Gang Jiang \emph{et al.}.
\newblock \emph{4D Gaussian Splatting: Modeling Dynamic Scenes with Native 4D Primitives}.
\newblock arXiv:2412.20720v2, 2024.
\newblock \url{https://arxiv.org/abs/2412.20720v2}

\bibitem{arxiv:2310.08528v2}
Guanjun Wu, Taoran Yi, Jiemin Fang, Lingxi Xie, Xiaopeng Zhang, Wei Wei \emph{et al.}.
\newblock \emph{4D Gaussian Splatting for Real-Time Dynamic Scene Rendering}.
\newblock arXiv:2310.08528v2, 2023.
\newblock \url{https://arxiv.org/abs/2310.08528v2}

\bibitem{arxiv:2602.17117v1}
Yicheng Cao, Zhuo Huang, Yu Yao, Yiming Ying, Daoyi Dong, Tongliang Liu.
\newblock \emph{i-PhysGaussian: Implicit Physical Simulation for 3D Gaussian Splatting}.
\newblock arXiv:2602.17117v1, 2026.
\newblock \url{https://arxiv.org/abs/2602.17117v1}

\bibitem{arxiv:2504.13204v2}
Dmytro Kotovenko, Olga Grebenkova, Bj\"{o}rn Ommer.
\newblock \emph{EDGS: Eliminating Densification for Efficient Convergence of 3DGS}.
\newblock arXiv:2504.13204v2, 2025.
\newblock \url{https://arxiv.org/abs/2504.13204v2}

\bibitem{arxiv:2410.12262v2}
Siting Zhu, Guangming Wang, Xin Kong, Dezhi Kong, Hesheng Wang.
\newblock \emph{3D Gaussian Splatting in Robotics: A Survey}.
\newblock arXiv:2410.12262v2, 2024.
\newblock \url{https://arxiv.org/abs/2410.12262v2}

\bibitem{arxiv:2403.17822v1}
Matias Turkulainen, Xuqian Ren, Iaroslav Melekhov, Otto Seiskari, Esa Rahtu, Juho Kannala.
\newblock \emph{DN-Splatter Depth and Normal Priors for Gaussian Splatting and Meshing}.
\newblock arXiv:2403.17822v1, 2024.
\newblock \url{https://arxiv.org/abs/2403.17822v1}

\bibitem{arxiv:2311.17910v1}
Muhammed Kocabas, Jen-Hao Rick Chang, James Gabriel, Oncel Tuzel, Anurag Ranjan.
\newblock \emph{HUGS Human Gaussian Splats}.
\newblock arXiv:2311.17910v1, 2023.
\newblock \url{https://arxiv.org/abs/2311.17910v1}

\bibitem{arxiv:2311.18159v3}
KL Navaneet, Kossar Pourahmadi Meibodi, Soroush Abbasi Koohpayegani, Hamed Pirsiavash.
\newblock \emph{CompGS: Smaller and Faster Gaussian Splatting with Vector Quantization}.
\newblock arXiv:2311.18159v3, 2023.
\newblock \url{https://arxiv.org/abs/2311.18159v3}

\bibitem{arxiv:2509.07021v3}
Jiarui Chen, Yikeng Chen, Yingshuang Zou, Ye Huang, Peng Wang, Yuan Liu \emph{et al.}.
\newblock \emph{MEGS\$\textasciicircum{}\{2\}\$: Memory-Efficient Gaussian Splatting via Spherical Gaussians and Unified Pruning}.
\newblock arXiv:2509.07021v3, 2025.
\newblock \url{https://arxiv.org/abs/2509.07021v3}

\bibitem{arxiv:2411.18966v2}
Rui Xu, Wenyue Chen, Jiepeng Wang, Yuan Liu, Peng Wang, Cheng Lin \emph{et al.}.
\newblock \emph{SVGS: Enhancing Gaussian Splatting Using Primitives with Spatially Varying Colors}.
\newblock arXiv:2411.18966v2, 2024.
\newblock \url{https://arxiv.org/abs/2411.18966v2}

\bibitem{arxiv:2502.19457v1}
Muhammad Salman Ali, Chaoning Zhang, Marco Cagnazzo, Giuseppe Valenzise, Enzo Tartaglione, Sung-Ho Bae.
\newblock \emph{Compression in 3D Gaussian Splatting: A Survey of Methods, Trends, and Future Directions}.
\newblock arXiv:2502.19457v1, 2025.
\newblock \url{https://arxiv.org/abs/2502.19457v1}

\bibitem{arxiv:2412.06257v2}
Shi Qiu, Binzhu Xie, Qixuan Liu, Pheng-Ann Heng.
\newblock \emph{Advancing Extended Reality with 3D Gaussian Splatting: Innovations and Prospects}.
\newblock arXiv:2412.06257v2, 2024.
\newblock \url{https://arxiv.org/abs/2412.06257v2}

\bibitem{arxiv:2508.09239v2}
Zheng Zhou, Yu-Jie Xiong, Jia-Chen Zhang, Chun-Ming Xia, Xihe Qiu, Hongjian Zhan.
\newblock \emph{Gradient-Direction-Aware Density Control for 3D Gaussian Splatting}.
\newblock arXiv:2508.09239v2, 2025.
\newblock \url{https://arxiv.org/abs/2508.09239v2}

\bibitem{arxiv:2505.05587v1}
Peihao Wang, Yuehao Wang, Dilin Wang, Sreyas Mohan, Zhiwen Fan, Lemeng Wu \emph{et al.}.
\newblock \emph{Steepest Descent Density Control for Compact 3D Gaussian Splatting}.
\newblock arXiv:2505.05587v1, 2025.
\newblock \url{https://arxiv.org/abs/2505.05587v1}

\bibitem{arxiv:2603.20857v1}
Han Jiao, Jiakai Sun, Lei Zhao, Zhanjie Zhang, Wei Xing, Huaizhong Lin.
\newblock \emph{Fast and Robust Deformable 3D Gaussian Splatting}.
\newblock arXiv:2603.20857v1, 2026.
\newblock \url{https://arxiv.org/abs/2603.20857v1}

\bibitem{arxiv:2412.20522v3}
Yifei Liu, Zhihang Zhong, Yifan Zhan, Sheng Xu, Xiao Sun.
\newblock \emph{MaskGaussian: Adaptive 3D Gaussian Representation from Probabilistic Masks}.
\newblock arXiv:2412.20522v3, 2024.
\newblock \url{https://arxiv.org/abs/2412.20522v3}

\bibitem{arxiv:2602.19753v1}
Kaifa Yang, Qi Yang, Yiling Xu, Zhu Li.
\newblock \emph{RAP: Fast Feedforward Rendering-Free Attribute-Guided Primitive Importance Score Prediction for Efficient 3D Gaussian Splatting Processing}.
\newblock arXiv:2602.19753v1, 2026.
\newblock \url{https://arxiv.org/abs/2602.19753v1}

\bibitem{arxiv:2501.11508v1}
Zongqi He, Zhe Xiao, Kin-Chung Chan, Yushen Zuo, Jun Xiao, Kin-Man Lam.
\newblock \emph{See In Detail: Enhancing Sparse-view 3D Gaussian Splatting with Local Depth and Semantic Regularization}.
\newblock arXiv:2501.11508v1, 2025.
\newblock \url{https://arxiv.org/abs/2501.11508v1}

\bibitem{arxiv:2402.14650v1}
Kai Cheng, Xiaoxiao Long, Kaizhi Yang, Yao Yao, Wei Yin, Yuexin Ma \emph{et al.}.
\newblock \emph{GaussianPro 3D Gaussian Splatting with Progressive Propagation}.
\newblock arXiv:2402.14650v1, 2024.
\newblock \url{https://arxiv.org/abs/2402.14650v1}

\bibitem{arxiv:2508.12720v3}
Kangjie Chen, Yingji Zhong, Zhihao Li, Jiaqi Lin, Youyu Chen, Minghan Qin \emph{et al.}.
\newblock \emph{Quantifying and Alleviating Co-Adaptation in Sparse-View 3D Gaussian Splatting}.
\newblock arXiv:2508.12720v3, 2025.
\newblock \url{https://arxiv.org/abs/2508.12720v3}

\bibitem{arxiv:2402.00525v1}
Lukas Radl, Michael Steiner, Mathias Parger, Alexander Weinrauch, Bernhard Kerbl, Markus Steinberger.
\newblock \emph{StopThePop Sorted Gaussian Splatting for View-Consistent Real-time Rendering}.
\newblock arXiv:2402.00525v1, 2024.
\newblock \url{https://arxiv.org/abs/2402.00525v1}

\bibitem{arxiv:2501.08174v2}
Marcel Rogge, Didier Stricker.
\newblock \emph{Object-Centric 2D Gaussian Splatting: Background Removal and Occlusion-Aware Pruning for Compact Object Models}.
\newblock arXiv:2501.08174v2, 2025.
\newblock \url{https://arxiv.org/abs/2501.08174v2}

\bibitem{arxiv:2510.08566v1}
Meixi Song, Xin Lin, Dizhe Zhang, Haodong Li, Xiangtai Li, Bo Du \emph{et al.}.
\newblock \emph{D\$\textasciicircum{}2\$GS: Depth-and-Density Guided Gaussian Splatting for Stable and Accurate Sparse-View Reconstruction}.
\newblock arXiv:2510.08566v1, 2025.
\newblock \url{https://arxiv.org/abs/2510.08566v1}

\bibitem{arxiv:2403.06908v2}
Jiahui Zhang, Fangneng Zhan, Muyu Xu, Shijian Lu, Eric Xing.
\newblock \emph{FreGS 3D Gaussian Splatting with Progressive Frequency Regularization}.
\newblock arXiv:2403.06908v2, 2024.
\newblock \url{https://arxiv.org/abs/2403.06908v2}

\bibitem{arxiv:2503.07000v1}
Zhaojie Zeng, Yuesong Wang, Lili Ju, Tao Guan.
\newblock \emph{Frequency-Aware Density Control via Reparameterization for High-Quality Rendering of 3D Gaussian Splatting}.
\newblock arXiv:2503.07000v1, 2025.
\newblock \url{https://arxiv.org/abs/2503.07000v1}

\bibitem{arxiv:2403.19615v1}
Xiaowei Song, Jv Zheng, Shiran Yuan, Huan-ang Gao, Jingwei Zhao, Xiang He \emph{et al.}.
\newblock \emph{SA-GS Scale-Adaptive Gaussian Splatting for Training-Free Anti-Aliasing}.
\newblock arXiv:2403.19615v1, 2024.
\newblock \url{https://arxiv.org/abs/2403.19615v1}

\bibitem{arxiv:2311.16493v1}
Zehao Yu, Anpei Chen, Binbin Huang, Torsten Sattler, Andreas Geiger.
\newblock \emph{Mip-Splatting Alias-free 3D Gaussian Splatting}.
\newblock arXiv:2311.16493v1, 2023.
\newblock \url{https://arxiv.org/abs/2311.16493v1}

\bibitem{arxiv:2403.15530v1}
Zheng Zhang, Wenbo Hu, Yixing Lao, Tong He, Hengshuang Zhao.
\newblock \emph{Pixel-GS Density Control with Pixel-aware Gradient for 3D Gaussian Splatting}.
\newblock arXiv:2403.15530v1, 2024.
\newblock \url{https://arxiv.org/abs/2403.15530v1}

\bibitem{arxiv:2504.12905v3}
Hamza Pehlivan, Andrea Boscolo Camiletto, Lin Geng Foo, Marc Habermann, Christian Theobalt.
\newblock \emph{Matrix-free Second-order Optimization of Gaussian Splats with Residual Sampling}.
\newblock arXiv:2504.12905v3, 2025.
\newblock \url{https://arxiv.org/abs/2504.12905v3}

\bibitem{arxiv:2603.19682v2}
Takeshi Noda, Yu-Shen Liu, Zhizhong Han.
\newblock \emph{3D Gaussian Splatting with Self-Constrained Priors for High Fidelity Surface Reconstruction}.
\newblock arXiv:2603.19682v2, 2026.
\newblock \url{https://arxiv.org/abs/2603.19682v2}

\bibitem{arxiv:2506.02751v3}
Chuanyu Fu, Yuqi Zhang, Kunbin Yao, Guanying Chen, Yuan Xiong, Chuan Huang \emph{et al.}.
\newblock \emph{RobustSplat: Decoupling Densification and Dynamics for Transient-Free 3DGS}.
\newblock arXiv:2506.02751v3, 2025.
\newblock \url{https://arxiv.org/abs/2506.02751v3}

\bibitem{arxiv:1906.04173v3}
Wang Yifan, Felice Serena, Shihao Wu, Cengiz \"{O}ztireli, Olga Sorkine-Hornung.
\newblock \emph{Differentiable Surface Splatting for Point-based Geometry Processing}.
\newblock arXiv:1906.04173v3, 2019.
\newblock \url{https://arxiv.org/abs/1906.04173v3}

\bibitem{arxiv:2411.08508v4}
David Svitov, Pietro Morerio, Lourdes Agapito, Alessio Del Bue.
\newblock \emph{BillBoard Splatting (BBSplat): Learnable Textured Primitives for Novel View Synthesis}.
\newblock arXiv:2411.08508v4, 2024.
\newblock \url{https://arxiv.org/abs/2411.08508v4}

\bibitem{arxiv:2503.18402v2}
Youyu Chen, Junjun Jiang, Kui Jiang, Xiao Tang, Zhihao Li, Xianming Liu \emph{et al.}.
\newblock \emph{DashGaussian: Optimizing 3D Gaussian Splatting in 200 Seconds}.
\newblock arXiv:2503.18402v2, 2025.
\newblock \url{https://arxiv.org/abs/2503.18402v2}

\bibitem{arxiv:2505.23158v2}
Jonas Kulhanek, Marie-Julie Rakotosaona, Fabian Manhardt, Christina Tsalicoglou, Michael Niemeyer, Torsten Sattler \emph{et al.}.
\newblock \emph{LODGE: Level-of-Detail Large-Scale Gaussian Splatting with Efficient Rendering}.
\newblock arXiv:2505.23158v2, 2025.
\newblock \url{https://arxiv.org/abs/2505.23158v2}

\bibitem{arxiv:2503.16924v2}
Joo Chan Lee, Jong Hwan Ko, Eunbyung Park.
\newblock \emph{Optimized Minimal 3D Gaussian Splatting}.
\newblock arXiv:2503.16924v2, 2025.
\newblock \url{https://arxiv.org/abs/2503.16924v2}

\bibitem{arxiv:2603.19234v2}
Zhilin Guo, Boqiao Zhang, Hakan Aktas, Kyle Fogarty, Jeffrey Hu, Nursena Koprucu Aslan \emph{et al.}.
\newblock \emph{Matryoshka Gaussian Splatting}.
\newblock arXiv:2603.19234v2, 2026.
\newblock \url{https://arxiv.org/abs/2603.19234v2}

\bibitem{arxiv:2603.25745v1}
Yixing Lao, Xuyang Bai, Xiaoyang Wu, Nuoyuan Yan, Zixin Luo, Tian Fang \emph{et al.}.
\newblock \emph{Less Gaussians, Texture More: 4K Feed-Forward Textured Splatting}.
\newblock arXiv:2603.25745v1, 2026.
\newblock \url{https://arxiv.org/abs/2603.25745v1}

\bibitem{arxiv:2410.08941v1}
Jaehoon Choi, Yonghan Lee, Hyungtae Lee, Heesung Kwon, Dinesh Manocha.
\newblock \emph{MeshGS: Adaptive Mesh-Aligned Gaussian Splatting for High-Quality Rendering}.
\newblock arXiv:2410.08941v1, 2024.
\newblock \url{https://arxiv.org/abs/2410.08941v1}

\bibitem{arxiv:2411.16392v4}
Ziyu Zhang, Binbin Huang, Hanqing Jiang, Liyang Zhou, Xiaojun Xiang, Shunhan Shen.
\newblock \emph{Quadratic Gaussian Splatting: High Quality Surface Reconstruction with Second-order Geometric Primitives}.
\newblock arXiv:2411.16392v4, 2024.
\newblock \url{https://arxiv.org/abs/2411.16392v4}

\bibitem{arxiv:2411.12788v2}
Guangchi Fang, Bing Wang.
\newblock \emph{Efficient Scene Modeling via Structure-Aware and Region-Prioritized 3D Gaussians}.
\newblock arXiv:2411.12788v2, 2024.
\newblock \url{https://arxiv.org/abs/2411.12788v2}

\bibitem{arxiv:2508.09977v5}
Shuting He, Peilin Ji, Yitong Yang, Changshuo Wang, Jiayi Ji, Yinglin Wang \emph{et al.}.
\newblock \emph{A Survey on 3D Gaussian Splatting Applications: Segmentation, Editing, and Generation}.
\newblock arXiv:2508.09977v5, 2025.
\newblock \url{https://arxiv.org/abs/2508.09977v5}

\bibitem{arxiv:2506.12400v2}
Hongbi Zhou, Zhangkai Ni.
\newblock \emph{Perceptual-GS: Scene-adaptive Perceptual Densification for Gaussian Splatting}.
\newblock arXiv:2506.12400v2, 2025.
\newblock \url{https://arxiv.org/abs/2506.12400v2}

\bibitem{arxiv:2411.14974v3}
Jan Held, Renaud Vandeghen, Abdullah Hamdi, Adrien Deliege, Anthony Cioppa, Silvio Giancola \emph{et al.}.
\newblock \emph{3D Convex Splatting: Radiance Field Rendering with 3D Smooth Convexes}.
\newblock arXiv:2411.14974v3, 2024.
\newblock \url{https://arxiv.org/abs/2411.14974v3}

\bibitem{arxiv:2604.03069v1}
Zicheng Zhang, Xiangting Meng, Ke Wu, Wenchao Ding.
\newblock \emph{SparseSplat: Towards Applicable Feed-Forward 3D Gaussian Splatting with Pixel-Unaligned Prediction}.
\newblock arXiv:2604.03069v1, 2026.
\newblock \url{https://arxiv.org/abs/2604.03069v1}

\bibitem{arxiv:2410.06245v1}
Shengji Tang, Weicai Ye, Peng Ye, Weihao Lin, Yang Zhou, Tao Chen \emph{et al.}.
\newblock \emph{HiSplat: Hierarchical 3D Gaussian Splatting for Generalizable Sparse-View Reconstruction}.
\newblock arXiv:2410.06245v1, 2024.
\newblock \url{https://arxiv.org/abs/2410.06245v1}

\bibitem{arxiv:2502.02283v6}
Zhihao Guo, Jingxuan Su, Chenghao Qian, Shenglin Wang, Jinlong Fan, Jing Zhang \emph{et al.}.
\newblock \emph{GP-GS: Gaussian Processes Densification for 3D Gaussian Splatting}.
\newblock arXiv:2502.02283v6, 2025.
\newblock \url{https://arxiv.org/abs/2502.02283v6}

\bibitem{arxiv:2508.09667v1}
Xingyilang Yin, Qi Zhang, Jiahao Chang, Ying Feng, Qingnan Fan, Xi Yang \emph{et al.}.
\newblock \emph{GSFixer: Improving 3D Gaussian Splatting with Reference-Guided Video Diffusion Priors}.
\newblock arXiv:2508.09667v1, 2025.
\newblock \url{https://arxiv.org/abs/2508.09667v1}

\bibitem{arxiv:2503.10860v2}
Avinash Paliwal, Xilong Zhou, Wei Ye, Jinhui Xiong, Rakesh Ranjan, Nima Khademi Kalantari.
\newblock \emph{RI3D: Few-Shot Gaussian Splatting With Repair and Inpainting Diffusion Priors}.
\newblock arXiv:2503.10860v2, 2025.
\newblock \url{https://arxiv.org/abs/2503.10860v2}

\bibitem{arxiv:2504.01503v2}
Ziteng Cui, Xuangeng Chu, Tatsuya Harada.
\newblock \emph{Luminance-GS: Adapting 3D Gaussian Splatting to Challenging Lighting Conditions with View-Adaptive Curve Adjustment}.
\newblock arXiv:2504.01503v2, 2025.
\newblock \url{https://arxiv.org/abs/2504.01503v2}

\bibitem{arxiv:2504.00219v1}
Han Zhou, Wei Dong, Jun Chen.
\newblock \emph{LITA-GS: Illumination-Agnostic Novel View Synthesis via Reference-Free 3D Gaussian Splatting and Physical Priors}.
\newblock arXiv:2504.00219v1, 2025.
\newblock \url{https://arxiv.org/abs/2504.00219v1}

\bibitem{arxiv:2501.05242v3}
Tianci Wen, Zhiang Liu, Yongchun Fang.
\newblock \emph{SEGS-SLAM: Structure-enhanced 3D Gaussian Splatting SLAM with Appearance Embedding}.
\newblock arXiv:2501.05242v3, 2025.
\newblock \url{https://arxiv.org/abs/2501.05242v3}

\bibitem{arxiv:2508.09479v2}
Xuejun Huang, Xinyi Liu, Yi Wan, Zhi Zheng, Bin Zhang, Mingtao Xiong \emph{et al.}.
\newblock \emph{SkySplat: Generalizable 3D Gaussian Splatting from Multi-Temporal Sparse Satellite Images}.
\newblock arXiv:2508.09479v2, 2025.
\newblock \url{https://arxiv.org/abs/2508.09479v2}

\bibitem{arxiv:2411.19454v2}
Jiepeng Wang, Yuan Liu, Peng Wang, Cheng Lin, Junhui Hou, Xin Li \emph{et al.}.
\newblock \emph{GausSurf: Geometry-Guided 3D Gaussian Splatting for Surface Reconstruction}.
\newblock arXiv:2411.19454v2, 2024.
\newblock \url{https://arxiv.org/abs/2411.19454v2}

\bibitem{arxiv:2601.17835v2}
Baowen Zhang, Chenxing Jiang, Heng Li, Shaojie Shen, Ping Tan.
\newblock \emph{Geometry-Grounded Gaussian Splatting}.
\newblock arXiv:2601.17835v2, 2026.
\newblock \url{https://arxiv.org/abs/2601.17835v2}

\bibitem{arxiv:2512.09162v3}
Kelian Baert, Mae Younes, Francois Bourel, Marc Christie, Adnane Boukhayma.
\newblock \emph{GTAvatar: Bridging Gaussian Splatting and Texture Mapping for Relightable and Editable Gaussian Avatars}.
\newblock arXiv:2512.09162v3, 2025.
\newblock \url{https://arxiv.org/abs/2512.09162v3}

\bibitem{arxiv:2412.12734v1}
Sebastian Weiss, Derek Bradley.
\newblock \emph{Gaussian Billboards: Expressive 2D Gaussian Splatting with Textures}.
\newblock arXiv:2412.12734v1, 2024.
\newblock \url{https://arxiv.org/abs/2412.12734v1}

\bibitem{arxiv:2412.01823v1}
Yunzhou Song, Heguang Lin, Jiahui Lei, Lingjie Liu, Kostas Daniilidis.
\newblock \emph{HDGS: Textured 2D Gaussian Splatting for Enhanced Scene Rendering}.
\newblock arXiv:2412.01823v1, 2024.
\newblock \url{https://arxiv.org/abs/2412.01823v1}

\bibitem{arxiv:2506.18575v3}
Kaifeng Sheng, Zheng Zhou, Yingliang Peng, Qianwei Wang.
\newblock \emph{2D Triangle Splatting for Direct Differentiable Mesh Training}.
\newblock arXiv:2506.18575v3, 2025.
\newblock \url{https://arxiv.org/abs/2506.18575v3}

\bibitem{arxiv:2511.19172v4}
Kehua Chen, Tianlu Mao, Xinzhu Ma, Hao Jiang, Zehao Li, Zihan Liu \emph{et al.}.
\newblock \emph{MetroGS: Efficient and Stable Reconstruction of Geometrically Accurate High-Fidelity Large-Scale Scenes}.
\newblock arXiv:2511.19172v4, 2025.
\newblock \url{https://arxiv.org/abs/2511.19172v4}

\bibitem{arxiv:2502.17531v2}
Hongyu Zhou, Zorah L\"{a}hner.
\newblock \emph{Laplace-Beltrami Operator for Gaussian Splatting}.
\newblock arXiv:2502.17531v2, 2025.
\newblock \url{https://arxiv.org/abs/2502.17531v2}

\bibitem{arxiv:2410.23213v1}
Muhammad Salman Ali, Sung-Ho Bae, Enzo Tartaglione.
\newblock \emph{ELMGS: Enhancing memory and computation scaLability through coMpression for 3D Gaussian Splatting}.
\newblock arXiv:2410.23213v1, 2024.
\newblock \url{https://arxiv.org/abs/2410.23213v1}

\bibitem{arxiv:2501.13045v1}
Yuang Shi, Simone Gasparini, G\'{e}raldine Morin, Chenggang Yang, Wei Tsang Ooi.
\newblock \emph{Sketch and Patch: Efficient 3D Gaussian Representation for Man-Made Scenes}.
\newblock arXiv:2501.13045v1, 2025.
\newblock \url{https://arxiv.org/abs/2501.13045v1}

\bibitem{arxiv:2501.13558v3}
Francesco Di Sario, Riccardo Renzulli, Marco Grangetto, Akihiro Sugimoto, Enzo Tartaglione.
\newblock \emph{GoDe: Gaussians on Demand for Progressive Level of Detail and Scalable Compression}.
\newblock arXiv:2501.13558v3, 2025.
\newblock \url{https://arxiv.org/abs/2501.13558v3}

\bibitem{arxiv:2410.13613v3}
Xinjie Zhang, Zhening Liu, Yifan Zhang, Xingtong Ge, Dailan He, Tongda Xu \emph{et al.}.
\newblock \emph{MEGA: Memory-Efficient 4D Gaussian Splatting for Dynamic Scenes}.
\newblock arXiv:2410.13613v3, 2024.
\newblock \url{https://arxiv.org/abs/2410.13613v3}

\bibitem{arxiv:2502.13196v2}
Pedro Martin, Ant\'{o}nio Rodrigues, Jo\~{a}o Ascenso, Maria Paula Queluz.
\newblock \emph{GS-QA: Comprehensive Quality Assessment Benchmark for Gaussian Splatting View Synthesis}.
\newblock arXiv:2502.13196v2, 2025.
\newblock \url{https://arxiv.org/abs/2502.13196v2}

\bibitem{arxiv:2510.08575v3}
Haofei Xu, Daniel Barath, Andreas Geiger, Marc Pollefeys.
\newblock \emph{ReSplat: Learning Recurrent Gaussian Splatting}.
\newblock arXiv:2510.08575v3, 2025.
\newblock \url{https://arxiv.org/abs/2510.08575v3}

\bibitem{arxiv:2512.15508v2}
Arthur Moreau, Richard Shaw, Michal Nazarczuk, Jisu Shin, Thomas Tanay, Zhensong Zhang \emph{et al.}.
\newblock \emph{Off The Grid: Detection of Primitives for Feed-Forward 3D Gaussian Splatting}.
\newblock arXiv:2512.15508v2, 2025.
\newblock \url{https://arxiv.org/abs/2512.15508v2}

\bibitem{arxiv:2312.12337v4}
David Charatan, Sizhe Li, Andrea Tagliasacchi, Vincent Sitzmann.
\newblock \emph{pixelSplat 3D Gaussian Splats from Image Pairs for Scalable Generalizable 3D Reconstruction}.
\newblock arXiv:2312.12337v4, 2023.
\newblock \url{https://arxiv.org/abs/2312.12337v4}

\bibitem{arxiv:2508.01171v2}
Ranran Huang, Krystian Mikolajczyk.
\newblock \emph{No Pose at All: Self-Supervised Pose-Free 3D Gaussian Splatting from Sparse Views}.
\newblock arXiv:2508.01171v2, 2025.
\newblock \url{https://arxiv.org/abs/2508.01171v2}

\bibitem{arxiv:2509.24421v5}
Yuanyuan Gao, Yuning Gong, Yifei Liu, Li Jingfeng, Dingwen Zhang, Yanci Zhang \emph{et al.}.
\newblock \emph{Proxy-GS: Unified Occlusion Priors for Training and Inference in Structured 3D Gaussian Splatting}.
\newblock arXiv:2509.24421v5, 2025.
\newblock \url{https://arxiv.org/abs/2509.24421v5}

\bibitem{arxiv:2406.17074v1}
Panagiotis Papantonakis, Georgios Kopanas, Bernhard Kerbl, Alexandre Lanvin, George Drettakis.
\newblock \emph{Reducing the Memory Footprint of 3D Gaussian Splatting}.
\newblock arXiv:2406.17074v1, 2024.
\newblock \url{https://arxiv.org/abs/2406.17074v1}

\bibitem{arxiv:2501.13975v2}
Lei Lan, Tianjia Shao, Zixuan Lu, Yu Zhang, Chenfanfu Jiang, Yin Yang.
\newblock \emph{3DGS\$\textasciicircum{}2\$: Near Second-order Converging 3D Gaussian Splatting}.
\newblock arXiv:2501.13975v2, 2025.
\newblock \url{https://arxiv.org/abs/2501.13975v2}

\bibitem{arxiv:2412.06299v1}
Jinbo Yan, Rui Peng, Luyang Tang, Ronggang Wang.
\newblock \emph{4D Gaussian Splatting with Scale-aware Residual Field and Adaptive Optimization for Real-time Rendering of Temporally Complex Dynamic Scenes}.
\newblock arXiv:2412.06299v1, 2024.
\newblock \url{https://arxiv.org/abs/2412.06299v1}

\bibitem{arxiv:2403.17898v1}
Kerui Ren, Lihan Jiang, Tao Lu, Mulin Yu, Linning Xu, Zhangkai Ni \emph{et al.}.
\newblock \emph{Octree-GS Towards Consistent Real-time Rendering with LOD-Structured 3D Gaussians}.
\newblock arXiv:2403.17898v1, 2024.
\newblock \url{https://arxiv.org/abs/2403.17898v1}

\bibitem{arxiv:2510.09997v2}
Zhigang Cheng, Mingchao Sun, Yu Liu, Zengye Ge, Luyang Tang, Mu Xu \emph{et al.}.
\newblock \emph{CLoD-GS: Continuous Level-of-Detail via 3D Gaussian Splatting}.
\newblock arXiv:2510.09997v2, 2025.
\newblock \url{https://arxiv.org/abs/2510.09997v2}

\bibitem{arxiv:2503.21226v2}
Yishai Lavi, Leo Segre, Shai Avidan.
\newblock \emph{Frequency-Aware Gaussian Splatting Decomposition}.
\newblock arXiv:2503.21226v2, 2025.
\newblock \url{https://arxiv.org/abs/2503.21226v2}

\bibitem{arxiv:2410.20815v3}
Jiawei Xu, Zexin Fan, Jian Yang, Jin Xie.
\newblock \emph{Grid4D: 4D Decomposed Hash Encoding for High-Fidelity Dynamic Gaussian Splatting}.
\newblock arXiv:2410.20815v3, 2024.
\newblock \url{https://arxiv.org/abs/2410.20815v3}

\bibitem{arxiv:2503.12307v1}
Jiahao Wu, Rui Peng, Zhiyan Wang, Lu Xiao, Luyang Tang, Jinbo Yan \emph{et al.}.
\newblock \emph{Swift4D:Adaptive divide-and-conquer Gaussian Splatting for compact and efficient reconstruction of dynamic scene}.
\newblock arXiv:2503.12307v1, 2025.
\newblock \url{https://arxiv.org/abs/2503.12307v1}

\bibitem{arxiv:2602.22571v1}
Tianyu Chen, Wei Xiang, Kang Han, Yu Lu, Di Wu, Gaowen Liu \emph{et al.}.
\newblock \emph{GIFSplat: Generative Prior-Guided Iterative Feed-Forward 3D Gaussian Splatting from Sparse Views}.
\newblock arXiv:2602.22571v1, 2026.
\newblock \url{https://arxiv.org/abs/2602.22571v1}

\bibitem{arxiv:2311.08581v1}
Wojciech Zielonka, Timur Bagautdinov, Shunsuke Saito, Michael Zollh\"{o}fer, Justus Thies, Javier Romero.
\newblock \emph{Drivable 3D Gaussian Avatars}.
\newblock arXiv:2311.08581v1, 2023.
\newblock \url{https://arxiv.org/abs/2311.08581v1}

\bibitem{arxiv:2403.05087v1}
Zhijing Shao, Zhaolong Wang, Zhuang Li, Duotun Wang, Xiangru Lin, Yu Zhang \emph{et al.}.
\newblock \emph{SplattingAvatar Realistic Real-Time Human Avatars with Mesh-Embedded Gaussian Splatting}.
\newblock arXiv:2403.05087v1, 2024.
\newblock \url{https://arxiv.org/abs/2403.05087v1}

\bibitem{arxiv:2503.06677v5}
Di Wu, Liu Liu, Zhou Linli, Anran Huang, Liangtu Song, Qiaojun Yu \emph{et al.}.
\newblock \emph{REArtGS: Reconstructing and Generating Articulated Objects via 3D Gaussian Splatting with Geometric and Motion Constraints}.
\newblock arXiv:2503.06677v5, 2025.
\newblock \url{https://arxiv.org/abs/2503.06677v5}

\bibitem{arxiv:2511.09827v2}
Aymen Mir, Jian Wang, Riza Alp Guler, Chuan Guo, Gerard Pons-Moll, Bing Zhou.
\newblock \emph{AHA! Animating Human Avatars in Diverse Scenes with Gaussian Splatting}.
\newblock arXiv:2511.09827v2, 2025.
\newblock \url{https://arxiv.org/abs/2511.09827v2}

\bibitem{arxiv:2407.11309v1}
Liuyue Xie, Joel Julin, Koichiro Niinuma, Laszlo A. Jeni.
\newblock \emph{Gaussian Splatting LK}.
\newblock arXiv:2407.11309v1, 2024.
\newblock \url{https://arxiv.org/abs/2407.11309v1}

\bibitem{arxiv:2411.18625v2}
Brian Chao, Hung-Yu Tseng, Lorenzo Porzi, Chen Gao, Tuotuo Li, Qinbo Li \emph{et al.}.
\newblock \emph{Textured Gaussians for Enhanced 3D Scene Appearance Modeling}.
\newblock arXiv:2411.18625v2, 2024.
\newblock \url{https://arxiv.org/abs/2411.18625v2}

\bibitem{arxiv:2507.17336v3}
Hyeongmin Lee, Kyungjune Baek.
\newblock \emph{Temporal Smoothness-Aware Rate-Distortion Optimized 4D Gaussian Splatting}.
\newblock arXiv:2507.17336v3, 2025.
\newblock \url{https://arxiv.org/abs/2507.17336v3}

\bibitem{arxiv:2412.13547v3}
Ankit Dhiman, Tao Lu, R Srinath, Emre Arslan, Angela Xing, Yuanbo Xiangli \emph{et al.}.
\newblock \emph{Turbo-GS: Accelerating 3D Gaussian Fitting for High-Quality Radiance Fields}.
\newblock arXiv:2412.13547v3, 2024.
\newblock \url{https://arxiv.org/abs/2412.13547v3}

\bibitem{arxiv:2411.06019v3}
Yangming Zhang, Wenqi Jia, Wei Niu, Miao Yin.
\newblock \emph{GaussianSpa: An "Optimizing-Sparsifying" Simplification Framework for Compact and High-Quality 3D Gaussian Splatting}.
\newblock arXiv:2411.06019v3, 2024.
\newblock \url{https://arxiv.org/abs/2411.06019v3}

\bibitem{arxiv:2505.10144v1}
Xuechang Tu, Lukas Radl, Michael Steiner, Markus Steinberger, Bernhard Kerbl, Fernando de la Torre.
\newblock \emph{VRSplat: Fast and Robust Gaussian Splatting for Virtual Reality}.
\newblock arXiv:2505.10144v1, 2025.
\newblock \url{https://arxiv.org/abs/2505.10144v1}

\bibitem{arxiv:2603.11531v1}
Xiaobiao Du, Yida Wang, Kun Zhan, Xin Yu.
\newblock \emph{Mobile-GS: Real-time Gaussian Splatting for Mobile Devices}.
\newblock arXiv:2603.11531v1, 2026.
\newblock \url{https://arxiv.org/abs/2603.11531v1}

\bibitem{arxiv:2402.07181v1}
Ben Fei, Jingyi Xu, Rui Zhang, Qingyuan Zhou, Weidong Yang, Ying He.
\newblock \emph{3D Gaussian as a New Vision Era A Survey}.
\newblock arXiv:2402.07181v1, 2024.
\newblock \url{https://arxiv.org/abs/2402.07181v1}

\bibitem{arxiv:2411.16816v3}
Georg Hess, Carl Lindstr\"{o}m, Maryam Fatemi, Christoffer Petersson, Lennart Svensson.
\newblock \emph{SplatAD: Real-Time Lidar and Camera Rendering with 3D Gaussian Splatting for Autonomous Driving}.
\newblock arXiv:2411.16816v3, 2024.
\newblock \url{https://arxiv.org/abs/2411.16816v3}

\bibitem{arxiv:2411.17190v5}
Gyeongjin Kang, Jisang Yoo, Jihyeon Park, Seungtae Nam, Hyeonsoo Im, Sangheon Shin \emph{et al.}.
\newblock \emph{SelfSplat: Pose-Free and 3D Prior-Free Generalizable 3D Gaussian Splatting}.
\newblock arXiv:2411.17190v5, 2024.
\newblock \url{https://arxiv.org/abs/2411.17190v5}

\bibitem{arxiv:2410.22128v2}
Sunghwan Hong, Jaewoo Jung, Heeseong Shin, Jisang Han, Jiaolong Yang, Chong Luo \emph{et al.}.
\newblock \emph{PF3plat: Pose-Free Feed-Forward 3D Gaussian Splatting}.
\newblock arXiv:2410.22128v2, 2024.
\newblock \url{https://arxiv.org/abs/2410.22128v2}

\bibitem{arxiv:2509.17246v2}
Ranran Huang, Krystian Mikolajczyk.
\newblock \emph{SPFSplatV2: Efficient Self-Supervised Pose-Free 3D Gaussian Splatting from Sparse Views}.
\newblock arXiv:2509.17246v2, 2025.
\newblock \url{https://arxiv.org/abs/2509.17246v2}

\bibitem{arxiv:2504.04294v1}
Zhisheng Huang, Peng Wang, Jingdong Zhang, Yuan Liu, Xin Li, Wenping Wang.
\newblock \emph{3R-GS: Best Practice in Optimizing Camera Poses Along with 3DGS}.
\newblock arXiv:2504.04294v1, 2025.
\newblock \url{https://arxiv.org/abs/2504.04294v1}

\bibitem{arxiv:2403.20309v1}
Zhiwen Fan, Wenyan Cong, Kairun Wen, Kevin Wang, Jian Zhang, Xinghao Ding \emph{et al.}.
\newblock \emph{InstantSplat Unbounded Sparse-view Pose-free Gaussian Splatting in 40 Seconds}.
\newblock arXiv:2403.20309v1, 2024.
\newblock \url{https://arxiv.org/abs/2403.20309v1}

\bibitem{arxiv:2503.06235v2}
Yang LI, Jinglu Wang, Lei Chu, Xiao Li, Shiu-hong Kao, Ying-Cong Chen \emph{et al.}.
\newblock \emph{StreamGS: Online Generalizable Gaussian Splatting Reconstruction for Unposed Image Streams}.
\newblock arXiv:2503.06235v2, 2025.
\newblock \url{https://arxiv.org/abs/2503.06235v2}

\bibitem{arxiv:2410.22817v2}
Zhiyuan Min, Yawei Luo, Jianwen Sun, Yi Yang.
\newblock \emph{Epipolar-Free 3D Gaussian Splatting for Generalizable Novel View Synthesis}.
\newblock arXiv:2410.22817v2, 2024.
\newblock \url{https://arxiv.org/abs/2410.22817v2}

\bibitem{arxiv:2505.15185v1}
Yifan Liu, Keyu Fan, Weihao Yu, Chenxin Li, Hao Lu, Yixuan Yuan.
\newblock \emph{MonoSplat: Generalizable 3D Gaussian Splatting from Monocular Depth Foundation Models}.
\newblock arXiv:2505.15185v1, 2025.
\newblock \url{https://arxiv.org/abs/2505.15185v1}

\bibitem{arxiv:2411.03637v1}
Rui Peng, Wangze Xu, Luyang Tang, Liwei Liao, Jianbo Jiao, Ronggang Wang.
\newblock \emph{Structure Consistent Gaussian Splatting with Matching Prior for Few-shot Novel View Synthesis}.
\newblock arXiv:2411.03637v1, 2024.
\newblock \url{https://arxiv.org/abs/2411.03637v1}

\bibitem{arxiv:2504.20378v1}
Jiang Wu, Rui Li, Yu Zhu, Rong Guo, Jinqiu Sun, Yanning Zhang.
\newblock \emph{Sparse2DGS: Geometry-Prioritized Gaussian Splatting for Surface Reconstruction from Sparse Views}.
\newblock arXiv:2504.20378v1, 2025.
\newblock \url{https://arxiv.org/abs/2504.20378v1}

\bibitem{arxiv:2403.14627v1}
Yuedong Chen, Haofei Xu, Chuanxia Zheng, Bohan Zhuang, Marc Pollefeys, Andreas Geiger \emph{et al.}.
\newblock \emph{MVSplat Efficient 3D Gaussian Splatting from Sparse Multi-View Images}.
\newblock arXiv:2403.14627v1, 2024.
\newblock \url{https://arxiv.org/abs/2403.14627v1}

\bibitem{arxiv:2602.06830v1}
Soonbin Lee, Yeong-Gyu Kim, Simon Sasse, Tomas M. Borges, Yago Sanchez, Eun-Seok Ryu \emph{et al.}.
\newblock \emph{GaussianPOP: Principled Simplification Framework for Compact 3D Gaussian Splatting via Error Quantification}.
\newblock arXiv:2602.06830v1, 2026.
\newblock \url{https://arxiv.org/abs/2602.06830v1}

\bibitem{arxiv:2606.09074v3}
Zhang Chen, Shuai Wan, Mengting Yu, Fuzheng Yang, Junhui Hou.
\newblock \emph{REFINE: Super-efficient 3D Gaussian Splatting Pruning via Rendering-Free Primitive Importance}.
\newblock arXiv:2606.09074v3, 2026.
\newblock \url{https://arxiv.org/abs/2606.09074v3}

\bibitem{arxiv:2601.20857v1}
Hongyu Zhou, Zisen Shao, Sheng Miao, Pan Wang, Dongfeng Bai, Bingbing Liu \emph{et al.}.
\newblock \emph{FreeFix: Boosting 3D Gaussian Splatting via Fine-Tuning-Free Diffusion Models}.
\newblock arXiv:2601.20857v1, 2026.
\newblock \url{https://arxiv.org/abs/2601.20857v1}

\bibitem{arxiv:2602.11638v3}
Hao Qin, Yukai Sun, Meng Wang, Ming Kong, Mengxu Lu, Qiang Zhu.
\newblock \emph{Variation-aware Flexible 3D Gaussian Editing}.
\newblock arXiv:2602.11638v3, 2026.
\newblock \url{https://arxiv.org/abs/2602.11638v3}

\bibitem{arxiv:2505.15208v4}
Wenjie Liu, Zhongliang Liu, Junwei Shu, Changbo Wang, Yang Li.
\newblock \emph{GT2-GS: Geometry-aware Texture Transfer for Gaussian Splatting}.
\newblock arXiv:2505.15208v4, 2025.
\newblock \url{https://arxiv.org/abs/2505.15208v4}

\bibitem{arxiv:2401.09720v2}
Mengtian Li, Shengxiang Yao, Zhifeng Xie, Keyu Chen.
\newblock \emph{GaussianBody Clothed Human Reconstruction via 3d Gaussian Splatting}.
\newblock arXiv:2401.09720v2, 2024.
\newblock \url{https://arxiv.org/abs/2401.09720v2}

\bibitem{arxiv:2510.12099v2}
Junfeng Ni, Yixin Chen, Zhifei Yang, Yu Liu, Ruijie Lu, Song-Chun Zhu \emph{et al.}.
\newblock \emph{G4Splat: Geometry-Guided Gaussian Splatting with Generative Prior}.
\newblock arXiv:2510.12099v2, 2025.
\newblock \url{https://arxiv.org/abs/2510.12099v2}

\bibitem{arxiv:2507.11061v3}
Hayeon Kim, Ji Ha Jang, Se Young Chun.
\newblock \emph{Robust 3D-Masked Part-level Editing in 3D Gaussian Splatting with Regularized Score Distillation Sampling}.
\newblock arXiv:2507.11061v3, 2025.
\newblock \url{https://arxiv.org/abs/2507.11061v3}

\bibitem{arxiv:2502.11801v2}
Sheng-Yu Huang, Zi-Ting Chou, Yu-Chiang Frank Wang.
\newblock \emph{3D Gaussian Inpainting with Depth-Guided Cross-View Consistency}.
\newblock arXiv:2502.11801v2, 2025.
\newblock \url{https://arxiv.org/abs/2502.11801v2}

\bibitem{arxiv:2402.06149v1}
Zhenglin Zhou, Fan Ma, Hehe Fan, Yi Yang.
\newblock \emph{HeadStudio Text to Animatable Head Avatars with 3D Gaussian Splatting}.
\newblock arXiv:2402.06149v1, 2024.
\newblock \url{https://arxiv.org/abs/2402.06149v1}

\bibitem{arxiv:2410.13862v3}
Haofei Xu, Songyou Peng, Fangjinhua Wang, Hermann Blum, Daniel Barath, Andreas Geiger \emph{et al.}.
\newblock \emph{DepthSplat: Connecting Gaussian Splatting and Depth}.
\newblock arXiv:2410.13862v3, 2024.
\newblock \url{https://arxiv.org/abs/2410.13862v3}

\bibitem{arxiv:2601.03824v3}
Wei Long, Haifeng Wu, Shiyin Jiang, Jinhua Zhang, Xinchun Ji, Shuhang Gu.
\newblock \emph{IDESplat: Iterative Depth Probability Estimation for Generalizable 3D Gaussian Splatting}.
\newblock arXiv:2601.03824v3, 2026.
\newblock \url{https://arxiv.org/abs/2601.03824v3}

\bibitem{arxiv:2512.17547v2}
Mehdi Hosseinzadeh, Shin-Fang Chng, Yi Xu, Simon Lucey, Ian Reid, Ravi Garg.
\newblock \emph{G3Splat: Geometrically Consistent Generalizable Gaussian Splatting}.
\newblock arXiv:2512.17547v2, 2025.
\newblock \url{https://arxiv.org/abs/2512.17547v2}

\bibitem{arxiv:2503.07819v2}
Joey Wilson, Marcelino Almeida, Sachit Mahajan, Martin Labrie, Maani Ghaffari, Omid Ghasemalizadeh \emph{et al.}.
\newblock \emph{POp-GS: Next Best View in 3D-Gaussian Splatting with P-Optimality}.
\newblock arXiv:2503.07819v2, 2025.
\newblock \url{https://arxiv.org/abs/2503.07819v2}

\bibitem{arxiv:2410.24204v3}
Kai Ye, Chong Gao, Guanbin Li, Wenzheng Chen, Baoquan Chen.
\newblock \emph{GeoSplatting: Towards Geometry Guided Gaussian Splatting for Physically-based Inverse Rendering}.
\newblock arXiv:2410.24204v3, 2024.
\newblock \url{https://arxiv.org/abs/2410.24204v3}

\bibitem{arxiv:2604.15941v1}
Haato Watanabe, Nobuyuki Umetani.
\newblock \emph{Neural Gabor Splatting: Enhanced Gaussian Splatting with Neural Gabor for High-frequency Surface Reconstruction}.
\newblock arXiv:2604.15941v1, 2026.
\newblock \url{https://arxiv.org/abs/2604.15941v1}

\bibitem{arxiv:2411.06976v2}
He Huang, Wenjie Huang, Qi Yang, Yiling Xu, Zhu li.
\newblock \emph{A Hierarchical Compression Technique for 3D Gaussian Splatting Compression}.
\newblock arXiv:2411.06976v2, 2024.
\newblock \url{https://arxiv.org/abs/2411.06976v2}

\bibitem{arxiv:2504.06815v1}
Hanxiao Sun, YuPeng Gao, Jin Xie, Jian Yang, Beibei Wang.
\newblock \emph{SVG-IR: Spatially-Varying Gaussian Splatting for Inverse Rendering}.
\newblock arXiv:2504.06815v1, 2025.
\newblock \url{https://arxiv.org/abs/2504.06815v1}

\bibitem{arxiv:2412.00905v2}
Youjia Zhang, Anpei Chen, Yumin Wan, Zikai Song, Junqing Yu, Yawei Luo \emph{et al.}.
\newblock \emph{Ref-GS: Directional Factorization for 2D Gaussian Splatting}.
\newblock arXiv:2412.00905v2, 2024.
\newblock \url{https://arxiv.org/abs/2412.00905v2}

\bibitem{arxiv:2502.09039v1}
Lingting Zhu, Guying Lin, Jinnan Chen, Xinjie Zhang, Zhenchao Jin, Zhao Wang \emph{et al.}.
\newblock \emph{Large Images are Gaussians: High-Quality Large Image Representation with Levels of 2D Gaussian Splatting}.
\newblock arXiv:2502.09039v1, 2025.
\newblock \url{https://arxiv.org/abs/2502.09039v1}

\bibitem{arxiv:2410.05111v3}
Qifeng Chen, Sheng Yang, Sicong Du, Tao Tang, Rengan Xie, Peng Chen \emph{et al.}.
\newblock \emph{LiDAR-GS:Real-time LiDAR Re-Simulation using Gaussian Splatting}.
\newblock arXiv:2410.05111v3, 2024.
\newblock \url{https://arxiv.org/abs/2410.05111v3}

\bibitem{arxiv:2506.07917v4}
Allen Tu, Haiyang Ying, Alex Hanson, Yonghan Lee, Tom Goldstein, Matthias Zwicker.
\newblock \emph{SpeeDe3DGS: Speedy Deformable 3D Gaussian Splatting with Temporal Pruning and Motion Grouping}.
\newblock arXiv:2506.07917v4, 2025.
\newblock \url{https://arxiv.org/abs/2506.07917v4}

\bibitem{arxiv:2504.00773v1}
Hyunwoo Park, Gun Ryu, Wonjun Kim.
\newblock \emph{DropGaussian: Structural Regularization for Sparse-view Gaussian Splatting}.
\newblock arXiv:2504.00773v1, 2025.
\newblock \url{https://arxiv.org/abs/2504.00773v1}

\bibitem{arxiv:2410.14189v1}
Wenyuan Zhang, Yu-Shen Liu, Zhizhong Han.
\newblock \emph{Neural Signed Distance Function Inference through Splatting 3D Gaussians Pulled on Zero-Level Set}.
\newblock arXiv:2410.14189v1, 2024.
\newblock \url{https://arxiv.org/abs/2410.14189v1}

\bibitem{arxiv:2411.15723v4}
Baixin Xu, Jiangbei Hu, Jiaze Li, Ying He.
\newblock \emph{GSurf: Learning Signed Distance Fields from Splatting Opaque Gaussians for High-quality 3D Reconstruction}.
\newblock arXiv:2411.15723v4, 2024.
\newblock \url{https://arxiv.org/abs/2411.15723v4}

\bibitem{arxiv:2311.17245v5}
Zhiwen Fan, Kevin Wang, Kairun Wen, Zehao Zhu, Dejia Xu, Zhangyang Wang.
\newblock \emph{LightGaussian Unbounded 3D Gaussian Compression with 15x Reduction and 200+ FPS}.
\newblock arXiv:2311.17245v5, 2023.
\newblock \url{https://arxiv.org/abs/2311.17245v5}

\bibitem{arxiv:2406.01597v2}
Henan Wang, Hanxin Zhu, Tianyu He, Runsen Feng, Jiajun Deng, Jiang Bian \emph{et al.}.
\newblock \emph{End-to-End Rate-Distortion Optimized 3D Gaussian Representation}.
\newblock arXiv:2406.01597v2, 2024.
\newblock \url{https://arxiv.org/abs/2406.01597v2}

\bibitem{arxiv:2410.01647v2}
Yang Cao, Yuanliang Ju, Dan Xu.
\newblock \emph{3DGS-DET: Empower 3D Gaussian Splatting with Boundary Guidance and Box-Focused Sampling for Indoor 3D Object Detection}.
\newblock arXiv:2410.01647v2, 2024.
\newblock \url{https://arxiv.org/abs/2410.01647v2}

\bibitem{arxiv:2409.18122v3}
Yuezhan Tao, Dexter Ong, Varun Murali, Igor Spasojevic, Pratik Chaudhari, Vijay Kumar.
\newblock \emph{RT-GuIDE: Real-Time Gaussian Splatting for Information-Driven Exploration}.
\newblock arXiv:2409.18122v3, 2024.
\newblock \url{https://arxiv.org/abs/2409.18122v3}

\bibitem{arxiv:2401.15318v1}
Yutao Feng, Xiang Feng, Yintong Shang, Ying Jiang, Chang Yu, Zeshun Zong \emph{et al.}.
\newblock \emph{Gaussian Splashing Dynamic Fluid Synthesis with Gaussian Splatting}.
\newblock arXiv:2401.15318v1, 2024.
\newblock \url{https://arxiv.org/abs/2401.15318v1}

\bibitem{arxiv:2604.21400v4}
Jinrang Jia, Zhenjia Li, Yifeng Shi.
\newblock \emph{You Only Gaussian Once: Controllable 3D Gaussian Splatting for Ultra-Densely Sampled Scenes}.
\newblock arXiv:2604.21400v4, 2026.
\newblock \url{https://arxiv.org/abs/2604.21400v4}

\bibitem{arxiv:2511.06830v3}
Tianang Chen, Jian Jin, Shilv Cai, Zhuangzi Li, Weisi Lin.
\newblock \emph{MUGSQA: Novel Multi-Uncertainty-Based Gaussian Splatting Quality Assessment Method, Dataset, and Benchmarks}.
\newblock arXiv:2511.06830v3, 2025.
\newblock \url{https://arxiv.org/abs/2511.06830v3}

\end{thebibliography}

\end{document}
